\chapter{Galois theory, finite and infinite}\label{ch:galois}

This chapter builds Galois theory from the definition of a field extension to the absolute Galois group $G_\Q=\Gal(\Qbar/\Q)$, its
continuous representations and the Chebotarev density theorem. The paper uses it in five places: the hypothesis of Theorem thm:main
(absolute irreducibility on $\Gal(\Qbar/\Q(\sqrt{\ell^*}))$), Proposition prop:ns (the determinant cuts out $\Q(\sqrt{-3})$, and the
$3$-division polynomial factors according to the orbits on the four lines), Lemma lem:old (Chebotarev and Brauer-Nesbitt), Lemmas
lem:tamagawa-exp and lem:tamfactor (Kummer theory and the Kummer cocycle), and Lemma lem:h0 (the cyclotomic character $\chi_\ell$ against
the quadratic character $\chi_{\ell^*}$). Group theory, profinite groups, rings, finite fields and representation theory are taken from
\Cref{ch:groups}, \Cref{ch:rings}, \Cref{ch:modules} and \Cref{ch:reps}; prime ideals and Frobenius elements in number fields are
\Cref{ch:ant}, local fields \Cref{ch:local}.

\section*{Core results at a glance}

Throughout, $F$ is a field, $\Fbar$ an algebraic closure, $F^{\mathrm{sep}}\subseteq\Fbar$ the separable closure and $G_F=\Gal(F^{\mathrm{sep}}/F)$;
$\Qbar\subset\C$ is the field of algebraic numbers, $G_\Q=\Gal(\Qbar/\Q)$ and $G_K=\Gal(\Qbar/K)$ for $K\subseteq\Qbar$; $\ell$ is a prime,
$\ell^*=(-1)^{(\ell-1)/2}\ell$ for odd $\ell$; $\mathcal{O}$ is the ring of integers of a finite extension of $\Q_\ell$, with uniformiser $\varpi$
and maximal ideal $\lambda=\varpi\mathcal{O}$.

\begin{enumerate}[label=\textbf{\thechapter.\arabic*}]
\item \textbf{Conventions.} Automorphisms compose as maps, $(\sigma\tau)(x)=\sigma(\tau(x))$; $\Frob_q$ is $x\mapsto x^q$ on $\Fbar_q$ and
$\Frob_p$ an arithmetic Frobenius, so $\chi_\ell(\Frob_p)=p$; $\chi_d$ ($d\in\Q^\times$ a non-square) is the character of $\Q(\sqrt d)$,
$\chi_{q^*}$ that of conductor $q$ (never $\chi_q$); $\chi_\ell$ is the $\ell$-adic cyclotomic character, $\bar\chi_\ell$ its reduction, $\chi_3^{\mathrm{cyc}}$
at the prime $3$; Kummer characters take values in $\mu_n$; $\Zl(1)=\varprojlim\mu_{\ell^k}$. \Cref{conv:galois:fields,conv:galois:frobenius,conv:galois:characters,conv:galois:reps}.
\item \textbf{Extensions and the tower law} (definition and proposition). $[L:F]=\dim_FL$; for $F\subseteq L\subseteq M$,
$[M:F]=[M:L][L:F]$, products of bases being bases. \emph{(proved here)} \Cref{def:galois:extension}, \Cref{prop:galois:tower}.
\item \textbf{Minimal polynomial.} If $\alpha$ is algebraic over $F$, the kernel of $F[x]\to L$, $x\mapsto\alpha$, is generated by a unique monic
irreducible $m_{\alpha,F}$; $F(\alpha)=F[\alpha]\simeq F[x]/(m_{\alpha,F})$, $[F(\alpha):F]=\deg m_{\alpha,F}$, and $g(\alpha)=0$ iff $m_{\alpha,F}\mid g$;
$\alpha$ is algebraic iff $[F(\alpha):F]<\infty$. \emph{(proved here)} \Cref{def:galois:algebraic}, \Cref{prop:galois:minimal-polynomial}.
\item \textbf{Algebraic extensions.} Finite extensions are algebraic; $F(\alpha_1,\dots,\alpha_r)$ with $\alpha_i$ algebraic is finite of degree at most
$\prod\deg m_{\alpha_i,F}$; algebraic elements form a subfield; algebraic over algebraic is algebraic. \emph{(proved here)} \Cref{prop:galois:algebraic-extensions}.
\item \textbf{Algebraic closure.} (a) Every field has an algebraic closure, unique up to (non-unique) isomorphism. (b) An embedding $F\to\Omega$,
$\Omega$ algebraically closed, extends to any algebraic $L/F$. (c) $\Qbar=\{z\in\C:z\text{ algebraic over }\Q\}$ is an algebraic closure of $\Q$,
countable. (d) Every $F$-embedding $\Fbar\to\Fbar$ is an automorphism, so every $F$-embedding $K\to\Fbar$ of a subextension $K\subseteq\Fbar$ extends to an element of
$\Aut(\Fbar/F)$. \emph{((a) and (b) for infinite $L$ proof sketched here; (b) for finite $L$, (c) and (d) given (b) proved here)} \Cref{thm:galois:algebraic-closure}.
\item \textbf{Splitting fields.} A splitting field of $f\in F[x]$ exists, of degree at most $n!$; an isomorphism $F\to F'$ extends to an isomorphism
of splitting fields of $f$ and $\phi(f)$, in at most $[L:F]$ ways, exactly $[L:F]$ if $f$ is separable. \emph{(proved here; families of polynomials
sketched)} \Cref{thm:galois:splitting-field}.
\item \textbf{Normal extensions.} For $L/F$ algebraic inside $\Fbar$, equivalent: every irreducible polynomial over $F$ with a root in $L$ splits in $L$;
$L$ is a splitting field of a family; every $F$-embedding $L\to\Fbar$ has image $L$. Normal closures exist; normality is not transitive.
\emph{(proved here)} \Cref{def:galois:normal}, \Cref{prop:galois:normal}.
\item \textbf{Separability.} An irreducible $f$ is inseparable iff $f'=0$ iff $\operatorname{char}F=p>0$ and $f\in F[x^p]$; $F$ is perfect iff
$\operatorname{char}F=0$ or $F=F^p$; $[L:F]_s=\#\mathrm{Emb}_F(L,\Fbar)$ divides $[L:F]$, with equality iff $L/F$ is separable, and is multiplicative in towers.
\emph{(proved here)} \Cref{def:galois:separable}, \Cref{prop:galois:separability}.
\item \textbf{Primitive element theorem.} A finite separable extension is simple; a finite extension is simple iff it has finitely many intermediate fields.
\emph{(proved here)} \Cref{thm:galois:primitive-element}.
\item \textbf{Separable closure.} $F^{\mathrm{sep}}$ is a field, normal and separable over $F$, equal to $\Fbar$ when $F$ is perfect.
\emph{(proved here)} \Cref{prop:galois:separable-closure}.
\item \textbf{Dedekind's lemma.} Distinct homomorphisms $G\to L^\times$ from a monoid are linearly independent over $L$. \emph{(proved here)} \Cref{thm:galois:dedekind}.
\item \textbf{Artin's theorem.} For a finite group $G$ of automorphisms of $L$, $L/L^G$ is finite, normal and separable of degree $\#G$, and
$\Aut(L/L^G)=G$. For any finite $L/F$, $\#\Aut(L/F)\le[L:F]$, with equality iff $L/F$ is normal and separable. \emph{(proved here)} \Cref{thm:galois:artin}.
\item \textbf{Galois extensions} (definition and theorem). $L/F$ algebraic is Galois if normal and separable; for $L/F$ finite: Galois iff splitting field
of a separable polynomial iff $L^{\Aut(L/F)}=F$ iff $\#\Aut(L/F)=[L:F]$. \emph{(proved here)} \Cref{def:galois:galois}, \Cref{thm:galois:galois-equivalences}.
\item \textbf{Fundamental theorem.} For $L/F$ finite Galois with group $G$, $H\mapsto L^H$ and $M\mapsto\Gal(L/M)$ are inverse inclusion-reversing
bijections; $[M:F]$ is the index of $\Gal(L/M)$; $\Gal(L/\sigma M)=\sigma\Gal(L/M)\sigma^{-1}$; $M/F$ is Galois iff $\Gal(L/M)$ is normal in $G$,
and then $G/\Gal(L/M)\simeq\Gal(M/F)$. \emph{(proved here)} \Cref{thm:galois:fundamental}.
\item \textbf{Translation.} For $L/F$ Galois and any extension $M/F$, the extension $LM/M$ is Galois and restriction is an isomorphism
$\Gal(LM/M)\simeq\Gal(L/L\cap M)$;
$\Gal(L_1L_2/F)\to\Gal(L_1/F)\times\Gal(L_2/F)$ is injective, bijective iff $L_1\cap L_2=F$. \emph{(proved here; the infinite case in
\Cref{cor:galois:translation-infinite})} \Cref{prop:galois:translation}.
\item \textbf{Finite fields.} Inside $\Fbar_p$, $\F_{q^n}$ is the unique extension of $\F_q$ of degree $n$; it is Galois with cyclic group generated
by $\Frob_q$; its subfields are the $\F_{q^d}$, $d\mid n$; an irreducible polynomial of degree $d$ over $\F_q$ has the roots $\alpha^{q^i}$,
$0\le i<d$. \emph{(proved here)} \Cref{thm:galois:finite-fields}.
\item \textbf{Cyclotomic polynomials modulo $p$.} For $p\nmid n$ and $f$ the order of $p$ in $(\Z/n)^\times$, $\Phi_n\bmod p$ is the product of
$\varphi(n)/f$ distinct irreducible polynomials of degree $f$, and $\F_p(\mu_n)=\F_{p^f}$. \emph{(proved here)} \Cref{prop:galois:cyclotomic-mod-p}.
\item \textbf{Quadratic extensions.} If $\operatorname{char}F\ne2$, the quadratic extensions are the $F(\sqrt d)$, in bijection with the nontrivial elements
of $F^\times/F^{\times2}$; two independent classes give a $V_4$-extension with three quadratic subfields. \emph{(proved here)} \Cref{prop:galois:quadratic}.
\item \textbf{Cyclotomic fields.} $\Phi_n$ is irreducible over $\Q$; $a\mapsto\sigma_a$, $\sigma_a(\zeta)=\zeta^a$ on $\mu_n$, is an isomorphism
$(\Z/n)^\times\simeq\Gal(\Q(\zeta_n)/\Q)$; $\Q(\zeta_m)\cap\Q(\zeta_n)=\Q(\zeta_{\gcd(m,n)})$, $\Q(\zeta_m)\Q(\zeta_n)=\Q(\zeta_{\lcm(m,n)})$;
$\Q(\zeta_3)=\Q(\sqrt{-3})$, $\Q(\zeta_8)=\Q(i,\sqrt2)$. \emph{(proved here)} \Cref{thm:galois:cyclotomic}.
\item \textbf{$\Q(\sqrt{\ell^*})$.} For odd $\ell$, $\Q(\sqrt{\ell^*})$ is the unique quadratic subfield of $\Q(\zeta_\ell)$, $\sigma_a(\sqrt{\ell^*})=
\bigl(\frac a\ell\bigr)\sqrt{\ell^*}$, so $\chi_{\ell^*}=\bar\chi_\ell^{(\ell-1)/2}$; $\chi_{\ell^*}=\bar\chi_\ell$ iff $\ell=3$; every quadratic field lies in
$\Q(\zeta_{|D|})$, $D$ its discriminant. \emph{(proved here)} \Cref{prop:galois:sqrt-lstar}.
\item \textbf{Norm and trace} (definition and proposition). $\Nm_{L/F}$, $\Tr_{L/F}$ are the determinant and trace of multiplication, the product and
sum of the conjugates; multiplicative, additive, transitive. \emph{(proved here)} \Cref{def:galois:norm-trace}, \Cref{prop:galois:norm-trace}.
\item \textbf{Hilbert's Theorem 90.} For $L/F$ cyclic with generator $\sigma$: $\Nm(b)=1$ iff $b=c/\sigma(c)$; $\Tr(b)=0$ iff $b=c-\sigma(c)$; for
$L/F$ finite Galois every $(a_\sigma)$ with $a_{\sigma\tau}=a_\sigma\sigma(a_\tau)$ is $a_\sigma=\sigma(c)/c$. \emph{(proved here)} \Cref{thm:galois:hilbert90}.
\item \textbf{Kummer characters.} If $\operatorname{char}F\nmid n$ and $\mu_n\subset F$, $\chi_a(\sigma)=\sigma(\alpha)/\alpha$ ($\alpha^n=a$) is a
homomorphism to $\mu_n$ independent of $\alpha$, with kernel $\Gal(L/F(\alpha))$; $[F(\alpha):F]$ is the order of $a$ in $F^\times/F^{\times n}$;
$\chi_a=1$ iff $a\in F^{\times n}$; $x^n-a$ is irreducible iff $a\notin F^{\times p}$ for every prime $p\mid n$. \emph{(proved here; Capelli's theorem
stated)} \Cref{prop:galois:kummer-character}.
\item \textbf{Cyclic extensions.} If $\operatorname{char}F\nmid n$ and $\mu_n\subset F$, every cyclic extension of degree $n$ is $F(\alpha)$ with
$\alpha^n\in F$. \emph{(proved here)} \Cref{thm:galois:kummer-cyclic}.
\item \textbf{Kummer theory.} If $\operatorname{char}F\nmid n$ and $\mu_n\subset F$, $\Delta\mapsto F(\Delta^{1/n})$ is a bijection from the groups
$F^{\times n}\subseteq\Delta\subseteq F^\times$ with $\Delta/F^{\times n}$ finite to the finite abelian extensions of exponent dividing $n$, and
$\Gal(F(\Delta^{1/n})/F)\times\Delta/F^{\times n}\to\mu_n$ is a perfect pairing. \emph{(proved here)} \Cref{thm:galois:kummer}.
\item \textbf{The affine shape.} Without $\mu_n\subset F$, $\sigma\mapsto\sm{u(\sigma)}{b(\sigma)}01$ embeds $\Gal(F(\mu_n,a^{1/n})/F)$ in
$\Z/n\rtimes(\Z/n)^\times$. \emph{(proved here)} \Cref{prop:galois:kummer-affine}.
\item \textbf{Kummer theory over $\F_q$.} For $\ell\nmid q$, $v=v_\ell(q-1)$, $u\in\F_q^\times$, $k\le v$: $\chi_u(\Frob_q)=u^{(q-1)/\ell^k}$; $u\mapsto
u^{(q-1)/\ell^v}$ identifies $\F_q^\times/(\F_q^\times)^{\ell^v}$ with $\mu_{\ell^v}(\F_q)$; the $\ell$-adic value $\kappa_u(\Frob_q)\in\Zl(1)/(q-1)$ is the class
of $u$, of valuation $\ge k$ iff $u\in(\F_q^\times)^{\ell^k}$. \emph{(proved here)} \Cref{prop:galois:kummer-finite-field}.
\item \textbf{Artin-Schreier theory.} In characteristic $p$, $x^p-x-a$ is split or irreducible with cyclic group of order $p$, and every cyclic
extension of degree $p$ is of this form; $\Hom_{\mathrm{cont}}(G_F,\Z/p)\simeq F/\wp(F)$. \emph{((a), (b) proved here; (c) proof sketched here)}
\Cref{thm:galois:artin-schreier}.
\item \textbf{Galois group of a polynomial.} $G_f$ acts faithfully on the roots; its orbits are the root sets of the irreducible factors, of sizes
$[F(r):F]$; over an intermediate field $M$ of a Galois $\Omega/F$, the irreducible factors over $M$ are the $\Gal(\Omega/M)$-orbits. \emph{(proved here)}
\Cref{prop:galois:polynomial-group}.
\item \textbf{Discriminant} (definition and proposition). For $f=a\prod_i(x-r_i)$,
\[
\disc(f)=a^{2n-2}\prod_{i<j}(r_i-r_j)^2=(-1)^{n(n-1)/2}a^{n-2}\prod_if'(r_i);
\]
if
$\operatorname{char}F\ne2$, $G_f\subseteq A_n$ iff $\disc f\in F^{\times2}$, and $F(\sqrt{\disc f})$ is cut out by the sign character; $\disc(x^3+px+q)=
-4p^3-27q^2$, $\disc(x^n-a)=(-1)^{(n-1)(n-2)/2}n^na^{n-1}$, $\disc(\Phi_\ell)=(-1)^{(\ell-1)/2}\ell^{\ell-2}$. \emph{(proved here)}
\Cref{def:galois:discriminant}, \Cref{prop:galois:discriminant}.
\item \textbf{Cubics.} An irreducible separable cubic has group $A_3$ or $S_3$ according as its discriminant is a square or not. \emph{(proved here)}
\Cref{prop:galois:cubic}.
\item \textbf{Quartics.} The resolvent cubic $R_f=y^3-by^2+(ac-4d)y-(a^2d-4bd+c^2)$ has $\disc R_f=\disc f$; an irreducible quartic has group $S_4$,
$A_4$, $V_4$, $D_4$ or $C_4$, read off from $R_f$, $\disc f$, and the factorisation of $f$ over $F(\sqrt{\disc f})$. \emph{(proved here)} \Cref{prop:galois:quartic}.
\item \textbf{Quartics and the four lines.} If the roots of a separable quartic $f$ correspond $G_F$-equivariantly to the four lines of $V=\F_3^2$ under
$\rho\colon G_F\to\GL(V)$ with image $H$, then $G_f$ is the image of $H$ in $\PGL_2(\F_3)\simeq S_4$, the factors of $f$ are the orbits of $H$ on the lines,
$F(\sqrt{\disc f})$ is cut out by $\det\rho$; the image is $\GL_2(\F_3)$, $N_{ns}(\F_3)$ or $N_s(\F_3)$ (when $H$ fixes no line and contains an element of
determinant $-1$ and trace $0$) according as $f$ is irreducible with group $S_4$, irreducible with group $D_4$, or a product of two irreducible quadratics.
\emph{(proved here)} \Cref{prop:galois:quartic-on-lines}.
\item \textbf{Nonempty inverse limits.} An inverse limit, over a directed set, of nonempty finite sets is nonempty. \emph{(proved here)}
\Cref{lem:galois:nonempty-limit}.
\item \textbf{Infinite Galois extensions.} $\Gal(L/F)=\varprojlim\Gal(K/F)$ over the finite Galois $K\subseteq L$, a profinite group in the Krull
topology. \emph{(proved here)} \Cref{prop:galois:infinite}.
\item \textbf{Krull's fundamental theorem.} Closed subgroups correspond to intermediate fields; $\Gal(L/L^H)=\overline H$; open subgroups correspond to
finite subextensions, normal closed subgroups to Galois subextensions, with $G/\Gal(L/M)\simeq\Gal(M/F)$ topologically. \emph{(proved here)} \Cref{thm:galois:krull}.
\item \textbf{The absolute Galois group} (definition and proposition). $G_K\subseteq G_F$ is open of index $[K:F]$ for $K/F$ finite separable, every open
subgroup is one, $G_K\trianglelefteq G_F$ iff $K/F$ is Galois; $G_{\F_q}\simeq\widehat\Z$ with $1\mapsto\Frob_q$; $G_\Q^{\mathrm{ab}}\simeq\widehat\Z^\times$.
\emph{(proved here, except the last statement, stated; reference)} \Cref{def:galois:absolute}, \Cref{prop:galois:absolute}.
\item \textbf{Absolute Kummer theory.} If $\operatorname{char}F\nmid n$ and $\mu_n\subset F$, $a\mapsto\chi_a$ is an isomorphism $F^\times/F^{\times n}\simeq
\Hom_{\mathrm{cont}}(G_F,\mu_n)$; for $n=2$, $\operatorname{char}F\ne2$: $F^\times/F^{\times2}\simeq\Hom_{\mathrm{cont}}(G_F,\{\pm1\})$. \emph{(proved here)}
\Cref{cor:galois:kummer-absolute}.
\item \textbf{Translation, infinite case.} The isomorphism $\Gal(LM/M)\simeq\Gal(L/L\cap M)$ holds for infinite $L/F$, as topological groups.
\emph{(proved here)} \Cref{cor:galois:translation-infinite}.
\item \textbf{The cyclotomic character} (definition and proposition). $\chi_\ell\colon G_\Q\to\Zl^\times$, $\sigma(\zeta)=\zeta^{\chi_\ell(\sigma)}$ on
$\mu_{\ell^\infty}$, is continuous and surjective with kernel $G_{\Q(\mu_{\ell^\infty})}$; $\bar\chi_\ell$ is surjective onto $\F_\ell^\times$ with kernel
$G_{\Q(\zeta_\ell)}$; over any $F$ of characteristic $\ne\ell$, $G_F$ acts on $\Zl(1)$ through $\chi_\ell$, and $\chi_\ell(\Frob_q)=q$ on $G_{\F_q}$.
\emph{(proved here)} \Cref{def:galois:cyclotomic-character}, \Cref{prop:galois:cyclotomic-character}.
\item \textbf{Complex conjugation.} $c\in G_\Q$ has order $2$, $\chi_\ell(c)=-1$, $\chi_d(c)=-1$ iff $d<0$; if $\det\rho(c)=-1$ for
$\rho\colon G_\Q\to\GL_2(k)$, $\operatorname{char}k\ne2$, then $\rho(c)\sim\operatorname{diag}(1,-1)$ has trace $0$ and $\rho$ is irreducible iff absolutely
irreducible. Every nontrivial element of finite order of $G_\Q$ is a complex conjugation. \emph{(proved here, except the last statement, stated;
reference)} \Cref{prop:galois:complex-conjugation}.
\item \textbf{Continuous representations} (definition and proposition). Over a discrete field, continuous means open kernel, and the field cut out is
finite Galois with group the image; over $\mathcal{O}$, continuous means continuous modulo every $\varpi^m$; every continuous representation over a
finite extension $E_\lambda$ of $\Q_\ell$ preserves a lattice; continuous complex representations of profinite groups have finite image.
\emph{(proved here)} \Cref{def:galois:continuous}, \Cref{prop:galois:continuous-reps}.
\item \textbf{Quadratic characters of $G_\Q$.} $d\mapsto\chi_d$ is an isomorphism $\Q^\times/\Q^{\times2}\simeq\Hom_{\mathrm{cont}}(G_\Q,\{\pm1\})$, $\ker\chi_d=
G_{\Q(\sqrt d)}$, $\chi_{dd'}=\chi_d\chi_{d'}$; for $p\nmid D$, $\chi_d(\Frob_p)=\bigl(\frac Dp\bigr)$ and $\chi_{\ell^*}(\Frob_p)=\bigl(\frac p\ell\bigr)$.
\emph{((a), (b) proved here; (c), (d) stated; reference)} \Cref{prop:galois:quadratic-characters}.
\item \textbf{Twists.} $\det(\rho\otimes\chi)=\chi^n\det\rho$; if $\rho\simeq\rho\otimes\chi$ in dimension $2$ then $\chi^2=1$; for
$\rho\colon G_\Q\to\GL_2(\Zl)$, $\ell$ odd, $\rho\bmod\ell^k\simeq(\rho\otimes\chi_\ell)\bmod\ell^k$ with $k\ge1$ forces $\ell=3$, $k=1$.
\emph{(proved here)} \Cref{prop:galois:twist}.
\item \textbf{Restriction to open subgroups.} $\rho(G_K)$ has index at most $[K:F]$ in $\rho(G_F)$; if $K$ is cut out by $\delta\circ\rho$ then
$\rho(G_K)=\rho(G_F)\cap\ker\delta$; for $E/\Q$, $\bar\rho_{E,\ell}(G_{\Q(\sqrt{\ell^*})})$ is the set of elements of the image with square determinant, the
determinant-one subgroup when $\ell=3$. \emph{(proved here)} \Cref{prop:galois:restriction}.
\item \textbf{Index two.} For $H\subseteq G$ open of index $2$ with character $\chi$ and $\rho$ two-dimensional, absolutely irreducible, $\operatorname{char}k\ne2$:
$\rho|_H$ is not absolutely irreducible iff $\rho\simeq\rho\otimes\chi$ iff $\rho\otimes\bar k\simeq\Ind_H^G\psi$ for a continuous character $\psi$ of $H$.
\emph{(proved here)} \Cref{prop:galois:index-two}.
\item \textbf{Reduction.} Reductions of two stable lattices have isomorphic semisimplifications; if one is irreducible the lattices are homothetic.
\emph{(proved here)} \Cref{prop:galois:reduction-ss}.
\item \textbf{The $\ell$-adic Kummer cocycle.} For $a\in F^\times$ and $\ell\ne\operatorname{char}F$, $\kappa_a(\sigma)=(\sigma(\alpha_k)/\alpha_k)_k\in\Zl(1)$
satisfies $\kappa_a(\sigma\tau)=\kappa_a(\sigma)+\chi_\ell(\sigma)\kappa_a(\tau)$, so, writing $\kappa_a=b\cdot e$ for a generator $e$ of $\Zl(1)$, $\sigma\mapsto\sm{\chi_\ell(\sigma)}{b(\sigma)}01$ is a continuous
representation; it changes by a coboundary with the root system; $\kappa_a(\sigma\tau\sigma^{-1})=\chi_\ell(\sigma)\kappa_a(\tau)$ for $\tau\in\ker\chi_\ell$;
if $\mu_{\ell^k}\subset K$, $\kappa_a\bmod\ell^k$ vanishes on $G_K$ iff $a\in K^{\times\ell^k}$. \emph{(proved here)} \Cref{prop:galois:l-adic-kummer}.
\item \textbf{Frobenius elements.} For $K/\Q$ finite Galois unramified at $p$, the Frobenius elements at the primes above $p$ form a conjugacy class
$\Frob_p$, trivial iff $p$ splits completely, compatible with quotients; $\Frob_p=\sigma_p$ in $\Q(\zeta_n)$; $\Frob_p$ acts on the roots of $f$ with the
cycle type of the factorisation of $f\bmod p$ (Dedekind). \emph{(stated; reference)} \Cref{thm:galois:frobenius}.
\item \textbf{Density} (definition). Dirichlet and natural density; density one; almost all. \Cref{def:galois:density}.
\item \textbf{Chebotarev's density theorem.} For $K/\Q$ finite Galois with group $G$ and a conjugacy class $C$, the primes with $\Frob_p=C$ have density
$\#C/\#G$. \emph{(stated; reference)} \Cref{thm:galois:chebotarev}.
\item \textbf{Density of Frobenius elements.} For $L/\Q$ Galois unramified outside a finite $S$ and $X$ of density one, the Frobenius elements at primes
over $X\setminus S$ are dense in $\Gal(L/\Q)$. \emph{(proved here, from Chebotarev)} \Cref{cor:galois:chebotarev-dense}.
\item \textbf{Chebotarev and Brauer-Nesbitt.} Two continuous semisimple $n$-dimensional representations of $G_\Q$, unramified outside a finite set, with the
same Frobenius characteristic polynomials on a set of density one, are isomorphic. \emph{(proved here)} \Cref{cor:galois:chebotarev-bn}.
\item \textbf{Congruent traces.} If $\rho_1,\rho_2\colon G_\Q\to\GL_2(\mathcal{O})$ have congruent traces and determinants of Frobenius modulo $\lambda$ on a
set of density one, then $\bar\rho_1^{ss}\simeq\bar\rho_2^{ss}$, and $\bar\rho_2\simeq\bar\rho_1$ if $\bar\rho_1$ is absolutely irreducible; (d) for a newform $h$
of weight $2$ on $\Gamma_0(M)$ and $\ell$ odd, $a_p(h)\equiv a_p(E)\pmod\lambda$ on a set of density one gives $\bar\rho_h\simeq\bar\rho_{E,\ell}\otimes k$ when $\bar\rho_{E,\ell}$ is irreducible.
\emph{((a) to (c) proved here; the newform statement (d) proved here, given the representations of newforms of \Cref{ch:modgalrep})}
\Cref{cor:galois:congruent-ap}.
\end{enumerate}

\section*{Where the paper uses this chapter}

\begin{itemize}
\item \textbf{The hypothesis of Theorem thm:main} (``$\bar\rho_{E,\ell}$ absolutely irreducible when restricted to $\Gal(\Qbar/\Q(\sqrt{\ell^*}))$''),
\textbf{the set $S_f$ of Section 4.2, Proposition prop:Sf and the quoted DDT Theorem 3.15}: $G_\Q$ and its open subgroups, \Cref{prop:galois:absolute};
the field $\Q(\sqrt{\ell^*})\subset\Q(\zeta_\ell)$, \Cref{prop:galois:sqrt-lstar}; restriction, \Cref{prop:galois:restriction}; the index-two equivalence,
\Cref{prop:galois:index-two}; irreducible equals absolutely irreducible, \Cref{prop:galois:complex-conjugation}.
\item \textbf{Proposition prop:ns} (the image $N_s(\F_3)$): $\det\bar\rho_{E,3}=\chi_{-3}$ cuts out $\Q(\sqrt{-3})$ and the image of
$\Gal(\Qbar/\Q(\sqrt{-3}))$ is the determinant-one subgroup, \Cref{prop:galois:restriction}; the induced case, \Cref{prop:galois:index-two}; the
factorisation of the $3$-division polynomial, \Cref{prop:galois:polynomial-group}, \Cref{prop:galois:discriminant}, \Cref{prop:galois:quartic-on-lines},
\Cref{ex:galois:ns-example}; complex conjugation, \Cref{prop:galois:complex-conjugation}.
\item \textbf{Lemma lem:h0} (the term $h_\ell(E)$) \textbf{and Section 1.3 (Conjecture conj:euler)}: $\chi_\ell(G_\Q)=\Zl^\times$,
\Cref{prop:galois:cyclotomic-character}; the determinant argument forcing $\ell=3$, $k=1$, \Cref{prop:galois:twist}; $\bar\chi_3=\chi_{-3}$ but
$\bar\chi_\ell\ne\chi_{\ell^*}$ for $\ell\ge5$, \Cref{prop:galois:sqrt-lstar}; induction from $\Q(\sqrt{-3})$, \Cref{prop:galois:index-two}.
\item \textbf{Lemma lem:old} (``by Chebotarev and Brauer-Nesbitt''): \Cref{thm:galois:frobenius}, \Cref{thm:galois:chebotarev},
\Cref{cor:galois:chebotarev-dense}, \Cref{cor:galois:chebotarev-bn}, \Cref{cor:galois:congruent-ap}, \Cref{prop:galois:reduction-ss}.
\item \textbf{Lemma lem:tamagawa-exp} (Kummer theory over $\Q_q^{\mathrm{ur}}$): \Cref{prop:galois:kummer-character}, \Cref{thm:galois:kummer},
\Cref{prop:galois:l-adic-kummer}(d).
\item \textbf{Lemma lem:tamfactor} (the Kummer cocycle of the Tate parameter): \Cref{prop:galois:l-adic-kummer}, \Cref{prop:galois:kummer-affine},
\Cref{prop:galois:kummer-finite-field}, \Cref{prop:galois:cyclotomic-character}.
\item \textbf{Section 1 (the table), Lemma lem:euler, Lemma lem:twist, Theorem thm:twist}: the characters $\chi_{-3}$, $\chi_{-4}$, $\chi_{q^*}$ and
$\chi_S$, \Cref{prop:galois:quadratic-characters}, \Cref{prop:galois:quadratic}; $a_p(f\otimes\chi)=\chi(p)a_p(f)$ as a twist, \Cref{prop:galois:twist}.
\item \textbf{Section 1.1 and Section 4} ($\rho\colon G_\Q\to\GL_2(\Zl)$, $\bar\rho$, $\Zl(1)$, $T(1)$): \Cref{prop:galois:continuous-reps},
\Cref{prop:galois:cyclotomic-character}.
\item \textbf{Section sec:data and the appendix} (the mod $3$ image labels checked against the $3$-division polynomial, the test \texttt{ns3}):
\Cref{prop:galois:quartic-on-lines}, \Cref{ex:galois:frobenius-on-lines}.
\end{itemize}

\section{Field extensions and algebraic elements}\label{sec:galois:extensions}

This section fixes the conventions of the chapter and introduces the basic objects: field extensions, their degree, algebraic elements and
their minimal polynomials, and the algebraic closures $\Fbar$ and $\Qbar$. Everything in the paper takes place inside $\Qbar$ or inside an
algebraic closure of $\Q_q$ or $\F_q$: the Galois group $G_\Q=\Gal(\Qbar/\Q)$ of Section 1.1 of the paper, the fields $\Q(\sqrt{\ell^*})$ and
$\Q(\sqrt{-3})$ of Theorem thm:main and Proposition prop:ns, and the unramified extension $\Q_q^{\mathrm{ur}}$ of Lemma lem:tamagawa-exp. Fields,
rings, ideals, polynomial rings, Gauss's lemma, Kronecker's construction $F[x]/(f)$ and finite fields are \Cref{ch:rings}.

\begin{convention}[Fields, closures, composition]\label{conv:galois:fields}
$\Qbar\subset\C$ is the field of complex numbers algebraic over $\Q$ (\Cref{thm:galois:algebraic-closure}); every algebraic extension of $\Q$ is taken inside
$\Qbar$, $G_\Q=\Gal(\Qbar/\Q)$ and $G_K=\Gal(\Qbar/K)\subseteq G_\Q$ for $K\subseteq\Qbar$. For a general field $F$ we fix an algebraic closure $\Fbar$ and
the separable closure $F^{\mathrm{sep}}\subseteq\Fbar$, and put $G_F=\Gal(F^{\mathrm{sep}}/F)$ (\Cref{def:galois:absolute}). Automorphisms are composed as
maps, $(\sigma\tau)(x)=\sigma(\tau(x))$, and groups act on the left. As in the rest of the book (Appendix B, convention B-C1), $p$, $q$, $\ell$ are primes,
except that in statements about finite fields $q$ is a power of a prime $p$; in the paper's applications $q$ is a bad prime of $E$ and $\F_q$ its
residue field, and the two meanings never occur in one display.
\end{convention}

\begin{convention}[Frobenius is arithmetic]\label{conv:galois:frobenius}
For a finite field $\F_q$, $\Frob_q\in\Gal(\Fbar_q/\F_q)$ is $x\mapsto x^q$. A Frobenius element $\Frob_p$ at a prime $p$ (\Cref{thm:galois:frobenius}, made
precise in \Cref{ch:ant}) is \emph{arithmetic}: it induces $x\mapsto x^p$ on residue fields, so it acts on $\mu_n$ ($p\nmid n$) by $\zeta\mapsto\zeta^p$
and $\chi_\ell(\Frob_p)=p$ for $\ell\ne p$. This is the convention of the proof of Lemma lem:tamfactor; the geometric Frobenius $\Frob_p^{-1}$ enters only
through $L$-factors, in later chapters.
\end{convention}

\begin{convention}[Characters]\label{conv:galois:characters}
\begin{enumerate}[label=(\alph*)]
\item For $d\in\Q^\times$ not a square, $\chi_d\colon G_\Q\to\{\pm1\}$ is $\sigma\mapsto\sigma(\sqrt d)/\sqrt d$. It depends only on the class of $d$
modulo squares, and its kernel is $G_{\Q(\sqrt d)}$. For an odd prime $q$, $q^*=(-1)^{(q-1)/2}q$ and $\chi_{q^*}$ is the quadratic character of conductor $q$,
the character of $\Q(\sqrt{q^*})$; the book never writes $\chi_q$ for it (the paper's $\chi_q$ is the book's $\chi_{q^*}$). $\chi_{-3}$ and $\chi_{-4}$ are the
characters of $\Q(\sqrt{-3})$ and $\Q(i)$. A numerical subscript always means such a quadratic character, written with the sign of $d$.
\item $\chi_\ell\colon G_\Q\to\Zl^\times$ is the $\ell$-adic cyclotomic character, $\sigma(\zeta)=\zeta^{\chi_\ell(\sigma)}$ on $\mu_{\ell^\infty}$
(\Cref{def:galois:cyclotomic-character}); $\bar\chi_\ell\colon G_\Q\to\F_\ell^\times$ is its reduction. A prime-letter subscript ($\chi_\ell$, $\chi_p$) always
means the cyclotomic character; at the number $3$ it is written $\chi_3^{\mathrm{cyc}}$, and its reduction $\bar\chi_3$. We shall prove
$\chi_{\ell^*}=\bar\chi_\ell^{(\ell-1)/2}$, which equals $\bar\chi_\ell$ only for $\ell=3$ (\Cref{prop:galois:sqrt-lstar}).
\item A Kummer character $\chi_a$ (\Cref{prop:galois:kummer-character}) takes values in the group $\mu_n$ of roots of unity, not in $\Z/n$; the two are
identified only after a primitive $n$-th root of unity is chosen. $\Zl(1)=\varprojlim_k\mu_{\ell^k}$, with transition maps $\zeta\mapsto\zeta^\ell$, is
written multiplicatively as a set of compatible systems $(\zeta_k)_k$ and additively when it receives cocycles: $(\zeta_k)+(\zeta'_k)=(\zeta_k\zeta'_k)$ and
$a\cdot(\zeta_k)=(\zeta_k^{a})$ for $a\in\Zl$ (with $\zeta_k^a$ meaning $\zeta_k^{a\bmod\ell^k}$).
\end{enumerate}
\end{convention}

\begin{convention}[Representations]\label{conv:galois:reps}
A representation of a profinite group $G$ is a continuous homomorphism $\rho\colon G\to\GL_n(k)$ (\Cref{def:galois:continuous}), where $k$ is a finite
field or $\Fbar_\ell$ with the discrete topology, or a finite extension $E_\lambda$ of $\Q_\ell$ with its $\ell$-adic topology, or the ring of integers
$\mathcal{O}$ of such an $E_\lambda$. For $\rho\colon G\to\GL_n(\mathcal{O})$, $\rho\bmod\varpi^m$ is the action on $\mathcal{O}^n/\varpi^m\mathcal{O}^n$ and
$\bar\rho=\rho\bmod\varpi$. Irreducibility, absolute irreducibility, $\rho^{ss}$, $\rho\otimes\chi$, $\Ind$ and $\Res$ are as in \Cref{ch:reps}.
``$\rho$ is absolutely irreducible on $G_K$'' means that $(\rho\otimes\bar k)|_{G_K}$ is irreducible.
\end{convention}

\begin{definition}[Field extensions]\label{def:galois:extension}
A \emph{field extension} $L/F$ is a field $L$ together with a subfield $F$. Then $L$ is an $F$-vector space, and its dimension $[L:F]=\dim_FL\in
\{1,2,3,\dots\}\cup\{\infty\}$ is the \emph{degree} of $L/F$; the extension is \emph{finite} if $[L:F]<\infty$. A field $M$ with $F\subseteq M\subseteq L$ is an
\emph{intermediate field}. For $S\subseteq L$, $F(S)$ is the smallest subfield of $L$ containing $F$ and $S$, and $F[S]$ the smallest subring; $F(\alpha_1,
\dots,\alpha_r)$ and $F[\alpha_1,\dots,\alpha_r]$ are the usual abbreviations. For two subfields $L_1,L_2$ of a field, the \emph{compositum} $L_1L_2$ is the
smallest subfield containing both. An \emph{$F$-embedding} of $L$ into a field $\Omega\supseteq F$ is a ring homomorphism $L\to\Omega$ that is the identity on
$F$ (it is automatically injective, $L$ being a field); $\mathrm{Emb}_F(L,\Omega)$ is the set of them, and $\Aut(L/F)$ the group of $F$-automorphisms of $L$.
\end{definition}

\begin{proposition}[Tower law]\label{prop:galois:tower}
Let $F\subseteq L\subseteq M$ be fields. If $(x_i)_{i\in I}$ is an $F$-basis of $L$ and $(y_j)_{j\in J}$ an $L$-basis of $M$, then $(x_iy_j)_{(i,j)\in I\times J}$
is an $F$-basis of $M$. Hence $[M:F]=[M:L]\,[L:F]$, with the convention $\infty\cdot n=n\cdot\infty=\infty$; in particular $M/F$ is finite iff $M/L$ and
$L/F$ are, and then $[L:F]$ and $[M:L]$ divide $[M:F]$.
\end{proposition}

\begin{proof}
Spanning: an element $z\in M$ is $z=\sum_jb_jy_j$ with $b_j\in L$ (finitely many nonzero), and each $b_j=\sum_ia_{ij}x_i$ with $a_{ij}\in F$; so
$z=\sum_{i,j}a_{ij}x_iy_j$. Independence: if $\sum_{i,j}a_{ij}x_iy_j=0$ with $a_{ij}\in F$, almost all zero, then $\sum_j\bigl(\sum_ia_{ij}x_i\bigr)y_j=0$ with
coefficients in $L$, so $\sum_ia_{ij}x_i=0$ for every $j$, and then $a_{ij}=0$ for all $i,j$. The degree formula is the count of $I\times J$.
\end{proof}

\begin{definition}[Algebraic elements]\label{def:galois:algebraic}
Let $L/F$ be an extension. An element $\alpha\in L$ is \emph{algebraic} over $F$ if $g(\alpha)=0$ for some nonzero $g\in F[x]$, and \emph{transcendental}
otherwise. The extension $L/F$ is \emph{algebraic} if every element of $L$ is algebraic over $F$.
\end{definition}

\begin{proposition}[Minimal polynomial]\label{prop:galois:minimal-polynomial}
Let $L/F$ be an extension and $\alpha\in L$ algebraic over $F$.
\begin{enumerate}[label=(\roman*)]
\item The kernel of the evaluation map $F[x]\to L$, $g\mapsto g(\alpha)$, is generated by a unique monic irreducible polynomial $m_{\alpha,F}$, the
\emph{minimal polynomial} of $\alpha$ over $F$. For $g\in F[x]$, $g(\alpha)=0$ iff $m_{\alpha,F}\mid g$; and $m_{\alpha,F}$ is the unique monic irreducible
polynomial in $F[x]$ with root $\alpha$.
\item $F[\alpha]=F(\alpha)\simeq F[x]/(m_{\alpha,F})$, and $1,\alpha,\dots,\alpha^{d-1}$ is an $F$-basis of $F(\alpha)$, where $d=\deg m_{\alpha,F}$; so
$[F(\alpha):F]=\deg m_{\alpha,F}$.
\item Conversely, if $\beta\in L$ and $[F(\beta):F]<\infty$, then $\beta$ is algebraic over $F$.
\item If $F\subseteq M\subseteq L$, then $m_{\alpha,M}$ divides $m_{\alpha,F}$ in $M[x]$.
\end{enumerate}
\end{proposition}

\begin{proof}
(i) The kernel $\mathfrak p$ is an ideal of the principal ideal domain $F[x]$ (\Cref{ch:rings}); it is nonzero as $\alpha$ is algebraic, and prime because
$F[x]/\mathfrak p$ embeds in the field $L$. A nonzero prime ideal of $F[x]$ is generated by an irreducible polynomial, unique up to a unit, hence unique
when monic. If $h$ is monic irreducible with $h(\alpha)=0$, then $m_{\alpha,F}\mid h$, and irreducibility forces $h=m_{\alpha,F}$.
(ii) The image of evaluation is $F[\alpha]$, so $F[\alpha]\simeq F[x]/(m_{\alpha,F})$, which is a field because $(m_{\alpha,F})$ is maximal (Kronecker's
construction, \Cref{ch:rings}, also for the basis). A subring of $L$ containing $F$ and $\alpha$ that is a field contains $F(\alpha)$; so $F[\alpha]=F(\alpha)$.
(iii) If $[F(\beta):F]=n<\infty$, the $n+1$ elements $1,\beta,\dots,\beta^n$ are linearly dependent over $F$, which is a nonzero polynomial relation.
(iv) $m_{\alpha,F}\in M[x]$ has the root $\alpha$; apply (i) over $M$.
\end{proof}

For example $m_{\sqrt2,\Q}=x^2-2$, $m_{2^{1/3},\Q}=x^3-2$ (Eisenstein, \Cref{ch:rings}), and $m_{\zeta_n,\Q}=\Phi_n$ (\Cref{thm:galois:cyclotomic}). The
number $\pi$ is transcendental over $\Q$ (Lindemann's theorem; we do not prove this, see \cite[Chapter 1]{MilneFields} for references); the book never uses
transcendental elements except through the countability argument of \Cref{thm:galois:algebraic-closure}(c).

\begin{proposition}[Algebraic extensions]\label{prop:galois:algebraic-extensions}
Let $L/F$ be an extension.
\begin{enumerate}[label=(\alph*)]
\item If $L/F$ is finite, it is algebraic.
\item If $\alpha_1,\dots,\alpha_r\in L$ are algebraic over $F$, then $F(\alpha_1,\dots,\alpha_r)=F[\alpha_1,\dots,\alpha_r]$ is finite over $F$, of degree at most
$\prod_i\deg m_{\alpha_i,F}$.
\item The elements of $L$ algebraic over $F$ form a subfield of $L$.
\item If $M/L$ and $L/F$ are algebraic, then $M/F$ is algebraic.
\item A finite extension is generated by finitely many algebraic elements.
\end{enumerate}
\end{proposition}

\begin{proof}
(a) For $\alpha\in L$, $F(\alpha)\subseteq L$ has finite degree; apply \Cref{prop:galois:minimal-polynomial}(iii).
(b) Induction on $r$. Put $F_{i}=F(\alpha_1,\dots,\alpha_i)$. The element $\alpha_{i+1}$ is a root of $m_{\alpha_{i+1},F}\in F_i[x]$, so it is algebraic over $F_i$
and $[F_{i+1}:F_i]=\deg m_{\alpha_{i+1},F_i}\le\deg m_{\alpha_{i+1},F}$ by \Cref{prop:galois:minimal-polynomial}(ii), (iv); also $F_{i+1}=F_i[\alpha_{i+1}]$. The tower
law gives the bound, and by induction $F_r=F[\alpha_1,\dots,\alpha_r]$.
(c) If $\alpha,\beta$ are algebraic, $F(\alpha,\beta)$ is finite by (b), so $\alpha\pm\beta$, $\alpha\beta$ and $\alpha/\beta$ ($\beta\ne0$), which lie in it,
are algebraic by (a).
(d) Let $\gamma\in M$, with $m_{\gamma,L}=x^m+c_{m-1}x^{m-1}+\dots+c_0$. The $c_i$ are algebraic over $F$, so $F'=F(c_0,\dots,c_{m-1})$ is finite over $F$ by (b), and
$\gamma$ is algebraic over $F'$, so $F'(\gamma)$ is finite over $F'$. By the tower law $F(\gamma)\subseteq F'(\gamma)$ is finite over $F$, and $\gamma$ is algebraic.
(e) An $F$-basis of $L$ generates $L$, and its elements are algebraic by (a).
\end{proof}

\begin{theorem}[Algebraic closure]\label{thm:galois:algebraic-closure}
A field $\Omega$ is \emph{algebraically closed} if every nonconstant polynomial in $\Omega[x]$ has a root in $\Omega$ (equivalently, splits into linear
factors). An \emph{algebraic closure} of $F$ is an algebraic extension of $F$ that is algebraically closed.
\begin{enumerate}[label=(\alph*)]
\item Every field $F$ has an algebraic closure $\Fbar$, and any two are $F$-isomorphic (by an isomorphism that is in general not unique).
\item Let $L/F$ be algebraic, $\Omega$ algebraically closed and $\sigma\colon F\to\Omega$ a field embedding. Then $\sigma$ extends to an embedding
$L\to\Omega$. If $L=F(\alpha)$, the extensions correspond bijectively to the distinct roots of $\sigma(m_{\alpha,F})$ in $\Omega$, through $\tau\mapsto\tau(\alpha)$.
\item $\Qbar=\{z\in\C:z\text{ is algebraic over }\Q\}$ is an algebraic closure of $\Q$. It is countable; hence $\C$ is not algebraic over $\Q$.
\item Every $F$-embedding $\Fbar\to\Fbar$ is an automorphism. Hence, for every subextension $K/F$ of $\Fbar/F$, every $F$-embedding $K\to\Fbar$ extends to an
element of $\Aut(\Fbar/F)$.
\end{enumerate}
\end{theorem}

\begin{proof}[Proof of (b) for finite $L$, and of (c)]
(b), $L=F(\alpha)$. If $\tau$ extends $\sigma$, then $0=\tau(m_{\alpha,F}(\alpha))=\sigma(m_{\alpha,F})(\tau(\alpha))$, so $\tau(\alpha)$ is a root of
$\sigma(m_{\alpha,F})$, and $\tau$ is determined by $\tau(\alpha)$ because $L=F[\alpha]$. Conversely, let $\beta\in\Omega$ be a root of
$\sigma(m_{\alpha,F})$ (one exists as $\Omega$ is algebraically closed). The homomorphism $F[x]\to\Omega$, $\sum a_ix^i\mapsto\sum\sigma(a_i)\beta^i$, kills
$m_{\alpha,F}$, so it factors through $F[x]/(m_{\alpha,F})\simeq F(\alpha)$ and gives $\tau$ with $\tau(\alpha)=\beta$. For a finite $L$, write
$L=F(\alpha_1,\dots,\alpha_r)$ (\Cref{prop:galois:algebraic-extensions}(e)) and extend step by step through $F(\alpha_1,\dots,\alpha_i)$.

(c) $\Qbar$ is a subfield of $\C$ by \Cref{prop:galois:algebraic-extensions}(c), algebraic over $\Q$ by definition. Let $g\in\Qbar[x]$ be nonconstant. By the
fundamental theorem of algebra $g$ has a root $z\in\C$. The coefficients of $g$ generate a finite extension $F'$ of $\Q$, and $z$ is algebraic over $F'$, so
$F'(z)/\Q$ is finite (tower law) and $z$ is algebraic over $\Q$: $z\in\Qbar$. Countability: $\Q[x]$ is countable and each nonzero polynomial has finitely many
roots in $\C$; as $\C$ is uncountable, some complex number is transcendental.
\end{proof}

\begin{proof}[Sketch of proof of (a), and of (b) for infinite $L$]
For (b) one applies Zorn's lemma to the set of pairs $(M,\tau)$ with $F\subseteq M\subseteq L$ and $\tau\colon M\to\Omega$
extending $\sigma$, ordered by extension: a maximal pair has $M=L$, since otherwise the case $M(\alpha)$ above extends it. For (a), one constructs an
algebraically closed field containing $F$ (Artin's construction: adjoin a root of every polynomial at once, in the quotient of a large polynomial ring by a
maximal ideal, and repeat countably many times, or once with an extra argument), and takes the algebraic elements in it; uniqueness follows from (b), an
$F$-embedding $\Fbar\to\Fbar'$ having algebraically closed image over which $\Fbar'$ is algebraic, hence being surjective. We do not give the details; see
\cite[Chapter V, Section 2]{Lang} or \cite[Section 13.4]{DummitFoote}.
\end{proof}

\begin{proof}[Proof of (d), given (b)]
Let $\sigma\colon\Fbar\to\Fbar$ be an $F$-embedding and $x\in\Fbar$. The set $R$ of roots of $m_{x,F}$ in $\Fbar$ is finite, and $\sigma$ maps $R$ injectively into $R$
(as $\sigma(m_{x,F})=m_{x,F}$); so $\sigma(R)=R$ and $x\in\sigma(\Fbar)$. Thus $\sigma$ is surjective. If $K\subseteq\Fbar$ and $\tau\colon K\to\Fbar$ is an $F$-embedding, then
$\Fbar/K$ is algebraic, so by (b) $\tau$ extends to an embedding $\Fbar\to\Fbar$, which is an automorphism by what we have shown.
\end{proof}

\begin{remark}\label{rem:galois:nonunique}
The non-uniqueness in (a) is the whole subject of this chapter: for perfect $F$ the $F$-automorphisms of $\Fbar$ are the elements of the absolute Galois
group $G_F$ (\Cref{def:galois:absolute}). For
$F=\Q$ we avoid all choices by working inside $\C$; the price is that one embedding $\Qbar\subset\C$ is singled out, and with it one element of $G_\Q$,
complex conjugation (\Cref{prop:galois:complex-conjugation}).
\end{remark}

\section{Splitting fields, normality, separability}\label{sec:galois:splitting}

A Galois extension is one that is normal (closed under conjugation) and separable (no repeated roots). This section develops the two notions and the
tools used everywhere afterwards: the extension of isomorphisms to splitting fields, the separable degree, and the primitive element theorem. In the
paper, every field ``cut out'' by a character or a representation (Proposition prop:ns: ``$\det\bar\rho_{E,3}=\chi_{-3}$ cuts out $\Q(\sqrt{-3})$'') is a
normal extension of $\Q$; the warning examples $\Q(2^{1/3})$ and $\Q(2^{1/4})$ below show that an extension generated by one root need not be.

\begin{theorem}[Splitting fields]\label{thm:galois:splitting-field}
Let $f\in F[x]$ have degree $n\ge1$. A \emph{splitting field} of $f$ over $F$ is an extension $L=F(r_1,\dots,r_n)$ such that $f=c\prod_{i=1}^n(x-r_i)$ in
$L[x]$, $c\in F^\times$.
\begin{enumerate}[label=(\alph*)]
\item A splitting field exists, and $[L:F]\le n!$.
\item Let $\phi\colon F\to F'$ be an isomorphism of fields, $L$ a splitting field of $f$ over $F$ and $L'$ a splitting field of $\phi(f)$ over $F'$. Then
$\phi$ extends to an isomorphism $L\to L'$. The number of extensions is at most $[L:F]$, and equal to $[L:F]$ if $f$ has $n$ distinct roots in $L$.
\item Any two splitting fields of $f$ over $F$ are $F$-isomorphic; inside a fixed $\Fbar$ the splitting field $F(\{r\in\Fbar:f(r)=0\})$ is unique.
\end{enumerate}
The same holds for a family $\mathcal{S}\subseteq F[x]$ of nonconstant polynomials: a splitting field is an extension generated by roots of the members of
$\mathcal{S}$ over which each of them splits; inside $\Fbar$ it is unique.
\end{theorem}

\begin{proof}
(a) Induction on $n$, the case $n=1$ being $L=F$. By Kronecker's construction (\Cref{ch:rings}) there is an extension $F(r_1)$ in which $f$ has a root
$r_1$, with $[F(r_1):F]\le n$ (the degree of the minimal polynomial of $r_1$, a factor of $f$). Then $f=(x-r_1)g$ with $g\in F(r_1)[x]$ of degree $n-1$, and a
splitting field $L$ of $g$ over $F(r_1)$, of degree at most $(n-1)!$, is one of $f$ over $F$. The tower law gives $[L:F]\le n\cdot(n-1)!$.

(b) Induction on $[L:F]$. If $[L:F]=1$, $f$ splits over $F$, so $\phi(f)$ splits over $F'$, $L'=F'$, and $\phi$ is the only extension. Otherwise some root
$r\in L$ of $f$ is not in $F$; let $g=m_{r,F}$, of degree $d>1$, a factor of $f$. Then $\phi(g)$ is an irreducible factor of $\phi(f)$, so it splits in
$L'$; let $r'_1,\dots,r'_e$ be its distinct roots in $L'$, with $e\le d$, and $e=d$ if $f$ (hence $g$) has distinct roots. Exactly as in the proof of
\Cref{thm:galois:algebraic-closure}(b), $\phi$ has exactly $e$ extensions $\phi_j\colon F(r)\to L'$, given by $r\mapsto r'_j$, and each is an isomorphism onto
$F'(r'_j)$. Now $L$ is a splitting field of $f$ over $F(r)$, $L'$ one of $\phi(f)=\phi_j(f)$ over $F'(r'_j)$, and $[L:F(r)]=[L:F]/d<[L:F]$. By induction each
$\phi_j$ extends to an isomorphism $L\to L'$, in at most $[L:F(r)]$ ways, exactly $[L:F(r)]$ if the roots are distinct. Every extension of $\phi$ to $L$
sends $r$ to some $r'_j$, so it is one of these. The total number is at most $e[L:F(r)]\le d[L:F(r)]=[L:F]$, with equality in the separable case. (That an
extension $\tau\colon L\to L'$ is surjective: $\tau(f)=\phi(f)=\phi(c)\prod(x-\tau(r_i))$, so the $\tau(r_i)$ are all the roots of $\phi(f)$ in $L'$, and they
generate $L'$.)

(c) The first statement is (b) with $\phi=\id$. Inside $\Fbar$, $f$ splits and any splitting field contained in $\Fbar$ is generated by all the roots of $f$ in
$\Fbar$.

For a family, the splitting field inside $\Fbar$ is the subfield generated by all roots in $\Fbar$ of all members of $\mathcal{S}$; existence and uniqueness up
to $F$-isomorphism of abstract ones follow from \Cref{thm:galois:algebraic-closure} (embed into $\Fbar$). For finite families, take the product polynomial.
\end{proof}

The bound in (a) improves to $[L:F]\mid n!$ (\Cref{exer:galois:nfactorial}).

\begin{definition}[Normal extensions]\label{def:galois:normal}
An algebraic extension $L/F$ with $L\subseteq\Fbar$ is \emph{normal} if it satisfies the equivalent conditions of \Cref{prop:galois:normal}. For $L/F$ finite, a
\emph{normal closure} of $L$ is a smallest normal extension of $F$ inside $\Fbar$ containing $L$.
\end{definition}

\begin{proposition}[Normality]\label{prop:galois:normal}
For an algebraic extension $F\subseteq L\subseteq\Fbar$ the following are equivalent:
\begin{enumerate}[label=(\roman*)]
\item every irreducible polynomial in $F[x]$ with a root in $L$ splits into linear factors in $L[x]$;
\item $L$ is the splitting field over $F$ of a family of polynomials;
\item every $F$-embedding $\tau\colon L\to\Fbar$ satisfies $\tau(L)=L$; that is, $\mathrm{Emb}_F(L,\Fbar)=\Aut(L/F)$.
\end{enumerate}
Moreover: if $L/F$ is normal and $F\subseteq M\subseteq L$, then $L/M$ is normal. A finite $L=F(\alpha_1,\dots,\alpha_r)$ has a unique normal closure inside
$\Fbar$, the splitting field of $\prod_i m_{\alpha_i,F}$, finite over $F$. Every extension of degree $2$ is normal, but normality is not transitive:
$\Q\subset\Q(\sqrt2)\subset\Q(2^{1/4})$, and $\Q(2^{1/3})/\Q$ is not normal.
\end{proposition}

\begin{proof}
(i)$\Rightarrow$(ii): $L$ is the splitting field of the family $\{m_{\alpha,F}:\alpha\in L\}$, each member splitting in $L$ by (i).
(ii)$\Rightarrow$(iii): let $L$ be the splitting field of $\mathcal{S}$, so $L$ is generated by the set $R$ of roots in $\Fbar$ of the members of $\mathcal{S}$,
and $R\subseteq L$. An $F$-embedding $\tau$ maps a root of $g\in\mathcal{S}$ to a root of $g$, so $\tau(R_g)\subseteq R_g$ for the finite set $R_g$ of roots of
$g$, with equality because $\tau$ is injective. Hence $\tau(L)$ is generated by $\tau(R)=R$, and $\tau(L)=L$.
(iii)$\Rightarrow$(i): let $g\in F[x]$ be irreducible with a root $\alpha\in L$, and $\beta\in\Fbar$ another root. By
\Cref{thm:galois:algebraic-closure}(b) there is an $F$-embedding $F(\alpha)\to\Fbar$ with $\alpha\mapsto\beta$, and it extends to an $F$-embedding
$\tau\colon L\to\Fbar$. By (iii) $\tau(L)=L$, so $\beta\in L$; all roots of $g$ lie in $L$.

If $L/F$ is normal and $F\subseteq M\subseteq L$, every $M$-embedding $L\to\Fbar$ is an $F$-embedding, so (iii) holds for $L/M$. For the normal closure: the
splitting field $N$ of $\prod_im_{\alpha_i,F}$ in $\Fbar$ is normal by (ii), finite by \Cref{thm:galois:splitting-field}(a), and contains $L$; a normal
extension of $F$ containing $L$ contains all roots of each $m_{\alpha_i,F}$ by (i), hence contains $N$.

Let $[L:F]=2$ and let $g\in F[x]$ be irreducible with a root $\alpha\in L$. If $\alpha\in F$, $g$ is linear. Otherwise $L=F(\alpha)$ and $g$ is (up to a
constant) $m_{\alpha,F}=x^2+bx+c$, whose second root $-b-\alpha$ lies in $L$; so (i) holds. Thus $\Q(\sqrt2)/\Q$ and $\Q(2^{1/4})/\Q(\sqrt2)$
are normal ($2^{1/4}$ has minimal polynomial $x^2-\sqrt2$ over $\Q(\sqrt 2)$, as $[\Q(2^{1/4}):\Q]=4$ by Eisenstein). But $x^4-2$ is irreducible over $\Q$ with the
root $2^{1/4}\in\Q(2^{1/4})\subset\R$ and the non-real root $i2^{1/4}$, so $\Q(2^{1/4})/\Q$ is not normal; likewise $x^3-2$ has one real root and two non-real ones.
\end{proof}

\begin{definition}[Separability]\label{def:galois:separable}
A polynomial $f\in F[x]$ is \emph{separable} if it has no repeated root in $\Fbar$. An algebraic element is \emph{separable} over $F$ if its minimal polynomial
is; an algebraic extension is separable if all its elements are. $F$ is \emph{perfect} if every irreducible polynomial over $F$ is separable. For $L/F$
finite, the \emph{separable degree} is $[L:F]_s=\#\mathrm{Emb}_F(L,\Fbar)$. The \emph{separable closure} of $F$ in $\Fbar$ is
$F^{\mathrm{sep}}=\{\alpha\in\Fbar:\alpha\text{ is separable over }F\}$.
\end{definition}

The formal derivative $f'$ satisfies the usual rules (\Cref{ch:rings}). Since the Euclidean algorithm computes $\gcd(f,f')$ inside $F[x]$, the gcd is the same
whether computed over $F$ or over $\Fbar$; and $r\in\Fbar$ is a repeated root of $f$ iff $f(r)=f'(r)=0$ (write $f=(x-r)h$, so $f'=h+(x-r)h'$). Hence
\emph{$f$ is separable iff $\gcd(f,f')=1$}.

\begin{proposition}[Separability]\label{prop:galois:separability}
\begin{enumerate}[label=(\alph*)]
\item An irreducible $f\in F[x]$ is inseparable iff $f'=0$, iff $\operatorname{char}F=p>0$ and $f\in F[x^p]$. If $\operatorname{char}F=p$, every irreducible $f$
can be written $f(x)=g(x^{p^e})$ with $e\ge0$ and $g$ irreducible and separable; then $f$ has exactly $\deg g$ distinct roots, each of multiplicity $p^e$.
\item $F$ is perfect iff $\operatorname{char}F=0$, or $\operatorname{char}F=p$ and $F=F^p$. Fields of characteristic $0$, finite fields and algebraically
closed fields are perfect; $\F_p(t)$ is not ($x^p-t$ is irreducible and inseparable).
\item Let $F\subseteq K\subseteq L$ be finite extensions inside $\Fbar$. Every $F$-embedding $K\to\Fbar$ has exactly $[L:K]_s$ extensions to $L$. Hence
$[L:F]_s=[L:K]_s[K:F]_s$.
\item For $L/F$ finite, $[L:F]_s$ divides $[L:F]$, and $[L:F]_s=[L:F]$ iff $L/F$ is separable. In particular $L=F(\alpha_1,\dots,\alpha_r)$ with all $\alpha_i$
separable is separable over $F$, and for $F\subseteq K\subseteq L$ finite, $L/F$ is separable iff $L/K$ and $K/F$ are.
\item Over a perfect field every algebraic extension is separable.
\end{enumerate}
\end{proposition}

\begin{proof}
(a) As $f$ is irreducible, $\gcd(f,f')$ is $1$ or $f$ (up to units); it is $f$ iff $f\mid f'$, iff $f'=0$ since $\deg f'<\deg f$. Writing
$f=\sum a_ix^i$, $f'=\sum ia_ix^{i-1}$ vanishes iff $ia_i=0$ for all $i$; in characteristic $0$ this forces $f$ constant, and in characteristic $p$ it
means $a_i=0$ unless $p\mid i$. For the second statement let $e$ be the largest integer with $f\in F[x^{p^e}]$ and $f(x)=g(x^{p^e})$. A factorisation of $g$
would give one of $f$, so $g$ is irreducible, and $g\notin F[x^p]$ by maximality of $e$, so $g$ is separable. If $g=\prod_j(x-\gamma_j)$ with distinct $\gamma_j$
(leading coefficient omitted) and $\beta_j\in\Fbar$ with $\beta_j^{p^e}=\gamma_j$, then $x^{p^e}-\gamma_j=(x-\beta_j)^{p^e}$ (the $p$-th power map is additive), so
$f=\prod_j(x-\beta_j)^{p^e}$ with distinct $\beta_j$.

(b) In characteristic $0$ apply (a). In characteristic $p$ with $F=F^p$: an inseparable irreducible $f=\sum a_ix^{pi}$ would be $\bigl(\sum b_ix^i\bigr)^p$ with
$b_i^p=a_i$, not irreducible. If $a\in F\setminus F^p$, let $\alpha^p=a$ in $\Fbar$; then $x^p-a=(x-\alpha)^p$, and a monic factor in $F[x]$ is $(x-\alpha)^k$
with $0<k<p$, whose coefficient of $x^{k-1}$ is $-k\alpha\notin F$ (as $\alpha\notin F$ and $k\in\F_p^\times$); so $x^p-a$ is irreducible and inseparable. For a
finite field the Frobenius $x\mapsto x^p$ is injective, hence surjective (\Cref{ch:rings}); over an algebraically closed field irreducible polynomials have
degree $1$. For $\F_p(t)$, $t\notin\F_p(t)^p$ (compare degrees in $t$).

(c) Fix an $F$-embedding $\sigma\colon K\to\Fbar$. If $L=K(\beta)$, the extensions of $\sigma$ correspond to the distinct roots of $\sigma(m_{\beta,K})$
(\Cref{thm:galois:algebraic-closure}(b)). Write $m_{\beta,K}=g(x^{p^e})$ as in (a) ($e=0$ in characteristic $0$). Then $\sigma(m_{\beta,K})=
\sigma(g)(x^{p^e})$, and $\sigma(g)$ is separable because $\gcd(\sigma(g),\sigma(g)')=\sigma(\gcd(g,g'))=1$; by the argument of (a), $\sigma(m_{\beta,K})$
has exactly $\deg g$ distinct roots, a number independent of $\sigma$. Taking $\sigma$ the inclusion, it is $[L:K]_s$. For general $L=K(\beta_1,\dots,\beta_s)$
extend in steps. The formula follows by counting pairs (embedding of $K$, extension to $L$).

(d) For $L=F(\alpha)$: $[L:F]_s$ is the number of distinct roots of $m_{\alpha,F}$, which is $\deg m_{\alpha,F}=[L:F]$ if $\alpha$ is separable and, by (a),
$\deg g=[L:F]/p^e$ with $e\ge1$ otherwise. So $[F(\alpha):F]_s$ divides $[F(\alpha):F]$, with equality iff $\alpha$ is separable. In general write
$L=F(\alpha_1,\dots,\alpha_r)$ and use (c) and the tower law along $F_i=F(\alpha_1,\dots,\alpha_i)$: $[L:F]_s=\prod[F_{i+1}:F_i]_s$ divides $\prod[F_{i+1}:F_i]=[L:F]$.
If $L/F$ is separable, each $\alpha_{i+1}$ is separable over $F$, hence over $F_i$ (its minimal polynomial over $F_i$ divides that over $F$), so equality holds
factor by factor. Conversely, if $[L:F]_s=[L:F]$ and $\alpha\in L$, then $[L:F(\alpha)]_s[F(\alpha):F]_s=[L:F(\alpha)][F(\alpha):F]$ with each factor on the left at
most the corresponding one on the right, so $[F(\alpha):F]_s=[F(\alpha):F]$ and $\alpha$ is separable. The last two statements follow from the equality
criterion and multiplicativity.

(e) Every irreducible polynomial over a perfect field is separable.
\end{proof}

\begin{theorem}[Primitive element]\label{thm:galois:primitive-element}
\begin{enumerate}[label=(\alph*)]
\item A finite separable extension $L/F$ is simple: $L=F(\theta)$ for some $\theta\in L$. If $F$ is infinite and $L=F(\alpha,\beta)$, one can take
$\theta=\alpha+c\beta$ for all but finitely many $c\in F$.
\item A finite extension $L/F$ is simple iff it has only finitely many intermediate fields.
\end{enumerate}
\end{theorem}

\begin{proof}
(a) If $F$ is finite, so is $L$, and $L^\times$ is cyclic (\Cref{ch:rings}); a generator $\theta$ of $L^\times$ satisfies $L=F(\theta)$. Let $F$ be infinite. By
induction on the number of generators it suffices to treat $L=F(\alpha,\beta)$, and we only need $\beta$ separable. Let $\alpha=\alpha_1,\dots,\alpha_n$ be
the roots of $m_{\alpha,F}$ and $\beta=\beta_1,\dots,\beta_m$ the distinct roots of $g=m_{\beta,F}$ in $\Fbar$. For each $i$ and each $j\ne1$ the equation
$\alpha_i+c\beta_j=\alpha+c\beta$ has at most one solution $c$; choose $c\in F$ avoiding these finitely many values and put $\theta=\alpha+c\beta$. The
polynomials $g(x)$ and $h(x)=m_{\alpha,F}(\theta-cx)$ lie in $F(\theta)[x]$ and both vanish at $\beta$. If $h(\beta_j)=0$ with $j\ne1$, then $\theta-c\beta_j=\alpha_i$
for some $i$, which $c$ was chosen to avoid. As $g$ has distinct roots, $\gcd(g,h)=x-\beta$ in $\Fbar[x]$, hence in $F(\theta)[x]$; so $\beta\in F(\theta)$, and
then $\alpha=\theta-c\beta\in F(\theta)$.

(b) Suppose $L=F(\theta)$ and let $M$ be intermediate, $g=m_{\theta,M}$, and $M_0\subseteq M$ the field generated over $F$ by the coefficients of $g$. Then $g$ is
irreducible over $M_0$ and $g(\theta)=0$, so $[L:M_0]=\deg g=[L:M]$ and $M=M_0$. Thus $M$ is determined by $g$, a monic divisor of $m_{\theta,F}$ in $L[x]$, of
which there are finitely many. Conversely, suppose there are finitely many intermediate fields. If $F$ is finite, (a) applies. If $F$ is infinite and
$L=F(\alpha,\beta)$, the fields $F(\alpha+c\beta)$, $c\in F$, cannot all differ, so $F(\alpha+c\beta)=F(\alpha+c'\beta)=:M$ for some $c\ne c'$; then $M$ contains
$(c-c')\beta$, hence $\beta$ and $\alpha$, and $L=M$ is simple. Induction on the number of generators finishes the proof.
\end{proof}

\begin{proposition}[Separable closure]\label{prop:galois:separable-closure}
$F^{\mathrm{sep}}$ is a subfield of $\Fbar$, normal and separable over $F$, and every separable algebraic extension of $F$ inside $\Fbar$ lies in it. If $F$
is perfect then $F^{\mathrm{sep}}=\Fbar$; in particular $\Q^{\mathrm{sep}}=\Qbar$, $\F_q^{\mathrm{sep}}=\Fbar_q$ and $\Q_p^{\mathrm{sep}}=\overline{\Q}_p$.
\end{proposition}

\begin{proof}
If $\alpha,\beta\in\Fbar$ are separable over $F$, then $F(\alpha,\beta)/F$ is separable by \Cref{prop:galois:separability}(d), so $\alpha\pm\beta$, $\alpha\beta$,
$\alpha/\beta$ are separable: $F^{\mathrm{sep}}$ is a field, separable over $F$ by definition. If $\alpha\in F^{\mathrm{sep}}$, all roots of $m_{\alpha,F}$ have the
same separable minimal polynomial, so they lie in $F^{\mathrm{sep}}$: condition (i) of \Cref{prop:galois:normal} holds. The rest is the definition and
\Cref{prop:galois:separability}(e).
\end{proof}

\section{Galois extensions and the fundamental theorem}\label{sec:galois:fundamental}

This section proves the fundamental theorem of Galois theory for finite extensions: subgroups of the Galois group correspond to intermediate fields.
Every ``fixed field'', ``cut out by'' and ``subgroup of index two'' in the paper is an instance: the statement of Proposition prop:ns that the image
of $\Gal(\Qbar/\Q(\sqrt{-3}))$ is the determinant-one subgroup combines the fundamental theorem with \Cref{prop:galois:restriction}. The translation
theorem at the end is what makes restriction of representations to $G_K$ well behaved, and in \Cref{ch:local} it says that an unramified base change
does not change the inertia group (the first sentence of the proof of Lemma lem:tamagawa-exp).

\begin{theorem}[Dedekind's lemma]\label{thm:galois:dedekind}
Let $G$ be a monoid, $L$ a field and $\chi_1,\dots,\chi_n\colon G\to L^\times$ distinct monoid homomorphisms. If $a_1,\dots,a_n\in L$ satisfy
$\sum_ia_i\chi_i(g)=0$ for all $g\in G$, then $a_1=\dots=a_n=0$. In particular distinct field embeddings $\sigma_1,\dots,\sigma_n\colon K\to L$ are linearly
independent over $L$ (apply the lemma to $G=K^\times$).
\end{theorem}

\begin{proof}
Suppose not, and take a nontrivial relation $\sum_ia_i\chi_i=0$ with the least number of nonzero coefficients; renumber so that $a_1,\dots,a_m\ne0$ and
the others vanish. Then $m\ge2$, since $\chi_1(g)\ne0$. Choose $h$ with $\chi_1(h)\ne\chi_2(h)$. For all $g$, $\sum_{i\le m}a_i\chi_i(hg)=\sum_{i\le m}
a_i\chi_i(h)\chi_i(g)=0$, and $\chi_1(h)\sum_{i\le m}a_i\chi_i(g)=0$. Subtracting, $\sum_{2\le i\le m}a_i(\chi_i(h)-\chi_1(h))\chi_i(g)=0$, a relation with fewer
nonzero coefficients whose coefficient at $i=2$ is nonzero: a contradiction.
\end{proof}

\begin{theorem}[Artin]\label{thm:galois:artin}
\begin{enumerate}[label=(\alph*)]
\item Let $L$ be a field, $G\subseteq\Aut(L)$ a finite group and $F=L^G=\{x\in L:\sigma(x)=x\text{ for all }\sigma\in G\}$. Then $L/F$ is finite, normal and
separable, $[L:F]=\#G$ and $\Aut(L/F)=G$. For $\alpha\in L$, $m_{\alpha,F}=\prod_{\beta\in G\alpha}(x-\beta)$, the product over the orbit.
\item For every finite extension $L/F$, $\#\Aut(L/F)\le[L:F]_s\le[L:F]$, and $\#\Aut(L/F)=[L:F]$ iff $L/F$ is normal and separable.
\end{enumerate}
\end{theorem}

\begin{proof}
(a) Let $\alpha\in L$ with orbit $G\alpha=\{\beta_1,\dots,\beta_k\}$, $k\le\#G$. The polynomial $P=\prod_j(x-\beta_j)$ has distinct roots, and every
$\sigma\in G$ permutes them, so its coefficients are fixed by $G$: $P\in F[x]$. If $g\in F[x]$ vanishes at $\alpha$, it vanishes at each $\sigma(\alpha)$, so $P\mid g$;
hence $P=m_{\alpha,F}$. Thus every $\alpha\in L$ is algebraic and separable over $F$ of degree at most $\#G$, and its minimal polynomial splits in $L$: $L/F$ is
algebraic, separable and normal. Every finite subextension $K/F$ of $L$ is separable, hence $K=F(\theta)$ (\Cref{thm:galois:primitive-element}), and
$[K:F]\le\#G$. Choose such a $K$ of maximal degree; for $\beta\in L$, $K(\beta)$ is again finite over $F$, so $K(\beta)=K$. Hence $L=K$ and $[L:F]\le\#G$. On the
other hand $G\subseteq\Aut(L/F)$, and by (b) $\#\Aut(L/F)\le[L:F]$. So $\#G\le\#\Aut(L/F)\le[L:F]\le\#G$, all are equal and $G=\Aut(L/F)$.

(b) Embed $L$ in $\Fbar$. Composing with the inclusion $L\subseteq\Fbar$ makes $\Aut(L/F)$ a subset of $\mathrm{Emb}_F(L,\Fbar)$, of size $[L:F]_s\le[L:F]$
(\Cref{prop:galois:separability}(d)). If $L/F$ is normal and separable, every $F$-embedding has image $L$ (\Cref{prop:galois:normal}), so $\Aut(L/F)=
\mathrm{Emb}_F(L,\Fbar)$ has $[L:F]_s=[L:F]$ elements. Conversely, if $\#\Aut(L/F)=[L:F]$, then $[L:F]_s=[L:F]$, so $L/F$ is separable, and
$\Aut(L/F)=\mathrm{Emb}_F(L,\Fbar)$, so every embedding has image $L$ and $L/F$ is normal.
\end{proof}

\begin{definition}[Galois extensions]\label{def:galois:galois}
An algebraic extension $L/F$ is \emph{Galois} if it is normal and separable. Its \emph{Galois group} is $\Gal(L/F)=\Aut(L/F)$.
\end{definition}

\begin{theorem}[Characterisations of finite Galois extensions]\label{thm:galois:galois-equivalences}
For a finite extension $L/F$ the following are equivalent:
\begin{enumerate}[label=(\roman*)]
\item $L/F$ is Galois;
\item $L$ is the splitting field over $F$ of a separable polynomial;
\item $L^{\Aut(L/F)}=F$;
\item $\#\Aut(L/F)=[L:F]$.
\end{enumerate}
If $L/F$ is Galois and $F\subseteq M\subseteq L$, then $L/M$ is Galois.
\end{theorem}

\begin{proof}
(i)$\Leftrightarrow$(iv) is \Cref{thm:galois:artin}(b). (i)$\Rightarrow$(ii): $L=F(\theta)$ (\Cref{thm:galois:primitive-element}); $m_{\theta,F}$ is separable,
splits in $L$ by normality, and its roots generate $L$. (ii)$\Rightarrow$(iv): by \Cref{thm:galois:splitting-field}(b) with $\phi=\id$, $L=L'$, the identity of
$F$ has exactly $[L:F]$ extensions to automorphisms of $L$. (i)$\Rightarrow$(iii): put $G=\Aut(L/F)$ and $F'=L^G\supseteq F$. By \Cref{thm:galois:artin}(a),
$[L:F']=\#G=[L:F]$, so $F'=F$. (iii)$\Rightarrow$(i): $G=\Aut(L/F)$ is finite (\Cref{thm:galois:artin}(b)), and \Cref{thm:galois:artin}(a) shows that $L/L^G=L/F$ is
normal and separable. Finally, if $L$ is the splitting field over $F$ of a separable $f$, it is the splitting field of $f$ over $M$.
\end{proof}

\begin{theorem}[Fundamental theorem of Galois theory]\label{thm:galois:fundamental}
Let $L/F$ be a finite Galois extension with group $G=\Gal(L/F)$.
\begin{enumerate}[label=(\alph*)]
\item The maps $H\mapsto L^H$ and $M\mapsto\Gal(L/M)$ are mutually inverse, inclusion-reversing bijections between the subgroups $H\le G$ and the
intermediate fields $F\subseteq M\subseteq L$.
\item $[L:M]=\#\Gal(L/M)$ and $[M:F]=[G:\Gal(L/M)]$.
\item For $\sigma\in G$, $\Gal(L/\sigma M)=\sigma\Gal(L/M)\sigma^{-1}$.
\item $M/F$ is Galois iff $\Gal(L/M)$ is normal in $G$. Then restriction $\sigma\mapsto\sigma|_M$ is a surjection $G\to\Gal(M/F)$ with kernel $\Gal(L/M)$,
so $G/\Gal(L/M)\simeq\Gal(M/F)$.
\item The compositum $M_1M_2$ corresponds to $\Gal(L/M_1)\cap\Gal(L/M_2)$, and $M_1\cap M_2$ to the subgroup generated by $\Gal(L/M_1)$ and $\Gal(L/M_2)$.
\end{enumerate}
\end{theorem}

\begin{proof}
(a), (b). For $H\le G$, \Cref{thm:galois:artin}(a) gives $[L:L^H]=\#H$ and $\Gal(L/L^H)=H$. For $M$ intermediate, $L/M$ is Galois
(\Cref{thm:galois:galois-equivalences}), so $L^{\Gal(L/M)}=M$ by (iii) there. Both maps reverse inclusions. The tower law gives $[M:F]=[L:F]/[L:M]=
\#G/\#\Gal(L/M)$.
(c) $\tau$ fixes $\sigma M$ iff $\sigma^{-1}\tau\sigma$ fixes $M$.
(d) The extension $M/F$ is separable. Every $F$-embedding $M\to\Fbar$ (with $L\subseteq\Fbar$) extends to an $F$-embedding $L\to\Fbar$
(\Cref{thm:galois:algebraic-closure}(b)), which has image $L$ because $L/F$ is normal, so it is an element of $G$. Hence the $F$-embeddings of $M$ are the
$\sigma|_M$, $\sigma\in G$, and by \Cref{prop:galois:normal}(iii) $M/F$ is normal iff $\sigma M=M$ for all $\sigma\in G$, iff (by (a) and (c))
$\sigma\Gal(L/M)\sigma^{-1}=\Gal(L/M)$ for all $\sigma$. In that case $\sigma\mapsto\sigma|_M$ is a homomorphism $G\to\Gal(M/F)$, surjective because every
automorphism of $M/F$ is an embedding and so extends, as just shown, with kernel $\Gal(L/M)$.
(e) $\sigma$ fixes $M_1M_2$ iff it fixes $M_1$ and $M_2$. The fixed field of the subgroup generated by $H_1$ and $H_2$ is $L^{H_1}\cap L^{H_2}$; apply (a).
\end{proof}

\begin{proposition}[Translation]\label{prop:galois:translation}
Let $L/F$ be finite Galois and $M/F$ any extension, with $L$ and $M$ contained in a common field. Then $LM/M$ is finite Galois, and restriction
$\Gal(LM/M)\to\Gal(L/F)$, $\sigma\mapsto\sigma|_L$, is injective with image $\Gal(L/L\cap M)$. Consequently:
\begin{enumerate}[label=(\alph*)]
\item $[LM:M]=[L:L\cap M]$; in particular $[LM:M]$ divides $[L:F]$.
\item If $L_1,L_2$ are finite Galois over $F$, so are $L_1L_2$ and $L_1\cap L_2$, and $\Gal(L_1L_2/F)\to\Gal(L_1/F)\times\Gal(L_2/F)$, $\sigma\mapsto(\sigma|_{L_1},
\sigma|_{L_2})$, is injective, with image the pairs that agree on $L_1\cap L_2$; it is bijective iff $L_1\cap L_2=F$.
\end{enumerate}
\end{proposition}

\begin{proof}
Let $L$ be the splitting field over $F$ of a separable $f$ with roots $r_1,\dots,r_n$. Then $LM=M(r_1,\dots,r_n)$ is the splitting field of $f$ over $M$, so
$LM/M$ is Galois. An element $\sigma\in\Gal(LM/M)$ permutes the $r_i$, so $\sigma(L)=L$ and $\sigma|_L\in\Gal(L/F)$; if $\sigma|_L=\id$, then $\sigma$ fixes $M$
and $L$, hence $LM$. Let $H$ be the image. An element $x\in L$ fixed by $H$ is fixed by $\Gal(LM/M)$, so $x\in M$ (\Cref{thm:galois:galois-equivalences}(iii));
hence $L^H=L\cap M$ (the inclusion $\supseteq$ being clear) and $H=\Gal(L/L\cap M)$ by \Cref{thm:galois:fundamental}(a).
(a) Count. (b) $L_1L_2$ is the splitting field of the product of the distinct irreducible factors of $f_1f_2$, which is separable; $L_1\cap L_2$ is normal by
\Cref{prop:galois:normal}(i) and separable. The map is injective as $L_1L_2$ is generated by $L_1$ and $L_2$. Its image lies in the set of pairs agreeing on
$L_1\cap L_2$, which has $\#\Gal(L_1/F)\cdot\#\Gal(L_2/L_1\cap L_2)$ elements (the second coordinate is any lift of the restriction of the first); and by (a),
$\#\Gal(L_1L_2/F)=[L_1L_2:L_1][L_1:F]=[L_2:L_1\cap L_2][L_1:F]$ is the same number.
\end{proof}

\section{Finite fields, quadratic and cyclotomic extensions}\label{sec:galois:examples}

The three families of Galois extensions that the paper uses constantly are worked out here. Finite fields give the residue fields $\F_q$ of
Lemmas lem:tamfactor and lem:euler and the honest model of a Frobenius element; quadratic fields give the characters $\chi_d$ of the table of Section 1,
of Lemma lem:twist and of Theorem thm:twist; cyclotomic fields give $\Q(\zeta_3)=\Q(\sqrt{-3})$ (Proposition prop:ns, Lemma lem:h0) and the field
$\Q(\sqrt{\ell^*})\subset\Q(\zeta_\ell)$ of the hypothesis of Theorem thm:main. The section ends with the splitting field of $x^3-2$, the model of a
non-abelian Galois group and of a quadratic subfield cut out by a sign.

\begin{theorem}[Galois theory of finite fields]\label{thm:galois:finite-fields}
Let $q=p^r$ and $n\ge1$, and work inside a fixed algebraic closure $\Fbar_p$; $\F_q$ is the unique subfield with $q$ elements.
\begin{enumerate}[label=(\alph*)]
\item $\F_{q^n}=\{x\in\Fbar_p:x^{q^n}=x\}$ is the splitting field of $x^{q^n}-x$ over $\F_q$, it is the unique extension of $\F_q$ of degree $n$ inside $\Fbar_p$,
and every extension of $\F_q$ of degree $n$ is $\F_q$-isomorphic to it.
\item $\F_{q^n}/\F_q$ is Galois, and $\Gal(\F_{q^n}/\F_q)$ is cyclic of order $n$, generated by $\Frob_q\colon x\mapsto x^q$.
\item The intermediate fields of $\F_{q^n}/\F_q$ are the $\F_{q^d}$ with $d\mid n$, and $\F_{q^d}$ is the fixed field of $\langle\Frob_q^d\rangle$.
\item If $g\in\F_q[x]$ is irreducible of degree $d$ with a root $\alpha$, then $\F_q(\alpha)=\F_{q^d}$ and $g$ has the $d$ distinct roots
$\alpha,\alpha^q,\dots,\alpha^{q^{d-1}}$ in $\F_{q^d}$. Every finite extension of a finite field is Galois with cyclic group.
\end{enumerate}
\end{theorem}

\begin{proof}
(a) The polynomial $x^{q^n}-x$ has derivative $-1$, so it is separable, and its $q^n$ roots in $\Fbar_p$ form a field (\Cref{ch:rings}: the roots of
$x^{p^m}-x$ are closed under the field operations). It contains $\F_q$, whose elements satisfy $x^q=x$. A subextension $K\subseteq\Fbar_p$ of degree $n$ over
$\F_q$ has $q^n$ elements, all satisfying $x^{q^n}=x$ (Lagrange in $K^\times$), so $K=\F_{q^n}$ by counting. An abstract extension of degree $n$ embeds over $\F_q$
into $\Fbar_p$ (\Cref{thm:galois:algebraic-closure}(b)), with image $\F_{q^n}$.
(b) $\Frob_q$ is an injective ring endomorphism of the finite field $\F_{q^n}$, hence an automorphism; it fixes $\F_q$. Its powers $\Frob_q^k$, $1\le k<n$,
are not the identity, since $x^{q^k}=x$ has at most $q^k<q^n$ solutions. So $\langle\Frob_q\rangle\subseteq\Aut(\F_{q^n}/\F_q)$ has order $n=[\F_{q^n}:\F_q]$,
and \Cref{thm:galois:artin}(b) gives equality and that the extension is Galois.
(c) The subgroups of the cyclic group of order $n$ are the $\langle\Frob_q^d\rangle$, $d\mid n$, of order $n/d$; the fixed field of $\Frob_q^d$ is
$\{x:x^{q^d}=x\}=\F_{q^d}$. Apply \Cref{thm:galois:fundamental}.
(d) $\F_q(\alpha)$ has degree $d$, so it is $\F_{q^d}$ by (a), Galois over $\F_q$. The roots of $g$ form the orbit of $\alpha$ under $\Gal(\F_{q^d}/\F_q)=
\langle\Frob_q\rangle$ (\Cref{thm:galois:artin}(a)), namely the $\Frob_q^i(\alpha)=\alpha^{q^i}$, $0\le i<d$, which are distinct since the orbit has $\deg g=d$
elements.
\end{proof}

\begin{proposition}[Cyclotomic polynomials modulo $p$]\label{prop:galois:cyclotomic-mod-p}
Let $n\ge1$, $p$ a prime with $p\nmid n$, $f$ the order of $p$ in $(\Z/n)^\times$, and $\bar\Phi_n\in\F_p[x]$ the reduction of $\Phi_n$. Then
$\bar\Phi_n=\prod_\zeta(x-\zeta)$, the product over the $\varphi(n)$ elements of order exactly $n$ in $\Fbar_p^\times$; it is the product of $\varphi(n)/f$ distinct
monic irreducible polynomials of degree $f$; and $\F_p(\mu_n)=\F_{p^f}$.
\end{proposition}

\begin{proof}
As $p\nmid n$, $x^m-1$ is separable over $\F_p$ for every $m\mid n$ (its derivative $mx^{m-1}$ is prime to it), so $\mu_n(\Fbar_p)$ is a group of order $n$, cyclic
(\Cref{ch:rings}), with exactly $\varphi(n)$ elements of order $n$. Reducing $x^m-1=\prod_{e\mid m}\Phi_e$ (\Cref{ch:rings}) modulo $p$, the roots of
$\prod_{e\mid m}\bar\Phi_e$ are the elements of order dividing $m$, each once. The roots of $\bar\Phi_n$ are therefore the roots of $x^n-1$ that are not roots of
$\bar\Phi_e$ for any proper divisor $e$ of $n$, that is, which have order dividing no proper divisor of $n$: the elements of order exactly $n$. Let $\zeta$ be one.
By \Cref{thm:galois:finite-fields}(d) its minimal polynomial over $\F_p$ has the roots $\zeta^{p^i}$, $0\le i<d$, where $d$ is the least integer with
$\zeta^{p^d}=\zeta$, that is, with $p^d\equiv1\pmod n$; so $d=f$ for every $\zeta$. Grouping the $\varphi(n)$ roots into these orbits gives $\varphi(n)/f$ distinct
irreducible factors of degree $f$, and $\F_p(\mu_n)=\F_p(\zeta)=\F_{p^f}$.
\end{proof}

This is the finite-field shadow of the decomposition of $p$ in $\Q(\zeta_n)$: in \Cref{ch:ant} it gives that $p\nmid n$ splits into $\varphi(n)/f$ primes of
residue degree $f$, and that $\Frob_p$ is $\sigma_p$ (\Cref{thm:galois:frobenius}).

\begin{proposition}[Quadratic extensions]\label{prop:galois:quadratic}
Let $\operatorname{char}F\ne2$.
\begin{enumerate}[label=(\alph*)]
\item Every extension $L/F$ of degree $2$ is Galois and equals $F(\sqrt d)$ for some $d\in F^\times\setminus F^{\times2}$; $\Gal(L/F)=\{1,\sigma\}$ with
$\sigma(\sqrt d)=-\sqrt d$.
\item Inside $\Fbar$, $F(\sqrt d)=F(\sqrt{d'})$ iff $d/d'\in F^{\times2}$. So $d\mapsto F(\sqrt d)$ is a bijection from the nontrivial elements of
$F^\times/F^{\times2}$ to the quadratic extensions of $F$ in $\Fbar$.
\item The quadratic fields are the $\Q(\sqrt d)$, $d\ne0,1$ a squarefree integer, pairwise distinct.
\item If the classes of $a,b\in F^\times$ in $F^\times/F^{\times2}$ are independent (that is, none of $a$, $b$, $ab$ is a square), then $F(\sqrt a,\sqrt b)/F$ is
Galois with group $\{\pm1\}^2\simeq V_4$, the element $(\epsilon_1,\epsilon_2)$ acting by $\sqrt a\mapsto\epsilon_1\sqrt a$, $\sqrt b\mapsto\epsilon_2\sqrt b$; its
quadratic subfields are exactly $F(\sqrt a)$, $F(\sqrt b)$ and $F(\sqrt{ab})$.
\end{enumerate}
\end{proposition}

\begin{proof}
(a) Let $\alpha\in L\setminus F$, so $L=F(\alpha)$ and $m_{\alpha,F}=x^2+bx+c$. Then $(\alpha+b/2)^2=b^2/4-c=:d$, and $d$ is not a square in $F$ (otherwise
$\alpha\in F$), so $L=F(\sqrt d)$ with $\sqrt d=\alpha+b/2$. The roots $\pm\sqrt d$ are distinct as $d\ne0$ and $2\ne0$, so $L$ is the splitting field of the
separable $x^2-d$, Galois of degree $2$, and the nontrivial automorphism sends $\sqrt d$ to the other root.
(b) If $d'=dt^2$, then $\sqrt{d'}=\pm t\sqrt d$. Conversely if $\sqrt{d'}=x+y\sqrt d$ with $x,y\in F$, squaring gives $d'=x^2+y^2d+2xy\sqrt d$, so $xy=0$; $y=0$ would
make $d'$ a square, so $x=0$ and $d'=y^2d$. With (a) this is the bijection.
(c) Every class in $\Q^\times/\Q^{\times2}$ has a unique squarefree integer representative.
(d) $\sqrt b\notin F(\sqrt a)$: otherwise $F(\sqrt b)=F(\sqrt a)$ (both of degree $2$), and $ab=a^2(b/a)$ would be a square by (b). So $[F(\sqrt a,\sqrt b):F]=4$; the
field is the splitting field of the separable $(x^2-a)(x^2-b)$, hence Galois, and an automorphism is determined by the signs $\epsilon_1,\epsilon_2$, all four
occurring by counting. The three subgroups of order $2$ fix $\sqrt b$, $\sqrt a$, $\sqrt{ab}$ respectively; by \Cref{thm:galois:fundamental} the three quadratic
subfields are $F(\sqrt a)$, $F(\sqrt b)$, $F(\sqrt{ab})$.
\end{proof}

\begin{theorem}[Cyclotomic fields]\label{thm:galois:cyclotomic}
Let $n\ge1$, $\zeta_n=e^{2\pi i/n}$, $\mu_n=\langle\zeta_n\rangle\subset\Qbar$, and $\Phi_n=\prod_{1\le a\le n,\ (a,n)=1}(x-\zeta_n^a)\in\Z[x]$ (\Cref{ch:rings}).
\begin{enumerate}[label=(\alph*)]
\item $\Phi_n$ is irreducible over $\Q$; so $[\Q(\zeta_n):\Q]=\varphi(n)$.
\item $\Q(\zeta_n)/\Q$ is Galois, and $a\mapsto\sigma_a$, where $\sigma_a(\zeta)=\zeta^a$ for all $\zeta\in\mu_n$, is an isomorphism
$(\Z/n)^\times\to\Gal(\Q(\zeta_n)/\Q)$, which does not depend on the choice of a primitive $n$-th root of unity.
\item For $m\mid n$, restriction $\Gal(\Q(\zeta_n)/\Q)\to\Gal(\Q(\zeta_m)/\Q)$ is reduction $(\Z/n)^\times\to(\Z/m)^\times$. For all $m,n$:
$\Q(\zeta_m)\Q(\zeta_n)=\Q(\zeta_{\lcm(m,n)})$ and $\Q(\zeta_m)\cap\Q(\zeta_n)=\Q(\zeta_{\gcd(m,n)})$; $\Q(\zeta_n)=\Q(\zeta_{2n})$ for odd $n$. Every subfield of
$\Q(\zeta_n)$ is Galois over $\Q$ with abelian group.
\item $\Q(\zeta_3)=\Q(\sqrt{-3})$, since $(2\zeta_3+1)^2=-3$; $\Q(\zeta_4)=\Q(i)$; $\Q(\zeta_8)=\Q(i,\sqrt2)$, since $\zeta_8+\zeta_8^{-1}=\sqrt2$, with the three
quadratic subfields $\Q(i)$, $\Q(\sqrt2)$, $\Q(\sqrt{-2})$.
\end{enumerate}
\end{theorem}

\begin{proof}
(a) Let $\zeta$ be a primitive $n$-th root of unity and $g$ its minimal polynomial over $\Q$, a monic factor of the monic $\Phi_n\in\Z[x]$; by Gauss's lemma
(\Cref{ch:rings}) $g\in\Z[x]$ and $\Phi_n=gh$ with $h\in\Z[x]$ monic. We show: \emph{if $p\nmid n$ is prime, then $g(\zeta^p)=0$.} If not, $\zeta^p$ (a primitive
$n$-th root of unity, as $p\nmid n$) is a root of $h$, so $\zeta$ is a root of $h(x^p)$, and $g\mid h(x^p)$ in $\Q[x]$, hence in $\Z[x]$ by Gauss's lemma:
$h(x^p)=g(x)k(x)$ with $k\in\Z[x]$. Reduce modulo $p$: since $\bar h$ has coefficients in $\F_p$, $\bar h(x^p)=\bar h(x)^p$. So $\bar g\bar k=\bar h^p$, and any
irreducible factor $\pi$ of $\bar g$ in $\F_p[x]$ divides $\bar h$. Then $\pi^2$ divides $\bar g\bar h=\bar\Phi_n$, which divides $x^n-1$ in $\F_p[x]$; but $x^n-1$
is separable over $\F_p$ because $p\nmid n$. This contradiction proves the claim. Now every primitive root is $\zeta^a$ with $(a,n)=1$, $a=p_1\cdots p_k$ a product
of primes not dividing $n$; applying the claim successively to $\zeta,\zeta^{p_1},\zeta^{p_1p_2},\dots$ (all primitive, all roots of the same $g$) gives
$g(\zeta^a)=0$. So all $\varphi(n)$ roots of $\Phi_n$ are roots of $g$, and $g=\Phi_n$.
(b) $\Q(\zeta_n)$ contains all primitive $n$-th roots of unity (powers of $\zeta_n$), so it is the splitting field of the separable $\Phi_n$, Galois of degree
$\varphi(n)$. An automorphism $\sigma$ sends $\zeta_n$ to a root of $\Phi_n$, so $\sigma(\zeta_n)=\zeta_n^{a}$ with $a\in(\Z/n)^\times$ determined by $\sigma$, and
then $\sigma(\zeta_n^k)=\zeta_n^{ak}$: $\sigma=\sigma_a$ acts on all of $\mu_n$ by $\zeta\mapsto\zeta^a$, which does not mention $\zeta_n$. As
$\sigma_a\sigma_b(\zeta)=\sigma_a(\zeta^b)=\zeta^{ab}$, $a\mapsto\sigma_a$ is an injective homomorphism, bijective by counting.
(c) $\zeta_m=\zeta_n^{n/m}$, so $\sigma_a(\zeta_m)=\zeta_m^a=\zeta_m^{a\bmod m}$. Let $g=\gcd(m,n)$, $s=\lcm(m,n)$ and $um+vn=g$ with $u,v\in\Z$. Then
$\zeta_m^v\zeta_n^u=e^{2\pi i(vn+um)/mn}=e^{2\pi ig/mn}=\zeta_s$, so $\Q(\zeta_m)\Q(\zeta_n)\supseteq\Q(\zeta_s)$, and the converse inclusion is clear. By
\Cref{prop:galois:translation}(a), $[\Q(\zeta_m)\cap\Q(\zeta_n):\Q]=[\Q(\zeta_m):\Q]\,[\Q(\zeta_n):\Q]/[\Q(\zeta_s):\Q]=\varphi(m)\varphi(n)/\varphi(s)$, and
$\varphi(m)\varphi(n)=\varphi(s)\varphi(g)$ (from $\varphi(k)=k\prod_{p\mid k}(1-1/p)$ and $mn=sg$, the primes dividing $s$ being those dividing $m$ or $n$, those
dividing $g$ those dividing both). So the intersection has degree $\varphi(g)$ and contains $\Q(\zeta_g)$: they are equal. For odd $n$, $\Q(\zeta_n)\subseteq
\Q(\zeta_{2n})$ and $\varphi(2n)=\varphi(n)$. Subgroups of the abelian group $(\Z/n)^\times$ are normal, so \Cref{thm:galois:fundamental}(d) applies.
(d) $(2\zeta_3+1)^2=4(\zeta_3^2+\zeta_3+1)-3=-3$, and both fields have degree $2$. $\zeta_8+\zeta_8^{-1}=2\cos(\pi/4)=\sqrt2$; $\Q(i,\sqrt2)\subseteq\Q(\zeta_8)$ and both
have degree $4$. The quadratic subfields follow from \Cref{prop:galois:quadratic}(d) with $a=-1$, $b=2$.
\end{proof}

\begin{proposition}[The field $\Q(\sqrt{\ell^*})$]\label{prop:galois:sqrt-lstar}
Let $\ell$ be an odd prime and $\ell^*=(-1)^{(\ell-1)/2}\ell$.
\begin{enumerate}[label=(\alph*)]
\item The Gauss sum $g=\sum_{a=1}^{\ell-1}\bigl(\frac a\ell\bigr)\zeta_\ell^a$ satisfies $g^2=\ell^*$ and $\sigma_a(g)=\bigl(\frac a\ell\bigr)g$ for $a\in(\Z/\ell)^\times$.
In particular $\sqrt{\ell^*}=\pm g\in\Q(\zeta_\ell)$.
\item $\Q(\sqrt{\ell^*})$ is the unique quadratic subfield of $\Q(\zeta_\ell)$, the fixed field of the subgroup of squares of $(\Z/\ell)^\times$, and
$\sigma_a(\sqrt{\ell^*})=\bigl(\frac a\ell\bigr)\sqrt{\ell^*}$. In terms of characters of $G_\Q$ (\Cref{def:galois:cyclotomic-character}): for $\sigma\in G_\Q$,
\[
\chi_{\ell^*}(\sigma)=\Bigl(\frac{\bar\chi_\ell(\sigma)}{\ell}\Bigr)=\bar\chi_\ell(\sigma)^{(\ell-1)/2},
\]
the last equality in $\F_\ell^\times\supseteq\{\pm1\}$ (Euler's criterion). Briefly, $\chi_{\ell^*}=\bar\chi_\ell^{(\ell-1)/2}$.
\item For $\ell=3$: $\Q(\sqrt{-3})=\Q(\zeta_3)$ and $\chi_{-3}=\bar\chi_3$. For $\ell\ge5$: $\chi_{\ell^*}\ne\bar\chi_\ell$, since $\bar\chi_\ell$ has order $\ell-1>2$.
\item Every quadratic field $\Q(\sqrt d)$, $d$ squarefree, lies in $\Q(\zeta_{|D|})$, where $D=d$ if $d\equiv1\pmod4$ and $D=4d$ otherwise is its discriminant.
In particular, for a finite set $S$ of odd primes, $\Q\bigl(\sqrt{\textstyle\prod_{q\in S}q^*}\bigr)\subseteq\Q(\zeta_{\prod_{q\in S}q})$.
\end{enumerate}
\end{proposition}

\begin{proof}
(a) $g^2=\ell^*$ is \Cref{ch:rings} (the quadratic Gauss sum). For $a$ prime to $\ell$, $\sigma_a(g)=\sum_b\bigl(\frac b\ell\bigr)\zeta_\ell^{ab}=\bigl(\frac a\ell\bigr)
\sum_b\bigl(\frac{ab}\ell\bigr)\zeta_\ell^{ab}=\bigl(\frac a\ell\bigr)g$, as $b\mapsto ab$ permutes $(\Z/\ell)^\times$ and $\bigl(\frac a\ell\bigr)^2=1$.
(b) $(\Z/\ell)^\times$ is cyclic of even order $\ell-1$ (\Cref{ch:rings}), so it has exactly one subgroup of index $2$, the squares; by
\Cref{thm:galois:fundamental} $\Q(\zeta_\ell)$ has exactly one quadratic subfield, the fixed field of the squares. By (a), $g$ is fixed exactly by the $\sigma_a$
with $\bigl(\frac a\ell\bigr)=1$, and $g\notin\Q$ (it is moved by $\sigma_a$ for a non-residue $a$), so that subfield is $\Q(g)=\Q(\sqrt{\ell^*})$. For
$\sigma\in G_\Q$, $\sigma|_{\Q(\zeta_\ell)}=\sigma_a$ with $a\equiv\bar\chi_\ell(\sigma)$ (\Cref{def:galois:cyclotomic-character}), so
$\chi_{\ell^*}(\sigma)=\sigma(g)/g=\bigl(\frac a\ell\bigr)$, and $\bigl(\frac a\ell\bigr)\equiv a^{(\ell-1)/2}\pmod\ell$.
(c) $\Q(\sqrt{-3})=\Q(\zeta_3)$ is \Cref{thm:galois:cyclotomic}(d); with $\ell=3$, $(\ell-1)/2=1$. For $\ell\ge5$, $\bar\chi_\ell$ is surjective onto the cyclic group
$\F_\ell^\times$ of order $\ell-1\ge4$ (\Cref{prop:galois:cyclotomic-character}), while $\chi_{\ell^*}^2=1$.
(d) Write $d=\pm2^e\prod_{i}q_i$ with $e\in\{0,1\}$ and distinct odd primes $q_i$; then $d=\epsilon2^e\prod_iq_i^*$ with $\epsilon=\pm1$, and $\sqrt{q_i^*}\in\Q(\zeta_{q_i})$
by (b). If $d\equiv1\pmod4$, then $e=0$ and $\epsilon=1$, because $\prod q_i^*\equiv1\pmod4$ while $-\prod q_i^*\equiv3\pmod 4$; so $\sqrt d\in\Q(\zeta_{|d|})$. If
$d\equiv3\pmod 4$, $e=0$ and $\sqrt d=i\prod\sqrt{q_i^*}$ up to sign, in $\Q(\zeta_{4|d|})$ ($i\in\Q(\zeta_4)$, \Cref{thm:galois:cyclotomic}(c)). If $d$ is even,
$\sqrt{\pm2}\in\Q(\zeta_8)$ and $\sqrt d\in\Q(\zeta_{8|d|/2})=\Q(\zeta_{4|d|})$. The last statement is the case $d=\prod_{q\in S}q^*\equiv1\pmod4$, up to the
squarefree representative, which here is $d$ itself.
\end{proof}

\begin{paperlink}
The hypothesis of Theorem thm:main concerns $\Gal(\Qbar/\Q(\sqrt{\ell^*}))=\ker\chi_{\ell^*}$, the subgroup of index $2$ of $G_\Q$ on which $\bar\chi_\ell$ takes square
values. For $\ell=3$ it is $\Gal(\Qbar/\Q(\zeta_3))=\ker\bar\chi_3$; this coincidence $\chi_{-3}=\bar\chi_3$ is why Lemma lem:h0 can phrase the exceptional case both as
$\bar\rho\simeq\bar\rho\otimes\bar\chi_3$ and as $\bar\rho\simeq\bar\rho\otimes\chi_{-3}$. For $\ell\ge5$ the two characters differ, and the paper's revised
Section 1.3 correctly states the criterion with $\bar\chi_\ell$ (the original version used $\chi_{\ell^*}$; see \Cref{prop:galois:twist}(c) and
\Cref{prop:galois:index-two}).
\end{paperlink}

\begin{example}[The splitting field of $x^3-2$]\label{ex:galois:x3minus2}
Let $\alpha=2^{1/3}\in\R$, $\zeta=\zeta_3$ and $L=\Q(\alpha,\zeta)$, the splitting field of $x^3-2$ (roots $\alpha$, $\zeta\alpha$, $\zeta^2\alpha$). As $[\Q(\alpha):\Q]=3$
and $[\Q(\zeta):\Q]=2$ divide $[L:\Q]\le6$, $[L:\Q]=6$. Let $\sigma\in\Gal(L/\Q)$ fix $\zeta$ and send $\alpha\mapsto\zeta\alpha$, and let $\tau$ be complex conjugation,
which fixes $\alpha$ and sends $\zeta\mapsto\zeta^2$ (both exist by \Cref{thm:galois:splitting-field}(b): the extensions of $\Q(\zeta)\to\Q(\zeta)$, identity or
conjugation, to $L$ are determined by the image of $\alpha$, any of the three roots). Then $\sigma^3=\tau^2=1$ and $\tau\sigma\tau(\alpha)=\tau(\zeta\alpha)=\zeta^2\alpha=
\sigma^{-1}(\alpha)$, $\tau\sigma\tau(\zeta)=\zeta=\sigma^{-1}(\zeta)$, so $\tau\sigma\tau=\sigma^{-1}$ and $\Gal(L/\Q)\simeq S_3$, acting faithfully on the three
roots. One computes $\sigma\tau(\zeta^2\alpha)=\zeta^4\cdot\zeta\alpha=\zeta^2\alpha$ and $\sigma^2\tau(\zeta\alpha)=\zeta^2\cdot\zeta^2\alpha=\zeta\alpha$; so each subgroup of
order $2$ is the stabiliser of the root it fixes, and the correspondence of \Cref{thm:galois:fundamental} is:
\begin{center}
\begin{tabular}{lll}
\toprule
subgroup & fixed field & degree over $\Q$\\
\midrule
$\{1\}$ & $L$ & $6$\\
$\langle\tau\rangle$ & $\Q(\alpha)$ & $3$\\
$\langle\sigma\tau\rangle$ & $\Q(\zeta^2\alpha)$ & $3$\\
$\langle\sigma^2\tau\rangle$ & $\Q(\zeta\alpha)$ & $3$\\
$\langle\sigma\rangle=A_3$ & $\Q(\zeta)=\Q(\sqrt{-3})$ & $2$\\
$S_3$ & $\Q$ & $1$\\
\bottomrule
\end{tabular}
\end{center}
The three cubic fields are conjugate (\Cref{thm:galois:fundamental}(c)) and not normal; $\Q(\sqrt{-3})$ is the only normal subfield other than $\Q$ and $L$, cut out by the sign of
the permutation of the roots, and $\disc(x^3-2)=-108=-3\cdot6^2$ (\Cref{prop:galois:discriminant}). This is the pattern of Proposition prop:ns in miniature: a
group of permutations of roots, and the quadratic subfield cut out by the sign, which there is $\det\bar\rho_{E,3}=\chi_{-3}$. The group is also the affine group
$\Z/3\rtimes(\Z/3)^\times$ of \Cref{prop:galois:kummer-affine}.
\end{example}

\section[Hilbert 90, Kummer theory, Artin-Schreier]{Cyclic extensions: Hilbert 90, Kummer, Artin-Schreier}\label{sec:galois:kummer}

Kummer theory describes the abelian extensions of exponent dividing $n$ of a field containing the $n$-th roots of unity: they are obtained by
adjoining $n$-th roots, and the Galois action on an $n$-th root $\alpha$ of $a$ is recorded by the character $\sigma\mapsto\sigma(\alpha)/\alpha$. The paper uses
it twice. In Lemma lem:tamagawa-exp the $\ell^t$-torsion of a Tate curve over $\Q_q^{\mathrm{ur}}$ is generated by $\mu_{\ell^t}$ and an $\ell^t$-th root of the Tate
parameter $q_E$, and ``by Kummer theory'' inertia acts trivially iff $q_E$ is an $\ell^t$-th power; in Lemma lem:tamfactor the Galois action on $T_\ell E$ is given by
the Kummer cocycle $\kappa$ of $q_E$, whose unramified part is the class of the unit part of $q_E$ in $\F_q^\times\otimes\Zl$. The field-theoretic statements are
proved here; the $\ell$-adic cocycle is \Cref{prop:galois:l-adic-kummer}, and the cohomological form $H^1(K,\mu_n)\simeq K^\times/K^{\times n}$ is \Cref{ch:galcoh}.

\begin{definition}[Norm and trace]\label{def:galois:norm-trace}
Let $L/F$ be finite separable of degree $n$ and $\sigma_1,\dots,\sigma_n$ its $F$-embeddings into $\Fbar$. For $a\in L$,
$\Nm_{L/F}(a)=\prod_i\sigma_i(a)$ and $\Tr_{L/F}(a)=\sum_i\sigma_i(a)$.
\end{definition}

\begin{proposition}[Norm and trace]\label{prop:galois:norm-trace}
Let $L/F$ be finite separable of degree $n$ and $a\in L$.
\begin{enumerate}[label=(\alph*)]
\item The characteristic polynomial of the $F$-linear map $m_a\colon x\mapsto ax$ of $L$ is $m_{a,F}^{[L:F(a)]}$; hence $\Nm_{L/F}(a)=\det m_a$ and
$\Tr_{L/F}(a)=\tr m_a$, and both lie in $F$.
\item $\Nm_{L/F}$ is multiplicative and $\Tr_{L/F}$ is $F$-linear; for $L/F$ Galois, $\Nm_{L/F}(a)=\prod_{\sigma\in\Gal(L/F)}\sigma(a)$.
\item For $F\subseteq K\subseteq L$: $\Nm_{L/F}=\Nm_{K/F}\circ\Nm_{L/K}$ and $\Tr_{L/F}=\Tr_{K/F}\circ\Tr_{L/K}$.
\item $\Tr_{L/F}$ is not identically zero.
\end{enumerate}
\end{proposition}

\begin{proof}
(a) Let $K=F(a)$, $d=[K:F]$, $m=[L:K]$, and $y_1,\dots,y_m$ a $K$-basis of $L$. In the $F$-basis $(a^iy_j)$ of $L$ (\Cref{prop:galois:tower}) the matrix of $m_a$ is
block diagonal with $m$ copies of the companion matrix of $m_{a,F}$, whose characteristic polynomial is $m_{a,F}$. On the other side, each of the $d$ embeddings
of $K$ (sending $a$ to the $d$ distinct roots $a_1,\dots,a_d$ of $m_{a,F}$) has exactly $m$ extensions to $L$ (\Cref{prop:galois:separability}(c), (d)), so
$\prod_i\sigma_i(a)=(\prod_ka_k)^m$ and $\sum_i\sigma_i(a)=m\sum_ka_k$, which are the determinant and the trace of the block matrix.
(b) Multiplicativity and linearity are clear from the definition (or from (a)). If $L/F$ is Galois, the embeddings are the elements of $\Gal(L/F)$.
(c) Every $F$-embedding of $L$ is obtained, exactly once, by choosing an $F$-embedding $\tau$ of $K$, extending it to an automorphism $\tilde\tau$ of $\Fbar$
(\Cref{thm:galois:algebraic-closure}(d)), and composing with one of the $K$-embeddings $\rho_1,\dots,\rho_m$ of $L$; so $\Nm_{L/F}(a)=\prod_\tau\prod_j\tilde\tau
\rho_j(a)=\prod_\tau\tau\bigl(\Nm_{L/K}(a)\bigr)$, and similarly for the trace.
(d) $\Tr_{L/F}=\sum_i\sigma_i$ is a nontrivial combination of distinct embeddings, nonzero by \Cref{thm:galois:dedekind}.
\end{proof}

\begin{theorem}[Hilbert's Theorem 90]\label{thm:galois:hilbert90}
\begin{enumerate}[label=(\alph*)]
\item Let $L/F$ be finite Galois with group $G$, and let $(a_\sigma)_{\sigma\in G}$ be elements of $L^\times$ with $a_{\sigma\tau}=a_\sigma\,\sigma(a_\tau)$ for all
$\sigma,\tau$. Then there is $c\in L^\times$ with $a_\sigma=\sigma(c)/c$ for all $\sigma$.
\item Let $L/F$ be cyclic of degree $n$ with $\Gal(L/F)=\langle\sigma\rangle$. For $b\in L^\times$: $\Nm_{L/F}(b)=1$ iff $b=c/\sigma(c)$ for some $c\in L^\times$.
\item In the same situation, for $b\in L$: $\Tr_{L/F}(b)=0$ iff $b=c-\sigma(c)$ for some $c\in L$.
\end{enumerate}
\end{theorem}

\begin{proof}
(a) By \Cref{thm:galois:dedekind} the map $\sum_{\tau\in G}a_\tau\tau$ is not identically zero; choose $\theta$ with $c_0=\sum_\tau a_\tau\tau(\theta)\ne0$. For
$\sigma\in G$, $\sigma(c_0)=\sum_\tau\sigma(a_\tau)\,\sigma\tau(\theta)=\sum_\tau a_\sigma^{-1}a_{\sigma\tau}\,\sigma\tau(\theta)=a_\sigma^{-1}c_0$. So $a_\sigma=c_0/\sigma(c_0)=
\sigma(c)/c$ with $c=c_0^{-1}$.
(b) If $b=c/\sigma(c)$, then $\Nm(b)=\Nm(c)/\Nm(\sigma c)=1$, the norm being the product over the group. Conversely, put $a_{\sigma^i}=b\,\sigma(b)\cdots\sigma^{i-1}(b)$
for $0\le i<n$ ($a_1=1$). The cocycle relation $a_{\sigma^{i+j}}=a_{\sigma^i}\sigma^i(a_{\sigma^j})$ holds when $i+j<n$ by definition, and when $i+j\ge n$ because
the product $b\sigma(b)\cdots\sigma^{i+j-1}(b)$ then contains the full norm $\prod_{k=0}^{n-1}\sigma^k(b)=1$ once. By (a), $b=a_\sigma=\sigma(d)/d$ for some $d$; put
$c=d^{-1}$.
(c) If $b=c-\sigma(c)$, $\Tr(b)=0$. Conversely, choose $\theta$ with $\Tr\theta\ne0$ (\Cref{prop:galois:norm-trace}(d)), put $s_i=b+\sigma(b)+\dots+\sigma^i(b)$
(so $s_{n-1}=\Tr b=0$ and $\sigma(s_i)=s_{i+1}-b$) and $c=(\Tr\theta)^{-1}\sum_{i=0}^{n-2}s_i\sigma^{i+1}(\theta)$. Then
\[
(\Tr\theta)\,\sigma(c)=\sum_{i=0}^{n-2}(s_{i+1}-b)\sigma^{i+2}(\theta)=\sum_{j=1}^{n-2}s_j\sigma^{j+1}(\theta)-b\sum_{j=2}^{n}\sigma^{j}(\theta),
\]
using $s_{n-1}=0$. Subtracting from $(\Tr\theta)c=s_0\sigma(\theta)+\sum_{j=1}^{n-2}s_j\sigma^{j+1}(\theta)$ and using $s_0=b$ gives $(\Tr\theta)(c-\sigma(c))=b\sum_{j=1}^n
\sigma^j(\theta)=b\Tr\theta$.
\end{proof}

\begin{proposition}[Kummer characters]\label{prop:galois:kummer-character}
Let $n\ge1$, $\operatorname{char}F\nmid n$, and assume $\mu_n\subset F$ (so $\mu_n$ is cyclic of order $n$). Let $a\in F^\times$, $\alpha\in\Fbar$ with
$\alpha^n=a$, and $L/F$ a Galois extension containing $\alpha$.
\begin{enumerate}[label=(\alph*)]
\item $\chi_a\colon\Gal(L/F)\to\mu_n$, $\sigma\mapsto\sigma(\alpha)/\alpha$, is a homomorphism that does not depend on the choice of $\alpha$; its kernel is
$\Gal(L/F(\alpha))$. Hence $F(\alpha)/F$ is Galois with cyclic group $\simeq\chi_a(\Gal(L/F))$, of order $d$ equal to the order of the class of $a$ in
$F^\times/F^{\times n}$; in particular $d\mid n$.
\item $\chi_{ab}=\chi_a\chi_b$, and $\chi_a=1$ iff $a\in F^{\times n}$. So if $L$ contains an $n$-th root of every element of $F^\times$ (for instance
$L=F^{\mathrm{sep}}$), $a\mapsto\chi_a$ is an injective homomorphism $F^\times/F^{\times n}\to\Hom(\Gal(L/F),\mu_n)$.
\item The minimal polynomial of $\alpha$ is $x^d-\alpha^d$, and $x^n-a$ is irreducible over $F$ iff $a\notin F^{\times p}$ for every prime $p\mid n$.
\end{enumerate}
Without the assumption $\mu_n\subset F$ the irreducibility criterion is Capelli's theorem: $x^n-a$ is irreducible iff $a\notin F^{\times p}$ for every prime
$p\mid n$ and $a\notin-4F^{\times4}$ when $4\mid n$ (for example $x^4+4=(x^2+2x+2)(x^2-2x+2)$). We do not prove this; see \cite[Chapter VI, Theorem 9.1]{Lang}.
\end{proposition}

\begin{proof}
(a) The polynomial $x^n-a$ is separable ($a\ne0$ and $n\ne0$ in $F$), with the $n$ roots $\zeta\alpha$, $\zeta\in\mu_n$, all in $F(\alpha)$; so $F(\alpha)$ is its
splitting field, Galois over $F$. For $\sigma\in\Gal(L/F)$, $\sigma(\alpha)^n=\sigma(a)=a$, so $\chi_a(\sigma)=\sigma(\alpha)/\alpha\in\mu_n$. Another root is
$\alpha'=\zeta\alpha$ with $\zeta\in\mu_n\subset F$ fixed by $\sigma$, so $\sigma(\alpha')/\alpha'=\sigma(\alpha)/\alpha$. Since $\chi_a(\tau)\in F$,
\[
\chi_a(\sigma\tau)=\frac{\sigma(\tau(\alpha))}{\alpha}=\frac{\sigma(\chi_a(\tau)\alpha)}{\alpha}=\chi_a(\tau)\frac{\sigma(\alpha)}\alpha=\chi_a(\sigma)\chi_a(\tau).
\]
The kernel consists of the $\sigma$ fixing $\alpha$, that is $\Gal(L/F(\alpha))$, and
\[
\Gal(F(\alpha)/F)\simeq\Gal(L/F)/\Gal(L/F(\alpha))\simeq\chi_a(\Gal(L/F)),
\]
a subgroup of the cyclic group $\mu_n$, of some order $d\mid n$. Then $\chi_a(\sigma)^d=1$ for all $\sigma$, so $\alpha^d$ is fixed by $\Gal(L/F)$, $\alpha^d\in F$, and
$a^d=(\alpha^d)^n\in F^{\times n}$. Conversely, if $a^k=b^n$ with $b\in F^\times$, then $(\alpha^k/b)^n=1$, so $\alpha^k\in\mu_nb\subseteq F$ and $\chi_a(\sigma)^k=
\sigma(\alpha^k)/\alpha^k=1$ for all $\sigma$, whence $d\mid k$. So $d$ is the order of $a$ modulo $F^{\times n}$.
(b) With $\alpha^n=a$, $\beta^n=b$, $(\alpha\beta)^n=ab$ and $\sigma(\alpha\beta)/\alpha\beta=\chi_a(\sigma)\chi_b(\sigma)$. $\chi_a=1$ iff $\alpha\in F$ iff $a\in F^{\times n}$
(if $a=b^n$, then $\alpha\in\mu_nb\subseteq F$).
(c) $\alpha$ is a root of $x^d-\alpha^d\in F[x]$, of degree $d=[F(\alpha):F]$. So $x^n-a$ is irreducible iff $d=n$, iff the class of $a$ in the group $F^\times/F^{\times n}$
of exponent $n$ has order exactly $n$, iff $a^{n/p}\notin F^{\times n}$ for every prime $p\mid n$. Finally $a^{n/p}\in F^{\times n}$ iff $a\in F^{\times p}$: if $a=b^p$,
$a^{n/p}=b^n$; if $a^{n/p}=b^n$, then $(a/b^p)^{n/p}=1$, so $a/b^p\in\mu_{n/p}=\{\zeta^p:\zeta\in\mu_n\}$ ($\mu_n$ being cyclic of order $n$), and $a$ is a $p$-th power.
\end{proof}

\begin{theorem}[Cyclic extensions]\label{thm:galois:kummer-cyclic}
Let $\operatorname{char}F\nmid n$ and $\mu_n\subset F$. Every cyclic extension $L/F$ of degree $n$ is of the form $L=F(c)$ with $c^n\in F^\times$.
\end{theorem}

\begin{proof}
Let $\Gal(L/F)=\langle\sigma\rangle$ and $\zeta$ a primitive $n$-th root of unity. As $\zeta\in F$, $\Nm_{L/F}(\zeta^{-1})=\zeta^{-n}=1$, so by
\Cref{thm:galois:hilbert90}(b) $\zeta^{-1}=c/\sigma(c)$ for some $c\in L^\times$: $\sigma(c)=\zeta c$. Then $\sigma^i(c)=\zeta^ic$, $0\le i<n$, are $n$ distinct conjugates
of $c$, so $[F(c):F]=n$ and $F(c)=L$; and $\sigma(c^n)=c^n$, so $c^n\in F$.
\end{proof}

\begin{theorem}[Kummer theory]\label{thm:galois:kummer}
Let $\operatorname{char}F\nmid n$ and $\mu_n\subset F$. For a subgroup $F^{\times n}\subseteq\Delta\subseteq F^\times$ let $F_\Delta=F(\Delta^{1/n})$ be the subfield of
$F^{\mathrm{sep}}$ generated by all $\alpha$ with $\alpha^n\in\Delta$.
\begin{enumerate}[label=(\alph*)]
\item If $\Delta/F^{\times n}$ is finite, $F_\Delta/F$ is a finite abelian extension of exponent dividing $n$, and
\[
\Gal(F_\Delta/F)\times\Delta/F^{\times n}\to\mu_n,\qquad(\sigma,a)\mapsto\sigma(\alpha)/\alpha\quad(\alpha^n=a),
\]
is a well defined bilinear pairing which is perfect: it induces isomorphisms $\Gal(F_\Delta/F)\simeq\Hom(\Delta/F^{\times n},\mu_n)$ and
$\Delta/F^{\times n}\simeq\Hom(\Gal(F_\Delta/F),\mu_n)$. In particular $[F_\Delta:F]=[\Delta:F^{\times n}]$.
\item $\Delta\mapsto F_\Delta$ is an inclusion-preserving bijection from the groups $\Delta$ with $\Delta/F^{\times n}$ finite to the finite abelian extensions of $F$ inside
$F^{\mathrm{sep}}$ of exponent dividing $n$; the inverse is $L\mapsto L^{\times n}\cap F^\times$. The cyclic extensions of degree $d$ correspond to the $\Delta$ with
$\Delta/F^{\times n}$ cyclic of order $d$.
\end{enumerate}
The form without finiteness, $F^\times/F^{\times n}\simeq\Hom_{\mathrm{cont}}(G_F,\mu_n)$, is \Cref{cor:galois:kummer-absolute}.
\end{theorem}

\begin{proof}
(a) Let $a_1,\dots,a_r$ represent generators of $\Delta/F^{\times n}$ and $\alpha_i^n=a_i$. If $\alpha^n=a\in\Delta$, write $a=b^n\prod a_i^{k_i}$ with $b\in F^\times$; then
$\alpha/(b\prod\alpha_i^{k_i})\in\mu_n\subset F$, so $F_\Delta=F(\alpha_1,\dots,\alpha_r)$, the splitting field of the separable $\prod_i(x^n-a_i)$, finite Galois. The
pairing is well defined by \Cref{prop:galois:kummer-character}(a) (independence of $\alpha$) and (b) ($\chi_a$ depends only on $a$ modulo $F^{\times n}$), and
bilinear by the homomorphism property in each variable. If $\sigma$ pairs trivially with all $a$, it fixes every $\alpha_i$, so $\sigma=1$; if $a$ pairs trivially
with all $\sigma$, then $\chi_a=1$ and $a\in F^{\times n}$. So $G=\Gal(F_\Delta/F)\to\Hom(\Delta/F^{\times n},\mu_n)$ and $\Delta/F^{\times n}\to\Hom(G,\mu_n)$ are
injective. The first shows that $G$ is abelian of exponent dividing $n$. For a finite abelian group $A$ of exponent dividing $n$, $\#\Hom(A,\mu_n)=\#A$ (write $A$ as
a product of cyclic groups $\Z/m$ with $m\mid n$, \Cref{ch:groups}, and note $\Hom(\Z/m,\mu_n)\simeq\mu_m$). Hence $\#G\le\#(\Delta/F^{\times n})\le\#G$, both injections
are bijections, and $[F_\Delta:F]=\#G=[\Delta:F^{\times n}]$.

(b) Let $L/F$ be finite abelian of exponent dividing $n$ and $\Delta_L=L^{\times n}\cap F^\times$. If $\beta\in L$ with $\beta^n=a\in F$, all $n$-th roots of $a$ are
$\zeta\beta\in L$; so $F_{\Delta_L}\subseteq L$, and in particular $\Delta_L/F^{\times n}$ is finite by (a) (it injects into $\Hom(\Gal(F_{\Delta_L}/F),\mu_n)$, the
injectivity argument of (a) not using finiteness). Conversely, $A=\Gal(L/F)$ is a finite abelian group, and the intersection of the subgroups $B$ with $A/B$ cyclic
is trivial (by the structure theorem, \Cref{ch:groups}, for $1\ne x\in A$ some projection to a cyclic factor is nonzero at $x$). By
\Cref{thm:galois:fundamental}(e), $L$ is the compositum of the fixed fields $L^B$, each cyclic of some degree $e\mid n$, hence $L^B=F(\beta)$ with $\beta^e\in F$
(\Cref{thm:galois:kummer-cyclic}, as $\mu_e\subseteq\mu_n\subset F$). Then $\beta^n=(\beta^e)^{n/e}\in F$, so $\beta\in F_{\Delta_L}$; hence $L=F_{\Delta_L}$. Finally
$\Delta_{F_\Delta}=\Delta$: the inclusion $\supseteq$ is clear; if $a\in F^\times\cap F_\Delta^{\times n}$, the group $\Delta'$ generated by $\Delta$ and $a$ has
$F_{\Delta'}=F_\Delta$, so $[\Delta':F^{\times n}]=[\Delta:F^{\times n}]$ by (a) and $a\in\Delta$. The two maps preserve inclusions. By (a), $\Gal(F_\Delta/F)$ is cyclic
of order $d$ iff its dual $\Delta/F^{\times n}$ is.
\end{proof}

\begin{proposition}[The affine shape without roots of unity]\label{prop:galois:kummer-affine}
Let $\operatorname{char}F\nmid n$, $a\in F^\times$, $\alpha^n=a$, $L=F(\mu_n,\alpha)$, and fix a primitive $n$-th root of unity $\zeta_n$. For $\sigma\in\Gal(L/F)$
define $u(\sigma)\in(\Z/n)^\times$ and $b(\sigma)\in\Z/n$ by $\sigma(\zeta_n)=\zeta_n^{u(\sigma)}$ and $\sigma(\alpha)=\zeta_n^{b(\sigma)}\alpha$. Then
\[
\sigma\longmapsto\begin{pmatrix}u(\sigma)&b(\sigma)\\0&1\end{pmatrix}
\]
is an injective homomorphism from $\Gal(L/F)$ to the affine group $\bigl\{\sm ub01:u\in(\Z/n)^\times,\ b\in\Z/n\bigr\}\simeq\Z/n\rtimes(\Z/n)^\times$. Equivalently,
$b(\sigma\tau)=b(\sigma)+u(\sigma)b(\tau)$: $b$ is a $1$-cocycle for the action of $\Gal(L/F)$ on $\Z/n$ through $u$. Its restriction to $\Gal(L/F(\mu_n))$ (where
$u=1$) is the Kummer character $\chi_a$, read in $\Z/n$ through $\zeta_n$.
\end{proposition}

\begin{proof}
$L$ is the splitting field of the separable $x^n-a$ (its roots $\zeta\alpha$ generate $\mu_n$ and $\alpha$), so it is Galois. We have
$\sigma\tau(\zeta_n)=\sigma(\zeta_n^{u(\tau)})=\zeta_n^{u(\sigma)u(\tau)}$ and
$\sigma\tau(\alpha)=\sigma(\zeta_n^{b(\tau)}\alpha)=\zeta_n^{u(\sigma)b(\tau)}\zeta_n^{b(\sigma)}\alpha$, which is the matrix product
$\sm{u(\sigma)}{b(\sigma)}01\sm{u(\tau)}{b(\tau)}01$. The map is injective since $\sigma$ is determined by $\sigma(\zeta_n)$ and $\sigma(\alpha)$.
\end{proof}

For $x^3-2$ over $\Q$ (\Cref{ex:galois:x3minus2}) the Galois group has order $6$ and embeds in the affine group of order $6$, so
$\Gal\simeq\Z/3\rtimes(\Z/3)^\times\simeq S_3$. The $\ell$-adic limit of this matrix, $\sm{\chi_\ell}{\kappa_a}01$, is \Cref{prop:galois:l-adic-kummer}; twisted by an
unramified character $\varepsilon$ of order at most $2$ it becomes the matrix $\sm{\chi_\ell\varepsilon}{\varepsilon\kappa}0\varepsilon$ of the proof of Lemma lem:tamfactor
(\Cref{ch:reduction}).

\begin{proposition}[Kummer theory over a finite field]\label{prop:galois:kummer-finite-field}
Let $q$ be a power of a prime $p$, $\ell\ne p$ a prime, $v=v_\ell(q-1)$, and $u\in\F_q^\times$.
\begin{enumerate}[label=(\alph*)]
\item For $1\le k\le v$, $\mu_{\ell^k}\subset\F_q$, and the Kummer character $\chi_u$ (with $n=\ell^k$) satisfies $\chi_u(\Frob_q)=u^{(q-1)/\ell^k}$. Hence
$u\in(\F_q^\times)^{\ell^k}$ iff $u^{(q-1)/\ell^k}=1$, and $u\mapsto u^{(q-1)/\ell^v}$ induces an isomorphism of cyclic groups of order $\ell^v$
\[
\F_q^\times/(\F_q^\times)^{\ell^v}\ \xrightarrow{\ \sim\ }\ \mu_{\ell^v}(\F_q);
\]
the left side is also $\F_q^\times\otimes\Zl$.
\item Let $(\alpha_k)_{k\ge0}$ be a compatible system of $\ell$-power roots of $u$ in $\Fbar_q$ ($\alpha_0=u$, $\alpha_{k+1}^\ell=\alpha_k$), and put
$\kappa_u(\Frob_q)=(\Frob_q(\alpha_k)/\alpha_k)_k=(\alpha_k^{q-1})_k\in\Zl(1)=\varprojlim_k\mu_{\ell^k}(\Fbar_q)$. Another compatible system changes $\kappa_u(\Frob_q)$ by an
element of $(q-1)\Zl(1)=\ell^v\Zl(1)$. The class of $\kappa_u(\Frob_q)$ in $\Zl(1)/\ell^v\Zl(1)\simeq\mu_{\ell^v}(\F_q)$ is $u^{(q-1)/\ell^v}$, and for $k\le v$ it lies in
$\ell^k\Zl(1)/\ell^v\Zl(1)$ iff $u\in(\F_q^\times)^{\ell^k}$.
\end{enumerate}
\end{proposition}

\begin{proof}
(a) $\F_q^\times$ is cyclic of order $q-1$ (\Cref{ch:rings}), divisible by $\ell^k$, so it contains $\mu_{\ell^k}$. If $\alpha^{\ell^k}=u$, then
$\chi_u(\Frob_q)=\alpha^q/\alpha=\alpha^{q-1}=(\alpha^{\ell^k})^{(q-1)/\ell^k}=u^{(q-1)/\ell^k}$. As $\Frob_q$ generates $\Gal(\F_q(\alpha)/\F_q)$
(\Cref{thm:galois:finite-fields}), $\chi_u=1$ iff $\chi_u(\Frob_q)=1$, and $\chi_u=1$ iff $u\in(\F_q^\times)^{\ell^k}$ (\Cref{prop:galois:kummer-character}(b)). The map
$u\mapsto u^{(q-1)/\ell^v}$ is a homomorphism to $\mu_{\ell^v}(\F_q)$ with kernel $(\F_q^\times)^{\ell^v}$ (case $k=v$), surjective because a generator $g$ of $\F_q^\times$
maps to an element of order $\ell^v$. Writing $q-1=\ell^vm$ with $\ell\nmid m$, $\F_q^\times\simeq\Z/\ell^v\times\Z/m$ and $\F_q^\times\otimes\Zl\simeq\Z/\ell^v$.
(b) $(\alpha_k^{q-1})^{\ell^k}=u^{q-1}=1$ and $(\alpha_{k+1}^{q-1})^\ell=\alpha_k^{q-1}$, so $(\alpha_k^{q-1})_k\in\Zl(1)$. Another system is $(\epsilon_k\alpha_k)$ with
$\epsilon=(\epsilon_k)\in\Zl(1)$, and $(\epsilon_k\alpha_k)^{q-1}=\epsilon_k^{q-1}\alpha_k^{q-1}$: in additive notation the value changes by $(q-1)\epsilon$; and
$(q-1)\Zl(1)=\ell^v\Zl(1)$ as $(q-1)/\ell^v\in\Zl^\times$. The projection to the $v$-th component identifies $\Zl(1)/\ell^v\Zl(1)$ with $\mu_{\ell^v}$, and the $v$-th
component of $\kappa_u(\Frob_q)$ is $\alpha_v^{q-1}=u^{(q-1)/\ell^v}$ by (a). Finally, $\ell^k\Zl(1)/\ell^v\Zl(1)$ corresponds to $\mu_{\ell^{v-k}}=\mu_{\ell^v}^{\ell^k}$, and
$u^{(q-1)/\ell^v}\in\mu_{\ell^{v-k}}$ iff $u^{(q-1)/\ell^k}=1$, iff $u\in(\F_q^\times)^{\ell^k}$ by (a).
\end{proof}

\begin{paperlink}
Lemma lem:tamfactor defines $c$ as the class of the reduction $\bar u$ of the unit part $u=q_Eq^{-n}$ of the Tate parameter in
$\F_q^\times/(\F_q^\times)^{\ell^v}$, identified with $\Z/\ell^v$ by a choice of generator, and states that $v_\ell(c)\ge k$ iff $\bar u$ is an $\ell^k$-th power, for
$k\le v$, with $v_\ell(0)$ read as $v$. That is \Cref{prop:galois:kummer-finite-field}(a), (b): the valuation of the class does not depend on the generator, and
the class is the value at $\Frob_q$ of the Kummer cocycle of $\bar u$, an element of $\Zl(1)/(q-1)=\F_q^\times\otimes\Zl$. In the proof of the lemma this class is the
unramified part $\kappa_u$ of the Kummer class $\kappa=n\kappa_q+\kappa_u$ of $q_E$; passing from $\F_q$ to $\Q_q$ needs Hensel's lemma and the inflation from $G_{\Q_q}/I_q$,
which are \Cref{ch:local} and \Cref{ch:galcoh}. In Lemma lem:tamagawa-exp the field is $\Q_q^{\mathrm{ur}}$, which contains $\mu_{\ell^t}$ because $q\ne\ell$
(\Cref{ch:local}), so \Cref{prop:galois:kummer-character}(b) applies over it: inertia, which is $G_{\Q_q^{\mathrm{ur}}}$, fixes an $\ell^t$-th root of $q_E$ iff $q_E$ is an
$\ell^t$-th power in $\Q_q^{\mathrm{ur}}$.
\end{paperlink}

\begin{theorem}[Artin-Schreier theory]\label{thm:galois:artin-schreier}
Let $\operatorname{char}F=p>0$ and $\wp(b)=b^p-b$, an additive map $F\to F$ with kernel $\F_p$.
\begin{enumerate}[label=(\alph*)]
\item For $a\in F$ the polynomial $x^p-x-a$ is separable; if $\alpha$ is a root, its roots are $\alpha+i$, $i\in\F_p$. It splits over $F$ if $a\in\wp(F)$, and is
irreducible otherwise; then $F(\alpha)/F$ is cyclic of degree $p$, with group generated by $\alpha\mapsto\alpha+1$.
\item Every cyclic extension of $F$ of degree $p$ is $F(\alpha)$ with $\alpha^p-\alpha\in F$.
\item $\Hom_{\mathrm{cont}}(G_F,\Z/p)\simeq F/\wp(F)$, the class of $a$ corresponding to $\sigma\mapsto\sigma(\alpha)-\alpha$ ($\wp(\alpha)=a$).
\end{enumerate}
For example $x^p-x-1$ is irreducible over $\F_p$ (no element of $\F_p$ is a root, as $\wp(\F_p)=0$), with splitting field $\F_{p^p}$.
\end{theorem}

\begin{proof}[Proof of (a) and (b)]
(a) The derivative is $-1$, so the polynomial is separable; $(\alpha+i)^p-(\alpha+i)-a=\alpha^p-\alpha-a+i^p-i=0$ for $i\in\F_p$. If one root lies in $F$, all do.
Otherwise $F(\alpha)$ contains all roots, so it is Galois over $F$, and $\sigma\mapsto\sigma(\alpha)-\alpha\in\F_p$ is an injective homomorphism (as in
\Cref{prop:galois:kummer-character}(a), $\F_p$ being fixed) into $\Z/p$ with nontrivial image; so the group is $\Z/p$, $[F(\alpha):F]=p$, and $x^p-x-a$ is irreducible.
(b) Let $\Gal(L/F)=\langle\sigma\rangle$ of order $p$. Then $\Tr_{L/F}(-1)=-p=0$, so \Cref{thm:galois:hilbert90}(c) gives $c$ with $-1=c-\sigma(c)$, that is $\sigma(c)=c+1$.
Then $c\notin F$, so $L=F(c)$, and $\sigma(\wp(c))=(c+1)^p-(c+1)=\wp(c)$, so $\wp(c)\in F$.
\end{proof}

\begin{proof}[Sketch of proof of (c)]
As in the proof of \Cref{cor:galois:kummer-absolute}, with (a), (b) and the additive pairing $(\sigma,a)\mapsto\sigma(\alpha)-\alpha$ in place of the
multiplicative one; see \cite[Chapter VI, Section 8]{Lang}.
\end{proof}

The name is shared with the Artin-Schreier theorem on real closed fields (\Cref{prop:galois:complex-conjugation}(c)); the two are unrelated. Artin-Schreier
extensions are the cyclic $p$-extensions that Kummer theory cannot reach in characteristic $p$; over the $p$-adic fields of \Cref{ch:local}, their analogues are
the wildly ramified extensions behind Remark rem:wild of the paper.

\section[The Galois group of a polynomial]{The Galois group of a polynomial; discriminants, cubics, quartics and the four lines}\label{sec:galois:polynomials}

The Galois group of a polynomial is the group of permutations of its roots induced by field automorphisms, and its orbits are the irreducible factors.
This is the dictionary behind the last sentence of Proposition prop:ns: the four roots of the $3$-division polynomial $\psi_3$ of $E$ are the $x$-coordinates of
the four lines of $E[3]$, so the factorisation of $\psi_3$ over $\Q$ records the orbits of the mod $3$ image on the lines, and its discriminant records the
determinant. It is also what the test \texttt{ns3} of the paper's appendix and the consistency check of Section sec:data compute. The discriminant, cubics and
quartics are developed only as far as this application needs.

\begin{proposition}[The Galois group of a polynomial]\label{prop:galois:polynomial-group}
Let $f\in F[x]$ be separable of degree $n\ge1$, $R\subseteq F^{\mathrm{sep}}$ its set of $n$ roots, $L=F(R)$ its splitting field and $G_f=\Gal(L/F)$, the
\emph{Galois group of $f$}.
\begin{enumerate}[label=(\alph*)]
\item $G_f$ acts faithfully on $R$; numbering the roots gives an injective homomorphism $G_f\to S_n$, whose image is well defined up to conjugation in $S_n$;
$\#G_f=[L:F]$.
\item The $G_f$-orbits on $R$ are the root sets of the monic irreducible factors of $f$ in $F[x]$: the orbit of $r$ is the set of roots of $m_{r,F}$; its size is
$[F(r):F]$, and the stabiliser of $r$ is $\Gal(L/F(r))$. In particular $f$ is irreducible iff $G_f$ is transitive on $R$, and then $n\mid\#G_f$.
\item Let $\Omega/F$ be a Galois extension containing $L$ (for instance $\Omega=F^{\mathrm{sep}}$). Then $\Gal(\Omega/F)$ acts on $R$ through its quotient $G_f$, with
the same orbits; and for every intermediate field $F\subseteq M\subseteq\Omega$, the irreducible factors of $f$ in $M[x]$ correspond to the orbits of
$\Gal(\Omega/M)$ on $R$.
\end{enumerate}
\end{proposition}

\begin{proof}
(a) An automorphism of $L/F$ permutes the roots of $f$ and is determined by its values on the generators $R$. Renumbering conjugates the image; $\#G_f=[L:F]$
as $L/F$ is Galois (\Cref{thm:galois:galois-equivalences}).
(b) Since $f$ is separable, $f=c\prod_jg_j$ with distinct monic irreducible $g_j$, whose root sets partition $R$. For $r\in R$ and $\sigma\in G_f$, $\sigma(r)$ is a
root of $m_{r,F}$; conversely, by \Cref{thm:galois:artin}(a) applied to $G_f$ (with $L^{G_f}=F$), $m_{r,F}=\prod_{\beta\in G_fr}(x-\beta)$. So the orbit of $r$ is the
root set of $m_{r,F}$, of size $\deg m_{r,F}=[F(r):F]$; $\sigma$ fixes $r$ iff it fixes $F(r)$. A transitive group has order divisible by the orbit size $n$.
(c) Restriction $\Gal(\Omega/F)\to G_f$ is surjective (\Cref{thm:galois:fundamental}(d) for finite $\Omega$; \Cref{thm:galois:krull}(d) in general), so the
orbits are the same. Over $M$: $f$ is separable in $M[x]$, $\Omega/M$ is Galois, and the first part applied over $M$ (with splitting field $M(R)\subseteq\Omega$)
gives the factors over $M$.
\end{proof}

\begin{definition}[Discriminant]\label{def:galois:discriminant}
Let $f=a\prod_{i=1}^n(x-r_i)$ with $a\ne0$ and roots $r_i$ in $\Fbar$. Put $\delta_f=\prod_{i<j}(r_i-r_j)$ (which depends on the numbering up to sign) and
$\disc(f)=a^{2n-2}\delta_f^2$. For monic $g=\prod_i(x-r_i)$ and $h=\prod_j(x-s_j)$, the resultant is $\operatorname{res}(g,h)=\prod_{i,j}(r_i-s_j)$ (written
$\operatorname{res}$ to avoid confusion with the restriction $\Res$).
\end{definition}

\begin{proposition}[Discriminants]\label{prop:galois:discriminant}
Let $f=a\prod(x-r_i)\in F[x]$ have degree $n$.
\begin{enumerate}[label=(\alph*)]
\item $\disc(f)=(-1)^{n(n-1)/2}a^{n-2}\prod_{i=1}^nf'(r_i)$; it lies in $F$, and it is nonzero iff $f$ is separable.
\item If $f$ is separable and $\sigma\in G_f\subseteq S_n$, then $\sigma(\delta_f)=\sgn(\sigma)\delta_f$.
\item Let $\operatorname{char}F\ne2$ and $f$ separable. Then $G_f\subseteq A_n$ iff $\disc(f)\in F^{\times2}$. The fixed field of $G_f\cap A_n$ is
$F(\delta_f)=F(\sqrt{\disc f})$, and the quadratic character of $\Gal(F^{\mathrm{sep}}/F)$ attached to it by \Cref{prop:galois:quadratic} is the sign of the
permutation of the roots: $\sigma(\sqrt{\disc f})=\sgn(\sigma|_R)\sqrt{\disc f}$.
\item For monic $g,h$: $\disc(gh)=\disc(g)\disc(h)\operatorname{res}(g,h)^2$.
\item $\disc(x^2+bx+c)=b^2-4c$; $\disc(x^3+px+q)=-4p^3-27q^2$; $\disc(x^3+bx^2+cx+d)=b^2c^2-4c^3-4b^3d-27d^2+18bcd$;
$\disc(x^n-a)=(-1)^{(n-1)(n-2)/2}n^na^{n-1}$; and for an odd prime $\ell$, $\disc(\Phi_\ell)=(-1)^{(\ell-1)/2}\ell^{\ell-2}=\ell^*\cdot\ell^{\ell-3}$.
\end{enumerate}
\end{proposition}

\begin{proof}
(a) Differentiating, $f'(r_i)=a\prod_{j\ne i}(r_i-r_j)$, so
\[
\prod_if'(r_i)=a^n\prod_{i\ne j}(r_i-r_j)=a^n(-1)^{n(n-1)/2}\delta_f^2
\]
(each pair $i<j$ appears twice, once with a sign). Hence $a^{2n-2}\delta_f^2=(-1)^{n(n-1)/2}a^{n-2}\prod f'(r_i)$. It vanishes iff two roots coincide. It lies in $F$: if $f$ is separable, $\delta_f^2$ is fixed by
$G_f$ (by (b)), hence in $L^{G_f}=F$; in general $\disc(f)$ is the value at the coefficients of $f$ of a universal polynomial with integer coefficients (a symmetric
polynomial in the roots, expressed through the elementary symmetric functions, which are $\pm$ the coefficients divided by $a$; the powers of $a$ clear the
denominators), and we do not need the inseparable case.
(b) $\sigma$ permutes the factors $r_i-r_j$ ($i<j$) up to sign, and the factor $r_i-r_j$ goes to $r_{\sigma(i)}-r_{\sigma(j)}$, which is $-(r_{\sigma(j)}-r_{\sigma(i)})$
exactly when $\sigma$ inverts the pair; the number of inversions has the parity of $\sgn\sigma$.
(c) As $\operatorname{char}F\ne2$ and $\delta_f\ne0$, $\delta_f\ne-\delta_f$. By (b), $\delta_f$ is fixed by $\sigma$ iff $\sigma$ is even; so $G_f\subseteq A_n$ iff
$\delta_f\in F$ (fundamental theorem), iff $\disc f=(a^{n-1}\delta_f)^2$ is a square in $F$ (if $\disc f=b^2$ with $b\in F$, then $a^{n-1}\delta_f=\pm b$). The fixed field
of $G_f\cap A_n$, of index $1$ or $2$, contains $\delta_f$, and equals $F(\delta_f)$ by degrees. The last formula is (b) with $\sqrt{\disc f}=a^{n-1}\delta_f$.
(d) The roots of $gh$ are those of $g$ and of $h$; the squared differences between roots of $g$ and roots of $h$ contribute $\prod_{i,j}(r_i-s_j)^2$.
(e) The quadratic case is $(r_1-r_2)^2=(r_1+r_2)^2-4r_1r_2$. For $f=x^3+px+q$, by (a) $\disc f=-\prod_i(3r_i^2+p)$. With $e_1=\sum r_i=0$, $e_2=p$, $e_3=r_1r_2r_3=-q$:
$\sum r_i^2=e_1^2-2e_2=-2p$, $\sum_{i<j}r_i^2r_j^2=e_2^2-2e_1e_3=p^2$, $\prod r_i^2=q^2$, so $\prod(3r_i^2+p)=27q^2+9p\cdot p^2+3p^2(-2p)+p^3=27q^2+4p^3$. For the general
cubic $f=x^3+bx^2+cx+d$ over a field of characteristic $0$, the substitution $x\mapsto x-b/3$ does not change the differences of the roots and gives $f(x-b/3)=x^3+px+q$ with
$p=c-b^2/3$ and $q=d-bc/3+2b^3/27$; expanding $-4p^3-27q^2$ gives $b^2c^2-4c^3-4b^3d-27d^2+18bcd$ (a routine expansion, which we also checked by computer algebra).
This proves the formula over every field of characteristic $0$. For the generic monic cubic, $\prod_{i<j}(r_i-r_j)^2$ is a symmetric polynomial in the roots with integer
coefficients, hence a polynomial with integer coefficients in $b,c,d$ (the theorem on symmetric polynomials, which holds over $\Z$; we do not prove it, see
\cite[Chapter IV, Section 6]{Lang}); the computation just made shows that this polynomial is $b^2c^2-4c^3-4b^3d-27d^2+18bcd$ in $\Q[b,c,d]$, hence in $\Z[b,c,d]$, and
specialising gives the formula over every field, of characteristic $3$ included. For $f=x^n-a$ with roots $r_i$: $f'(r_i)=nr_i^{n-1}$ and $\prod r_i=(-1)^{n+1}a$, so
$\prod f'(r_i)=n^n(-1)^{(n+1)(n-1)}a^{n-1}$ and $\disc f=(-1)^{n(n-1)/2+n^2-1}n^na^{n-1}$; the exponent differs from $(n-1)(n-2)/2$ by $n^2+n-2=(n+2)(n-1)$, which is even.
For $\Phi_\ell$ (degree $\ell-1$, roots $\zeta^a$, $1\le a\le\ell-1$, $\zeta=\zeta_\ell$): differentiating $(x-1)\Phi_\ell(x)=x^\ell-1$ at $\zeta^a$ gives
$(\zeta^a-1)\Phi_\ell'(\zeta^a)=\ell\zeta^{-a}$. Now $\prod_a\zeta^{-a}=\zeta^{-\ell(\ell-1)/2}=1$ ($\ell$ odd) and $\prod_a(\zeta^a-1)=\prod_a(1-\zeta^a)=\Phi_\ell(1)=\ell$ (an even
number of sign changes), so $\prod_a\Phi'_\ell(\zeta^a)=\ell^{\ell-1}/\ell=\ell^{\ell-2}$, and $\disc\Phi_\ell=(-1)^{(\ell-1)(\ell-2)/2}\ell^{\ell-2}=(-1)^{(\ell-1)/2}\ell^{\ell-2}$,
since $\ell-2$ is odd.
\end{proof}

Since $\ell^{\ell-3}$ is a square, $\disc(\Phi_\ell)\in\ell^*\Q^{\times2}$, and (c) gives a second proof that $\sqrt{\ell^*}\in\Q(\zeta_\ell)$: the Galois group $(\Z/\ell)^\times$
of $\Phi_\ell$ acts on the roots $\zeta^a$ by multiplication, and the sign of $x\mapsto ax$ on $(\Z/\ell)^\times$ is the Legendre symbol (Zolotarev's lemma). The values
were also checked with PARI for $\ell\le13$.

\begin{proposition}[Cubics]\label{prop:galois:cubic}
Let $\operatorname{char}F\ne2$ and $f\in F[x]$ a separable cubic. If $f$ is irreducible, $G_f=A_3\simeq C_3$ when $\disc f\in F^{\times2}$ and $G_f=S_3$ otherwise. If
$f$ is a linear factor times an irreducible quadratic $g$, $G_f\simeq C_2$, with splitting field $F(\sqrt{\disc g})=F(\sqrt{\disc f})$. Examples over $\Q$:
$x^3-2$ has $\disc=-108$ and $G_f=S_3$ (\Cref{ex:galois:x3minus2}); $x^3-3x+1$ has $\disc=81$ and $G_f=C_3$, its roots are $2\cos(2\pi k/9)$ for $k=1,2,4$, and its
splitting field is the real cubic subfield $\Q(\zeta_9+\zeta_9^{-1})$ of $\Q(\zeta_9)$.
\end{proposition}

\begin{proof}
An irreducible cubic has transitive $G_f$ (\Cref{prop:galois:polynomial-group}(b)), so $G_f\in\{A_3,S_3\}$, distinguished by \Cref{prop:galois:discriminant}(c). In the
reducible case $L$ is the splitting field of $g$, of degree $2$, and $\disc f=\disc g\cdot\operatorname{res}^2$ with $\operatorname{res}\in F^\times$ by
\Cref{prop:galois:discriminant}(d) ($f$ made monic). For $x^3-3x+1$: it has no rational root ($\pm1$ give $-1$ and $3$), so it is irreducible; its discriminant is
$-4(-3)^3-27=81=9^2$. With $z=\zeta_9+\zeta_9^{-1}=2\cos(2\pi/9)$, $z^3=\zeta_9^3+\zeta_9^{-3}+3z=-1+3z$ (as $\zeta_3+\zeta_3^{-1}=-1$), so $z$ is a root; the others are its
images $\zeta_9^2+\zeta_9^{-2}$ and $\zeta_9^4+\zeta_9^{-4}$ under $\sigma_2$ and $\sigma_4$. Hence the splitting field is $\Q(z)$, the fixed field of $\{\sigma_{\pm1}\}$ in
$\Q(\zeta_9)$, of degree $\varphi(9)/2=3$.
\end{proof}

\begin{proposition}[Quartics]\label{prop:galois:quartic}
Let $\operatorname{char}F\ne2$ and $f=x^4+ax^3+bx^2+cx+d\in F[x]$ separable with roots $r_1,\dots,r_4$. Put $\theta_1=r_1r_2+r_3r_4$, $\theta_2=r_1r_3+r_2r_4$,
$\theta_3=r_1r_4+r_2r_3$, one for each of the three partitions of the roots into two pairs, and $R_f(y)=\prod_{i=1}^3(y-\theta_i)$, the \emph{resolvent cubic}.
\begin{enumerate}[label=(\alph*)]
\item $R_f(y)=y^3-by^2+(ac-4d)y-(a^2d-4bd+c^2)$ and $\disc R_f=\disc f$; in particular the $\theta_i$ are distinct.
\item $S_4$ acts on the three partitions through a surjection $S_4\to S_3$ with kernel
\[
V_4=\{1,(12)(34),(13)(24),(14)(23)\};
\]
for any separable $f$,
$\theta_i\in F$ iff $G_f$ stabilises the $i$-th partition, and the degrees of the irreducible factors of $f$ are the orbit sizes of $G_f$.
\item If $f$ is irreducible, $G_f$ is conjugate to one of $S_4$, $A_4$, $D_4$, $C_4$, $V_4$, and: $R_f$ irreducible and $\disc f\notin F^{\times2}$ gives $S_4$;
$R_f$ irreducible and $\disc f\in F^{\times2}$ gives $A_4$; $R_f$ split into linear factors gives $V_4$; $R_f$ with exactly one root in $F$ gives $D_4$ or $C_4$, and
$C_4$ iff $f$ is reducible over $F(\sqrt{\disc f})$ (the criterion of Kappe and Warren \cite{KappeWarren89}).
\end{enumerate}
\end{proposition}

\begin{proof}
(a) With $e_1=-a$, $e_2=b$, $e_3=-c$, $e_4=d$ the elementary symmetric functions of the roots, direct expansion gives $\theta_1+\theta_2+\theta_3=e_2$,
$\sum_{i<j}\theta_i\theta_j=e_1e_3-4e_4$ and $\theta_1\theta_2\theta_3=e_1^2e_4+e_3^2-4e_2e_4$ (each is a short computation; we also checked them by computer algebra),
which is the formula. Next $\theta_1-\theta_2=(r_1-r_4)(r_2-r_3)$, $\theta_1-\theta_3=(r_1-r_3)(r_2-r_4)$, $\theta_2-\theta_3=(r_1-r_2)(r_3-r_4)$, whose product is $\pm\delta_f$;
so $\disc R_f=\delta_f^2=\disc f$.
(b) A permutation of the roots permutes the partitions; the kernel consists of the permutations preserving each partition, which is $V_4$, and the image has
order $24/4=6$. By (a) the $\theta_i$ are distinct, so $\sigma\in G_f$ fixes $\theta_i$ iff it stabilises the $i$-th partition, and $\theta_i\in F$ iff all of $G_f$ does.
The last statement is \Cref{prop:galois:polynomial-group}(b).
(c) $G_f$ is transitive, so $4\mid\#G_f$. The subgroups of $S_4$ of order $24$, $12$, $8$ are $S_4$, $A_4$ (the only subgroup of index $2$) and the Sylow $2$-subgroups
$D_4$ (\Cref{ch:groups}); a subgroup of order $4$ is cyclic, generated by a $4$-cycle (transitive), or a Klein group, either the normal $V_4$ (transitive) or
$\langle(ij),(kl)\rangle$ for disjoint transpositions (not transitive), since two commuting transpositions are disjoint. The image of $G_f$ in $S_3$ is $S_3$, $A_3$,
trivial for $S_4$, $A_4$, $V_4$, and of order $2$ for $D_4$ and $C_4$ (which contain $V_4$, respectively $\{1,(13)(24)\}$ for $C_4=\langle(1234)\rangle$, and stabilise one
partition). By (b), $R_f$ is irreducible iff the image is transitive on the partitions ($S_4$, $A_4$, told apart by the discriminant, $A_4=S_4\cap A_4$), splits iff the
image is trivial ($V_4$), and has exactly one root iff $G_f\in\{D_4,C_4\}$. In the last case $G_f\cap A_4$ is $V_4$ for $D_4$ (transitive) and
$\langle(13)(24)\rangle$ for $C_4$ (not transitive); by \Cref{prop:galois:discriminant}(c) and \Cref{prop:galois:polynomial-group}(c), $G_f\cap A_4$ is the group of $f$
over $F(\sqrt{\disc f})$, so $f$ stays irreducible there iff $G_f=D_4$.
\end{proof}

The next proposition is the Galois-theoretic form of the last sentence of Proposition prop:ns. Its group theory is \Cref{ch:groups}: $\GL_2(\F_3)$ acts on the set
$\PP(V)$ of the four lines of $V=\F_3^2$, the kernel of the action is the group of scalars $\{\pm1\}$, the induced map $\PGL_2(\F_3)\to S_4$ is an isomorphism, the sign of
the permutation induced by $g$ is $\det g\in\{\pm1\}=\F_3^\times$, and the three properties of the paper (surjective determinant, no stable line, an element of
determinant $-1$ and trace $0$) leave exactly the three images $\GL_2(\F_3)$, $N_{ns}(\F_3)$, $N_s(\F_3)$. The table below was also checked by enumerating the $55$
subgroups of $\GL_2(\F_3)$ in $16$ conjugacy classes.

\begin{proposition}[Quartics with group acting on the four lines]\label{prop:galois:quartic-on-lines}
Let $V$ be a two-dimensional $\F_3$-vector space, $\rho\colon G_F\to\GL(V)\simeq\GL_2(\F_3)$ a continuous homomorphism with image $H$, and $f\in F[x]$ a separable
quartic whose root set $R$ is in a $G_F$-equivariant bijection with $\PP(V)$.
\begin{enumerate}[label=(\alph*)]
\item $G_F$ acts on $R$ through $H\to\PGL_2(\F_3)\simeq S_4$, so $G_f$ is isomorphic to the image $\bar H$ of $H$ in $\PGL_2(\F_3)$, and the splitting field of $f$ is the
fixed field of $\rho^{-1}(\{\pm1\})$.
\item The irreducible factors of $f$ over $F$ correspond to the orbits of $H$ on the four lines, with degrees the orbit sizes.
\item The sign of $\sigma\in G_F$ on $R$ is $\det\rho(\sigma)$. Hence, if $\operatorname{char}F\ne2$, $F(\sqrt{\disc f})$ is the field cut out by $\det\circ\rho$, and
$\disc f\in F^{\times2}$ iff $\det(H)=1$.
\item For the subgroups $H$ with $\det(H)=\F_3^\times$, up to conjugacy, the orbit sizes (which by (b) are the degrees of the irreducible factors of $f$)
and the group $G_f$ are:
\begin{center}
\begin{tabular}{llccl}
\toprule
$H$ & label & $\#H$ & orbits on lines & $G_f$\\
\midrule
$\GL_2(\F_3)$ & & $48$ & $4$ & $S_4$\\
$N_{ns}(\F_3)$ & 3Nn & $16$ & $4$ & $D_4$\\
$C_{ns}(\F_3)$ & 3Cn & $8$ & $4$ & $C_4$\\
$N_s(\F_3)$ & 3Ns & $8$ & $2+2$ & $V_4$, not normal in $S_4$\\
$B(\F_3)$ & 3B & $12$ & $1+3$ & $S_3$\\
two of order $6$ & 3B.1.1, 3B.1.2 & $6$ & $1+3$ & $S_3$\\
$C_s(\F_3)$ & 3Cs & $4$ & $1+1+2$ & $C_2$\\
$\langle\operatorname{diag}(1,-1)\rangle$ & 3Cs.1.1 & $2$ & $1+1+2$ & $C_2$\\
\bottomrule
\end{tabular}
\end{center}
So $f$ is irreducible in the first three cases, a product of two irreducible quadratics for $N_s(\F_3)$, a linear factor times an irreducible cubic in the
Borel cases, and two linear factors times an irreducible quadratic in the last two.
\item If $H=N_s(\F_3)$, the two quadratic factors have distinct splitting fields $K_1\ne K_2$, and the third quadratic subfield of $K_1K_2$ is $F(\sqrt{\disc f})$, the
field cut out by $\det\circ\rho$. The resolvent cubic of $f$ (made monic) has exactly one root in $F$ when $H$ is $N_s(\F_3)$ or $N_{ns}(\F_3)$, and none when
$H=\GL_2(\F_3)$.
\item If $H$ fixes no line and contains an element of determinant $-1$ and trace $0$, then $H$ is conjugate to $\GL_2(\F_3)$, $N_{ns}(\F_3)$ or $N_s(\F_3)$, and $f$ is
irreducible in the first two cases and a product of two irreducible quadratics in the third.
\end{enumerate}
\end{proposition}

\begin{proof}
(a) The action of $\sigma$ on $R$ corresponds to the action of $\rho(\sigma)$ on $\PP(V)$, which depends only on the image of $\rho(\sigma)$ in $\PGL_2(\F_3)$, and
$\PGL_2(\F_3)$ acts faithfully. So the kernel of $G_F\to\mathrm{Sym}(R)$ is $\rho^{-1}(\{\pm1\})$, whose fixed field is the splitting field $F(R)$ (an element of $G_F$
fixes $F(R)$ iff it fixes every root), and $G_f\simeq G_F/\rho^{-1}(\{\pm1\})\simeq\bar H$.
(b) \Cref{prop:galois:polynomial-group}(c) with $\Omega=F^{\mathrm{sep}}$.
(c) The sign of the permutation of $\PP(V)$ induced by $g$ is $\det g$ (\Cref{ch:groups}). Apply \Cref{prop:galois:discriminant}(c).
(d) The orders, orbits and images are the group theory of \Cref{ch:groups} (and our enumeration). $G_f$ is the image $\bar H=H/(H\cap\{\pm1\})$; all these $H$ contain
$-1$ except the two groups of order $6$ (which contain an element of order $3$ and one of determinant $-1$, and meet $\{\pm1\}$ trivially) and $\langle\operatorname{diag}
(1,-1)\rangle$, so $\#\bar H=\#H/2$ in the other cases. The transitive images of orders $24$, $8$, $4$ are $S_4$, $D_4$ and $C_4$ (for $C_{ns}(\F_3)$, cyclic of order $8$, the
image is cyclic of order $4$); for $N_s(\F_3)$ the image has order $4$ with two orbits of size two, hence is $\langle(12),(34)\rangle$ for a suitable numbering, the Klein
group that is not normal in $S_4$. The factorisations follow from (b).
(e) Number the lines so that $\bar H=\langle(12),(34)\rangle$ and $G_f=\bar H$ acts on $R$ in the same way. $K_2=F(r_3)$, the splitting field of the factor with roots
$r_3,r_4$, is the fixed field of the stabiliser $\langle(12)\rangle$ of $r_3$; likewise $K_1=F(r_1)$ is the fixed field of $\langle(34)\rangle$. These are two of the three
subgroups of order $2$ of $\bar H\simeq V_4$, so $K_1\ne K_2$; the third, $\langle(12)(34)\rangle=\bar H\cap A_4$, has fixed field $F(\sqrt{\disc f})$ by
\Cref{prop:galois:discriminant}(c). By \Cref{prop:galois:quartic}(b) the roots of $R_f$ in $F$ correspond to the partitions stabilised by $G_f$: $\langle(12),(34)\rangle$
stabilises $\{12|34\}$ and swaps the other two; $D_4$ stabilises exactly one partition; $S_4$ none.
(f) The group-theoretic statement is \Cref{ch:groups}; the determinant is then surjective (it takes the values $1$ and $-1$), and (d) applies.
\end{proof}

\begin{paperlink}
In Proposition prop:ns, $F=\Q$, $V=E[3]$, $\rho=\bar\rho_{E,3}$, and $f=\psi_3$ is the $3$-division polynomial (leading coefficient $3$; irreducibility is meant up to
this constant). The equivariant bijection $R\leftrightarrow\PP(E[3])$ is $\{\pm P\}\mapsto x(P)$, which is injective because $x(P)=x(Q)$ iff $Q=\pm P$
(\Cref{ch:ellcurves}). The inputs ``$\det\bar\rho_{E,3}=\bar\chi_3=\chi_{-3}$'' (\Cref{ch:ellcurves}, Weil pairing, and \Cref{prop:galois:sqrt-lstar}) and ``complex
conjugation has determinant $-1$ and trace $0$'' (\Cref{prop:galois:complex-conjugation}) give the hypotheses of (f) when $E[3]$ is irreducible. Then (f) is the
paper's last sentence: $\psi_3$ is a product of two irreducible quadratics exactly for the image $N_s(\F_3)$, and irreducible otherwise. By (c),
$\disc\psi_3\in-3\cdot\Q^{\times2}$ for every $E/\Q$; for 11a1, $\disc\psi_3=-3^3\cdot11^{10}$, and for 121a1, $-3^3\cdot11^4$. Reducible $E[3]$ always gives a
rational line, hence a linear factor of $\psi_3$ (the patterns $1+3$ and $1+1+2$); the pattern $2+2$ of $N_s(\F_3)$ occurs although $E[3]$ is irreducible, because the
factorisation detects the orbits of the image, not reducibility. The argument uses neither Zywina's classification nor the absence of complex multiplication.
The function \texttt{ns3} of the paper's appendix returns $1$ iff PARI's factorisation of \texttt{elldivpol(E,3)} has exactly two factors, the first of degree $2$;
as PARI lists the factors by increasing degree, the possible patterns are those of (d), and $\det\bar\rho_{E,3}$ is surjective for every $E/\Q$, the test
returns $1$ exactly when the image of $\bar\rho_{E,3}$ is $N_s(\F_3)$, with no hypothesis on $E[3]$.
\end{paperlink}

\begin{example}[An image $N_s(\F_3)$, and a surjective image]\label{ex:galois:ns-example}
Let $E$ be 6912d1, $y^2=x^3-6x-8$ ($j=-1728$, which is not a CM $j$-invariant; $\Delta=-2^9\cdot3^3$). PARI gives
\[
\psi_3=3x^4-36x^2-96x-36=3(x^2-4x-2)(x^2+4x+6),
\]
two irreducible quadratics with discriminants $24$ and $-8$, so $K_1=\Q(\sqrt6)$, $K_2=\Q(\sqrt{-2})$, and the third quadratic subfield of $K_1K_2$ is
$\Q(\sqrt{-12})=\Q(\sqrt{-3})$, as \Cref{prop:galois:quartic-on-lines}(e) predicts; the resolvent cubic of $\psi_3/3$ is $(y-4)(y^2+16y+112)$, with exactly one rational
root. Since $\det\bar\rho_{E,3}$ is surjective and only $N_s(\F_3)$ has the pattern $2+2$ in the table of (d), the image is $N_s(\F_3)$ (Sutherland's 3Ns, as in Cremona's
tables). The group $N_s(\F_3)$ acts absolutely irreducibly on $\Fbar_3^2$: a line
stable under it is an eigenline of $\operatorname{diag}(1,-1)$, hence a coordinate axis, and the antidiagonal elements swap the two axes. The image of
$G_{\Q(\sqrt{-3})}$ is the determinant-one subgroup $\{\pm1,\pm J\}$ (\Cref{prop:galois:restriction}(c)), which is abelian, so $\bar\rho|_{G_{\Q(\sqrt{-3})}}$ is not absolutely
irreducible, and $\bar\rho\simeq\bar\rho\otimes\chi_{-3}$ by \Cref{prop:galois:index-two}. For the primes $p\ne3$ of good reduction (here $p\ge5$) one has
$\tr\bar\rho_{E,3}(\Frob_p)\equiv a_p\pmod3$ (\Cref{ch:ellcurves}, \Cref{ch:galrep}). So $a_p\equiv\chi_{-3}(\Frob_p)a_p\pmod3$, and $a_p\equiv0\pmod3$ whenever
$\chi_{-3}(\Frob_p)=-1$, that is for $p\equiv2\pmod 3$ (\Cref{prop:galois:quadratic-characters}(c)). This part of the example looks ahead: it uses
\Cref{prop:galois:restriction}, \Cref{prop:galois:index-two} and \Cref{prop:galois:quadratic-characters} of Section 8 and the Frobenius elements of Section 9
(\Cref{thm:galois:frobenius}); nothing in Sections 6 to 8 depends on it. Indeed $a_5=3$, $a_{11}=3$, $a_{17}=-6$, $a_{23}=6$, $a_{29}=6$,
$a_{41}=6$, $a_{47}=12$, $a_{53}=-3$, $a_{59}=12$, $a_{71}=-6$; note $9\nmid a_5$. By contrast 11a1, $y^2+y=x^3-x^2-10x-20$, has
\[
\psi_3=3x^4-4x^3-60x^2-237x-21,
\]
irreducible, with irreducible resolvent cubic and $\disc\psi_3=-3^3\cdot11^{10}\notin\Q^{\times2}$: $G_{\psi_3}=S_4$ and the image of $\bar\rho_{E,3}$ is $\GL_2(\F_3)$. The paper's own
example 726a1 is \Cref{exer:galois:726}.
The computations are reproduced in PARI/GP by:
\begin{lstlisting}[language=GP]
E = ellinit([0,0,0,-6,-8]); f = elldivpol(E,3);
factor(f)                       \\ [x^2-4*x-2,1; x^2+4*x+6,1]
[poldisc(x^2-4*x-2), poldisc(x^2+4*x+6)]          \\ [24, -8]
[[p, ellap(E,p)] | p <- primes([5,71]), p%3 == 2]
F = elldivpol(ellinit([0,-1,1,-10,-20]),3);  \\ 11a1
[factor(F), factor(poldisc(F))]
\end{lstlisting}
\end{example}

\section[Infinite Galois theory, \texorpdfstring{$G_\Q$}{GQ}, the cyclotomic character]{Infinite Galois theory, \texorpdfstring{$G_\Q$}{GQ}, the cyclotomic character, complex conjugation}\label{sec:galois:infinite}

The Galois groups of the paper are infinite: $G_\Q$ in Section 1.1 and Section 4, the decomposition and inertia groups $G_{\Q_q}\supseteq I_q$ in Lemmas
lem:tamagawa-exp and lem:tamfactor. This section extends the fundamental theorem to infinite Galois extensions, where subgroups must be \emph{closed} in the Krull
topology, defines the absolute Galois group and the cyclotomic character $\chi_\ell$, proves $\chi_\ell(G_\Q)=\Zl^\times$ (the fact that forces $\ell=3$ in Lemma
lem:h0), and studies complex conjugation, the element of $G_\Q$ that makes ``irreducible'' and ``absolutely irreducible'' coincide for $\bar\rho_{E,\ell}$. Inverse
limits, profinite groups and their basic properties (open subgroups are closed of finite index, open normal subgroups form a basis at $1$, continuous images of
compact groups are closed) are \Cref{ch:groups}. One compactness statement is used repeatedly, here and in \Cref{ch:ant}, so we isolate it first.

\begin{lemma}[Inverse limits of nonempty finite sets]\label{lem:galois:nonempty-limit}
Let $(I,\le)$ be a directed set (a partially ordered set in which any two elements have a common upper bound), and let $(X_i)_{i\in I}$ be nonempty finite sets with maps
$f_{ij}\colon X_j\to X_i$ for $i\le j$, such that $f_{ii}$ is the identity and $f_{ij}\circ f_{jk}=f_{ik}$ for $i\le j\le k$. Then the inverse limit
\[
\varprojlim_{i\in I}X_i=\Bigl\{(x_i)_{i\in I}\in\prod_{i\in I}X_i:\ f_{ij}(x_j)=x_i\text{ for all }i\le j\Bigr\}
\]
is nonempty.
\end{lemma}

\begin{proof}
Give each $X_i$ the discrete topology; it is compact, being finite, so $P=\prod_iX_i$ is compact by Tychonoff's theorem (\Cref{ch:groups}). For $i\le j$ the set
$C_{ij}=\{x\in P:f_{ij}(x_j)=x_i\}$ is closed: it is the preimage of a subset of the finite discrete space $X_i\times X_j$ under the continuous projection $P\to X_i\times X_j$.
The inverse limit is $\bigcap_{i\le j}C_{ij}$. By compactness it suffices to show that finitely many of the $C_{ij}$, say $C_{i_1j_1},\dots,C_{i_rj_r}$, have a common point.
As $I$ is directed, there is $m\in I$ with $j_s\le m$ (hence $i_s\le m$) for $s=1,\dots,r$. Choose $y\in X_m$ (it is nonempty), put $x_i=f_{im}(y)$ for every $i$ among the
$i_s$, $j_s$, and choose $x_i\in X_i$ arbitrarily for the other $i$. Then $f_{i_sj_s}(x_{j_s})=f_{i_sj_s}f_{j_sm}(y)=f_{i_sm}(y)=x_{i_s}$ for each $s$, so $x$ lies in all $r$
sets. Hence $\bigcap_{i\le j}C_{ij}\ne\emptyset$.
\end{proof}

Finiteness matters: an inverse limit of nonempty \emph{infinite} sets can be empty (take $X_n=\{n,n+1,n+2,\dots\}$, $n\in\N$, with the inclusions as maps).

\begin{proposition}[Infinite Galois extensions]\label{prop:galois:infinite}
Let $L/F$ be a Galois extension of arbitrary degree and $\mathcal{K}$ the set of finite Galois subextensions $K/F$ of $L$.
\begin{enumerate}[label=(\alph*)]
\item $L=\bigcup_{K\in\mathcal{K}}K$; $\mathcal{K}$ is directed by inclusion (any two members lie in a third); and restriction is a group isomorphism
\[
\Gal(L/F)\xrightarrow{\ \sim\ }\varprojlim_{K\in\mathcal{K}}\Gal(K/F),\qquad\sigma\mapsto(\sigma|_K)_K .
\]
\item The \emph{Krull topology} on $\Gal(L/F)$ is the one in which the subgroups $\Gal(L/K)$, $K\in\mathcal{K}$, form a basis of open neighbourhoods of $1$ (and their
cosets a basis of the topology). It is the inverse-limit topology transported by (a), so $\Gal(L/F)$ is a profinite group: compact, Hausdorff and totally
disconnected. The subgroups $\Gal(L/M)$ with $M\subseteq L$ finite over $F$ are open.
\item The restriction maps $\Gal(L/F)\to\Gal(K/F)$, $K\in\mathcal{K}$, are surjective. If $L/F$ is finite, the Krull topology is discrete.
\end{enumerate}
\end{proposition}

\begin{proof}
(a) For $\alpha\in L$, $m_{\alpha,F}$ is separable and splits in $L$ (normality), so its splitting field is a member of $\mathcal{K}$ containing $\alpha$. If
$K_1,K_2\in\mathcal{K}$, then $K_1K_2\in\mathcal{K}$ (\Cref{prop:galois:translation}(b)). Each $\sigma\in\Gal(L/F)$ maps $K\in\mathcal{K}$ onto itself
(\Cref{prop:galois:normal}(iii)), so $\sigma|_K\in\Gal(K/F)$, and the restrictions are compatible. The map is injective because $L$ is the union of the $K$.
Surjective: a compatible family $(\sigma_K)$ defines a map $\sigma\colon L\to L$, $\sigma(x)=\sigma_K(x)$ for any $K\ni x$ (well defined by compatibility and
directedness), which is a field automorphism over $F$ (its inverse is given by the family $(\sigma_K^{-1})$).
(b) In $\varprojlim\Gal(K/F)$ with the topology induced by the product of discrete groups, a basis of neighbourhoods of $1$ is given by the kernels of the
projections to finitely many factors $\Gal(K_1/F),\dots,\Gal(K_r/F)$; under (a) such a kernel is $\Gal(L/K_1\cdots K_r)$, and $K_1\cdots K_r\in\mathcal{K}$. So the
two topologies agree, and the inverse limit of finite groups is profinite (\Cref{ch:groups}). If $M/F$ is finite, its normal closure $M'$ in $L$ is in $\mathcal{K}$
(\Cref{prop:galois:normal}), and $\Gal(L/M)\supseteq\Gal(L/M')$ is a union of cosets of an open subgroup.
(c) Let $K_0\in\mathcal{K}$ and $\tau\in\Gal(K_0/F)$. The set $\mathcal{K}_0$ of $K\in\mathcal{K}$ containing $K_0$ is directed by inclusion (by (a)). For $K\in\mathcal{K}_0$ the
set $X_K$ of $\sigma\in\Gal(K/F)$ with $\sigma|_{K_0}=\tau$ is finite and nonempty (\Cref{thm:galois:fundamental}(d)), and for $K'\subseteq K''$ in $\mathcal{K}_0$ restriction
maps $X_{K''}$ to $X_{K'}$, compatibly with composition. By \Cref{lem:galois:nonempty-limit} there is a compatible family $(\sigma_K)_{K\in\mathcal{K}_0}$. Every $K\in\mathcal{K}$
lies in $KK_0\in\mathcal{K}_0$, so the members of $\mathcal{K}_0$ exhaust $L$, and the family defines as in (a) an element $\sigma\in\Gal(L/F)$, with $\sigma|_{K_0}=\tau$. (Alternatively,
extend $\tau$ to an $F$-embedding $L\to\Fbar$ by \Cref{thm:galois:algebraic-closure}(b); its image is $L$ by normality.) For $L/F$ finite, $\{1\}=\Gal(L/L)$ is open.
\end{proof}

\begin{theorem}[Krull's fundamental theorem]\label{thm:galois:krull}
Let $L/F$ be Galois with group $G=\Gal(L/F)$ in the Krull topology.
\begin{enumerate}[label=(\alph*)]
\item For every intermediate field $M$, $L/M$ is Galois, $\Gal(L/M)$ is a closed subgroup of $G$, and $L^{\Gal(L/M)}=M$.
\item For every subgroup $H\le G$, $\Gal(L/L^H)=\overline H$, the closure of $H$.
\item Hence $H\mapsto L^H$ and $M\mapsto\Gal(L/M)$ are mutually inverse inclusion-reversing bijections between the \emph{closed} subgroups of $G$ and the
intermediate fields.
\item Open subgroups correspond to the intermediate fields $M$ finite over $F$, and then $[G:\Gal(L/M)]=[M:F]$.
\item $M/F$ is Galois iff $\Gal(L/M)$ is normal in $G$; then restriction induces an isomorphism of topological groups $G/\Gal(L/M)\simeq\Gal(M/F)$ (quotient
topology on the left, Krull topology on the right).
\item $\Gal(L/\sigma M)=\sigma\Gal(L/M)\sigma^{-1}$ for $\sigma\in G$.
\end{enumerate}
\end{theorem}

\begin{proof}
(a) $L/M$ is separable, and normal by \Cref{prop:galois:normal}. For $x\in L$, the stabiliser of $x$ in $G$ contains the open subgroup $\Gal(L/K)$ for any $K\in
\mathcal{K}$ containing $x$, so it is open, hence closed; $\Gal(L/M)$ is the intersection of the stabilisers of the $x\in M$, hence closed. If $x\in L\setminus M$, then
$m_{x,M}$ has degree $>1$, is separable and splits in $L$, so it has a root $y\ne x$ in $L$. Let $K\subseteq L$ be the splitting field of $m_{x,M}$ over $M$, finite Galois over $M$; by
\Cref{prop:galois:polynomial-group}(b) some $\tau\in\Gal(K/M)$ has $\tau(x)=y$, and by \Cref{prop:galois:infinite}(c), applied to $L/M$, $\tau$ extends to an element of
$\Gal(L/M)$, which moves $x$. So $L^{\Gal(L/M)}=M$.
(b) $H\subseteq\Gal(L/L^H)$, which is closed by (a), so $\overline H\subseteq\Gal(L/L^H)$. Conversely, let $\sigma\in\Gal(L/L^H)$ and let $\sigma\Gal(L/K)$, $K\in\mathcal{K}$,
be a basic neighbourhood of $\sigma$. Let $H_K\le\Gal(K/F)$ be the image of $H$ under restriction. Its fixed field in $K$ is $K\cap L^H$, which $\sigma$ fixes; so
$\sigma|_K\in\Gal(K/K^{H_K})=H_K$ by the finite fundamental theorem. Thus some $h\in H$ has $h|_K=\sigma|_K$, that is $h\in\sigma\Gal(L/K)$. Every neighbourhood of $\sigma$
meets $H$, so $\sigma\in\overline H$.
(c) follows from (a) and (b).
(d) If $[M:F]<\infty$, $\Gal(L/M)$ is open (\Cref{prop:galois:infinite}(b)). If $\Gal(L/M)$ is open, it contains some $\Gal(L/K)$, $K\in\mathcal{K}$, so $M\subseteq K$
is finite over $F$. Then $\Gal(L/M)$ is the preimage of $\Gal(K/M)$ under the surjection $G\to\Gal(K/F)$, and its index is $[\Gal(K/F):\Gal(K/M)]=[M:F]$.
(e) and (f). (f) is immediate. As in the finite case, every $F$-embedding $M\to\Fbar$ extends to one of $L$ (\Cref{thm:galois:algebraic-closure}(b)), which lies in $G$,
so $M/F$ is normal iff $\sigma M=M$ for all $\sigma\in G$, iff $\Gal(L/M)$ is normal by (f) and (c). Then restriction $G\to\Gal(M/F)$ is a homomorphism, surjective
(extension of automorphisms), with kernel $\Gal(L/M)$. It is continuous: the preimage of a basic open subgroup $\Gal(M/K)$ ($K\subseteq M$ finite Galois over $F$) is
$\Gal(L/K)$, open. The induced bijection $G/\Gal(L/M)\to\Gal(M/F)$ is continuous from a compact space to a Hausdorff space, hence a homeomorphism.
\end{proof}

\begin{definition}[The absolute Galois group]\label{def:galois:absolute}
The \emph{absolute Galois group} of $F$ is $G_F=\Gal(F^{\mathrm{sep}}/F)$ with the Krull topology (\Cref{prop:galois:separable-closure}). For $K\subseteq F^{\mathrm{sep}}$,
$G_K=\Gal(F^{\mathrm{sep}}/K)\subseteq G_F$. In particular $G_\Q=\Gal(\Qbar/\Q)$. A homomorphism from $G_F$ \emph{cuts out} the fixed field of its kernel when
the kernel is closed.
\end{definition}

\begin{proposition}[The absolute Galois group]\label{prop:galois:absolute}
\begin{enumerate}[label=(\alph*)]
\item For $K/F$ finite inside $F^{\mathrm{sep}}$, $G_K$ is an open subgroup of $G_F$ of index $[K:F]$; every open subgroup is of this form; $\sigma G_K\sigma^{-1}=
G_{\sigma K}$; and $G_K$ is normal iff $K/F$ is Galois, in which case $G_F/G_K\simeq\Gal(K/F)$.
\item $G_{\F_q}=\Gal(\Fbar_q/\F_q)\simeq\widehat\Z=\varprojlim\Z/n$, by an isomorphism sending $\Frob_q$ to $1$; $\Frob_q$ generates a dense subgroup, and the open
subgroups are the $n\widehat\Z=G_{\F_{q^n}}$, $n\ge1$.
\item For a number field $K\subseteq\Qbar$, $G_K$ is open of index $[K:\Q]$ in $G_\Q$. So $G_{\Q(\sqrt d)}$ has index $2$ for $d$ not a square, and $G_{\Q(\zeta_n)}$ is open
and normal of index $\varphi(n)$, with $G_\Q/G_{\Q(\zeta_n)}\simeq(\Z/n)^\times$.
\item $G_\Q$ is uncountable and non-abelian. Its maximal abelian quotient is $\Gal(\Q^{\mathrm{ab}}/\Q)\simeq\widehat\Z^\times$, where
$\Q^{\mathrm{ab}}=\Q(\mu_\infty)$ is the union of the cyclotomic fields (the Kronecker-Weber theorem).
\end{enumerate}
\end{proposition}

\begin{proof}[Proof of (a) to (c), and of the first sentence of (d)]
(a) is \Cref{thm:galois:krull}(d), (e), (f) with $L=F^{\mathrm{sep}}$. (b) Every finite extension of $\F_q$ in $\Fbar_q$ is an $\F_{q^n}$, Galois with group
$\Z/n$ generated by $\Frob_q$ (\Cref{thm:galois:finite-fields}); $\F_{q^m}\subseteq\F_{q^n}$ iff $m\mid n$, with restriction $\Z/n\to\Z/m$ the reduction. By
\Cref{prop:galois:infinite}(a), $G_{\F_q}\simeq\varprojlim_n\Z/n$ (divisibility order), which is $\widehat\Z$ of \Cref{ch:groups} (the factorials are cofinal), with
$\Frob_q\mapsto(1\bmod n)_n=1$. The image of $\langle\Frob_q\rangle$ in every finite quotient is everything, so it is dense. The open subgroups are the $G_{\F_{q^n}}$ by
(a), and $G_{\F_{q^n}}$ is the kernel of $\widehat\Z\to\Z/n$, which is $n\widehat\Z$. (c) is (a) with $F=\Q$, and \Cref{thm:galois:cyclotomic}. (d) By (c) and
\Cref{prop:galois:quadratic}(c), $G_\Q$ has infinitely many distinct subgroups of index $2$; a profinite group with infinitely many open subgroups is infinite, and an
infinite profinite group is uncountable (it is a compact Hausdorff space without isolated points, the open subgroups having finite index). It is non-abelian since
$\Gal(\Q(2^{1/3},\zeta_3)/\Q)\simeq S_3$ is a quotient (\Cref{ex:galois:x3minus2}). The rest of (d) is the Kronecker-Weber theorem, which we do not prove; see
\Cref{ch:ant} and \cite[Chapter 14]{Washington}.
\end{proof}

\begin{example}[Two warnings]\label{ex:galois:warnings}
(1) The subgroup $\Frob_p^{\Z}$ of $G_{\F_p}\simeq\widehat\Z$ is dense and not closed: it is countable, while $\widehat\Z$ is uncountable. Its fixed field is $\F_p$, the same as
that of $G_{\F_p}$; this is why \Cref{thm:galois:krull}(c) requires closed subgroups.
(2) A subgroup of finite index need not be open. Let $L=\Q(\sqrt{-1},\sqrt2,\sqrt3,\sqrt5,\sqrt7,\dots)$, the compositum of all quadratic fields. The fields $\Q(\Delta^{1/2})$ with
$\Delta/\Q^{\times2}$ finite are finite Galois and cofinal among the finite subextensions of $L$ (each of these lies in a compositum of finitely many quadratic fields), and
$\Gal(\Q(\Delta^{1/2})/\Q)\simeq\Hom(\Delta/\Q^{\times2},\{\pm1\})$ compatibly with restriction (\Cref{thm:galois:kummer}(a) with $n=2$). Passing to the limit
(\Cref{prop:galois:infinite}(a), where the limit may be taken over this cofinal family) gives $\Gal(L/\Q)\simeq\varprojlim_\Delta\Hom(\Delta/\Q^{\times2},\{\pm1\})=
\Hom(\Q^\times/\Q^{\times2},\{\pm1\})\simeq\prod_{\N}\Z/2$, as $\Q^\times/\Q^{\times2}$ is the union of the $\Delta/\Q^{\times2}$ and is an $\F_2$-vector space with basis
$-1$ and the primes. The open subgroups of index $2$ are the kernels of the countably many continuous characters, which are the
evaluations at the elements of $\Q^\times/\Q^{\times2}$; but $\prod_\N\Z/2$ has uncountably many subgroups of index $2$ (kernels of linear forms on an
$\F_2$-vector space of uncountable dimension, whose existence uses the axiom of choice), so most are not open. We do not use this.
\end{example}

\begin{corollary}[Absolute Kummer theory]\label{cor:galois:kummer-absolute}
Let $\operatorname{char}F\nmid n$ and $\mu_n\subset F$. Then $a\mapsto\chi_a$, $\chi_a(\sigma)=\sigma(\alpha)/\alpha$ ($\alpha^n=a$), is an isomorphism
\[
F^\times/F^{\times n}\xrightarrow{\ \sim\ }\Hom_{\mathrm{cont}}(G_F,\mu_n),
\]
$\mu_n$ being discrete. After a primitive $\zeta_n$ is chosen (so $\mu_n\simeq\Z/n$), the surjective continuous homomorphisms $G_F\to\Z/n$ correspond to the elements of
exact order $n$ of $F^\times/F^{\times n}$. For $n=2$ and $\operatorname{char}F\ne2$: $F^\times/F^{\times2}\simeq\Hom_{\mathrm{cont}}(G_F,\{\pm1\})$, the class of $d$
corresponding to the character of $F(\sqrt d)$.
\end{corollary}

\begin{proof}
$\chi_a$ is continuous: its kernel $G_{F(\alpha)}$ is open. It is a homomorphism in $a$ and injective modulo $F^{\times n}$ by \Cref{prop:galois:kummer-character}(b)
with $L=F^{\mathrm{sep}}$. Surjective: a continuous $\chi\colon G_F\to\mu_n$ has open kernel $G_K$, $K/F$ finite Galois with group a subgroup of $\mu_n$, abelian of
exponent dividing $n$. By \Cref{thm:galois:kummer}(b), $K=F_\Delta$ with $\Delta/F^{\times n}$ finite, and by (a) of that theorem the character of $\Gal(K/F)$ induced by
$\chi$ is $\sigma\mapsto\sigma(\alpha)/\alpha$ for some $a=\alpha^n\in\Delta$; so $\chi=\chi_a$. The surjective homomorphisms are those of order $n$ in the group
$\Hom_{\mathrm{cont}}(G_F,\mu_n)$, which correspond to the elements of order $n$. For $n=2$, $\mu_2=\{\pm1\}\subset F$ always, and $\chi_d(\sigma)=\sigma(\sqrt d)/\sqrt d$.
\end{proof}

\begin{corollary}[Translation, infinite case]\label{cor:galois:translation-infinite}
Let $L/F$ be Galois (of any degree) and $M/F$ any extension, inside a common field. Then $LM/M$ is Galois and restriction $\Gal(LM/M)\to\Gal(L/L\cap M)$ is an
isomorphism of topological groups. In particular, for $K/F$ finite separable and $L\subseteq F^{\mathrm{sep}}$ Galois over $F$, $\Gal(LK/K)\simeq\Gal(L/L\cap K)$.
\end{corollary}

\begin{proof}
$LM$ is generated over $M$ by roots of separable polynomials over $F$ that split in $L$, so $LM/M$ is normal and separable. Restriction is a continuous
homomorphism $\Gal(LM/M)\to\Gal(L/F)$ (the preimage of $\Gal(L/K)$ is $\Gal(LM/MK)$), injective as before. Its image $H$ is compact, hence closed, and as in
\Cref{prop:galois:translation} $L^H=L\cap M$; by \Cref{thm:galois:krull}(c), $H=\Gal(L/L\cap M)$. A continuous bijection from a compact space to a Hausdorff space is a
homeomorphism.
\end{proof}

\begin{definition}[The cyclotomic character]\label{def:galois:cyclotomic-character}
Let $\ell$ be a prime and $F$ a field of characteristic $\ne\ell$. For $k\ge1$ the group $\mu_{\ell^k}=\mu_{\ell^k}(F^{\mathrm{sep}})$ is cyclic of order $\ell^k$
($x^{\ell^k}-1$ is separable), and $G_F$ acts on it by automorphisms, $\sigma(\zeta)=\zeta^{a_k(\sigma)}$ with $a_k(\sigma)\in(\Z/\ell^k)^\times$. The
\emph{$\ell$-adic cyclotomic character} is
\[
\chi_\ell\colon G_F\to\Zl^\times=\varprojlim_k(\Z/\ell^k)^\times,\qquad\chi_\ell(\sigma)=(a_k(\sigma))_k,
\]
so that $\sigma(\zeta)=\zeta^{\chi_\ell(\sigma)}$ for every $\zeta\in\mu_{\ell^\infty}=\bigcup_k\mu_{\ell^k}$; $\bar\chi_\ell=\chi_\ell\bmod\ell\colon G_F\to\F_\ell^\times$. The
\emph{Tate module of $\mu$} is $\Zl(1)=\varprojlim_k\mu_{\ell^k}$ (transition maps $\zeta\mapsto\zeta^\ell$), and for a $\Zl$-module $M$ with $G_F$-action,
$M(1)=M\otimes_{\Zl}\Zl(1)$.
\end{definition}

\begin{proposition}[The cyclotomic character]\label{prop:galois:cyclotomic-character}
\begin{enumerate}[label=(\alph*)]
\item $\chi_\ell$ is a well defined continuous homomorphism, and it is the unique map $G_F\to\Zl^\times$ with $\sigma(\zeta)=\zeta^{\chi_\ell(\sigma)}$ on $\mu_{\ell^\infty}$.
$\chi_\ell\bmod\ell^k$ is trivial iff $\mu_{\ell^k}\subset F$, and $\ker\chi_\ell=G_{F(\mu_{\ell^\infty})}$.
\item For $F=\Q$: $\chi_\ell$ is surjective, and $\Gal(\Q(\mu_{\ell^\infty})/\Q)\simeq\Zl^\times$ as topological groups; $\chi_\ell(\sigma)\bmod\ell^k$ is the $a$ with
$\sigma|_{\Q(\zeta_{\ell^k})}=\sigma_a$.
\item For $F=\Q$: $\bar\chi_\ell$ is surjective onto $\F_\ell^\times$ with kernel $G_{\Q(\zeta_\ell)}$, and for odd $\ell$, $\chi_{\ell^*}=\bigl(\frac{\bar\chi_\ell}{\ell}\bigr)=
\bar\chi_\ell^{(\ell-1)/2}$ (\Cref{prop:galois:sqrt-lstar}).
\item For $F=\F_q$ with $\ell\nmid q$: $\chi_\ell(\Frob_q)=q$.
\item $\Zl(1)$ is a free $\Zl$-module of rank one, on which $\sigma\in G_F$ acts by multiplication by $\chi_\ell(\sigma)$; $\Zl(1)/\ell^k\Zl(1)\simeq\mu_{\ell^k}$, by projection to
the $k$-th component.
\end{enumerate}
\end{proposition}

\begin{proof}
(a) If $\zeta\in\mu_{\ell^{k+1}}$ then $\zeta^\ell\in\mu_{\ell^k}$ and $\sigma(\zeta^\ell)=\zeta^{\ell a_{k+1}(\sigma)}$, so $a_{k+1}(\sigma)\equiv a_k(\sigma)\pmod{\ell^k}$: the
$a_k(\sigma)$ form an element of $\Zl^\times$. Each $a_k$ is multiplicative in $\sigma$ and constant on cosets of the open subgroup $G_{F(\mu_{\ell^k})}$; as the
subgroups $1+\ell^k\Zl$ form a basis of neighbourhoods of $1$ in $\Zl^\times$, $\chi_\ell$ is continuous. Uniqueness is clear, and the statements on kernels are the
definitions.
(b) By \Cref{thm:galois:cyclotomic}, $G_\Q\to\Gal(\Q(\zeta_{\ell^k})/\Q)\simeq(\Z/\ell^k)^\times$ is surjective for each $k$, and this composite is $\chi_\ell\bmod\ell^k$.
So $\chi_\ell(G_\Q)$ is dense in $\Zl^\times$; being compact, it is closed, hence equal to $\Zl^\times$. By \Cref{thm:galois:krull}(e), $\Gal(\Q(\mu_{\ell^\infty})/\Q)\simeq
G_\Q/\ker\chi_\ell\simeq\Zl^\times$, topologically (continuous bijection of compact Hausdorff groups).
(c) The case $k=1$ of (b), and \Cref{prop:galois:sqrt-lstar}(b).
(d) $\Frob_q(\zeta)=\zeta^q$.
(e) Choose $\zeta_{\ell^k}$ inductively: $\zeta_\ell$ a generator of $\mu_\ell$, and $\zeta_{\ell^{k+1}}$ an $\ell$-th root of $\zeta_{\ell^k}$, which has order $\ell^{k+1}$. Then
$a\mapsto(\zeta_{\ell^k}^{a})_k$ is an isomorphism $\Zl\to\Zl(1)$ (its components $\Z/\ell^k\to\mu_{\ell^k}$ are isomorphisms), and $\sigma(\zeta_{\ell^k}^{a})=
\zeta_{\ell^k}^{a\chi_\ell(\sigma)}$. The kernel of the projection to $\mu_{\ell^k}$ corresponds to $\ell^k\Zl$.
\end{proof}

In the additive notation of \Cref{conv:galois:characters}(c), $\sigma\cdot z=\chi_\ell(\sigma)z$ for $z\in\Zl(1)$; so $\Frob_q$ acts on $\Zl(1)$ by $q$, which is the
statement ``$\Frob$ acts on $\Zl(1)$ by $q$'' in the proof of Lemma lem:tamfactor (with $F=\Q_q$ and $\Frob$ acting through the residue field, \Cref{ch:local}).

\begin{proposition}[Complex conjugation]\label{prop:galois:complex-conjugation}
Let $c\in G_\Q$ be the restriction to $\Qbar$ of complex conjugation (\Cref{conv:galois:fields}).
\begin{enumerate}[label=(\alph*)]
\item $c$ is an element of order $2$. Every embedding $\Qbar\to\C$ is $\iota\circ\sigma$ with $\iota$ the inclusion and $\sigma\in G_\Q$, and complex conjugation transported
by it is $\sigma^{-1}c\sigma$. A \emph{complex conjugation} is an element of the conjugacy class of $c$.
\item $\chi_\ell(c)=-1$ for every $\ell$; $\chi_d(c)=-1$ iff $d<0$. So $\chi_{\ell^*}(c)=(-1)^{(\ell-1)/2}$ and $\chi_{-3}(c)=\chi_{-4}(c)=-1$.
\item (Artin and Schreier.) Every element of $G_\Q$ of finite order other than $1$ is a complex conjugation; in particular it has order $2$.
\item Let $k$ be a field with $\operatorname{char}k\ne2$ and $\rho\colon G_\Q\to\GL_2(k)$ a homomorphism with $\det\rho(c)=-1$ (for instance $\det\rho=\bar\chi_\ell$ with
$\ell$ odd, or $\det\rho=\chi_\ell$). Then $\rho(c)$ is conjugate to $\operatorname{diag}(1,-1)$, so it has determinant $-1$ and trace $0$, and $\rho$ is irreducible over
$k$ iff it is absolutely irreducible.
\end{enumerate}
\end{proposition}

\begin{proof}
(a) The conjugate of an algebraic number is algebraic (it is a root of the same polynomial with rational coefficients), so complex conjugation maps $\Qbar$ to
itself; it has order $2$ because $c(i)=-i$. An embedding $\tau\colon\Qbar\to\C$ has image algebraic over $\Q$, so contained in $\Qbar$, and equal to it
($\Qbar/\tau(\Qbar)$ is algebraic and $\tau(\Qbar)$ is algebraically closed); so $\tau=\iota\circ\sigma$ with $\sigma\in G_\Q$, and $x\mapsto\tau^{-1}(\overline{\tau(x)})$
is $\sigma^{-1}c\sigma$.
(b) For $\zeta\in\mu_{\ell^\infty}\subset\C$, $c(\zeta)=\bar\zeta=\zeta^{-1}$. If $d>0$, $\sqrt d$ is real and fixed; if $d<0$, $\sqrt d$ is purely imaginary and
$c(\sqrt d)=-\sqrt d$. Apply to $\ell^*$, $-3$, $-4$.
(d) $\rho(c)^2=1$, so the minimal polynomial of $\rho(c)$ divides $X^2-1=(X-1)(X+1)$, which has distinct roots as $\operatorname{char}k\ne2$; hence $\rho(c)$ is
diagonalisable with eigenvalues in $\{\pm1\}$, and $\det\rho(c)=-1$ forces the eigenvalues $1$ and $-1$. The equivalence is the case $n=2$ of the distinct-eigenvalue
criterion of \Cref{ch:reps}: a line of $\bar k^2$ stable under $\rho(G_\Q)$ is an eigenline of $\rho(c)$, and both eigenlines are defined over $k$. For
$\det\rho=\bar\chi_\ell$, $\bar\chi_\ell(c)=-1\ne1$ in $\F_\ell$ as $\ell$ is odd.
\end{proof}

We do not prove (c). It follows from the theorem of Artin and Schreier that a field $K$ with $[\bar K:K]$ finite and greater than $1$ has characteristic $0$ and
$\bar K=K(i)$ (\cite[Chapter VI, Section 9]{Lang}), together with the theory of real closed fields (\cite[Chapter XI]{Lang}), which shows that such a $K$ is real closed;
both locations are pointers we have not re-verified. It is not used in the paper; what is used, in Proposition prop:ns and Lemmas lem:h0, lem:old and lem:dfg21, is
(d) and the fact that $\Gal(\C/\R)=\{1,c\}$ has order $2$ (so $H^1(\R,M)$ is killed by $2$, \Cref{ch:galcoh}).

\section[Continuous representations]{Continuous representations of profinite groups}\label{sec:galois:reps}

The paper's objects are continuous representations of $G_\Q$ and of its local subgroups: $\rho=\rho_{E,\ell}\colon G_\Q\to\GL_2(\Zl)$ on $T=T_\ell E$, its reductions
$\rho\bmod\ell^k$ and $\bar\rho$, and the characters $\chi_\ell$, $\chi_d$, $\varepsilon$. This section translates the representation theory of \Cref{ch:reps}
(which is about abstract groups) to profinite groups: what continuity means, the field cut out by a representation, twists, restriction to an open subgroup such as
$G_{\Q(\sqrt{\ell^*})}$ (the hypothesis of Theorem thm:main), the index-two criterion used in Lemma lem:h0 and Proposition prop:ns, independence of $\bar\rho^{ss}$ from the
lattice (Lemma lem:old), and the $\ell$-adic Kummer cocycle (Lemma lem:tamfactor).

\begin{definition}[Continuous representations]\label{def:galois:continuous}
Let $G$ be a profinite group. A finite field, $\Fbar_\ell$, or any field regarded as discrete, gives $\GL_n(k)$ the discrete topology. A finite extension $E_\lambda$ of
$\Q_\ell$, with ring of integers $\mathcal{O}$ and uniformiser $\varpi$ (\Cref{ch:local}), gives $\GL_n(E_\lambda)\subseteq E_\lambda^{n^2}$ the subspace topology of the $\ell$-adic topology;
the subgroups $1+\varpi^m\mathrm{M}_n(\mathcal{O})$, $m\ge1$, form a basis of neighbourhoods of $1$, and $\GL_n(\mathcal{O})$ is an open subgroup. A
\emph{representation} of $G$ is a continuous homomorphism $\rho\colon G\to\GL_n(k)$ for such a $k$, or $\rho\colon G\to\GL_n(\mathcal{O})$. For $G=G_F$, the field
\emph{cut out} by $\rho$ is $F_\rho=(F^{\mathrm{sep}})^{\ker\rho}$; $\rho$ \emph{factors through} $\Gal(F_\rho/F)$.
\end{definition}

\begin{proposition}[Continuous representations]\label{prop:galois:continuous-reps}
Let $G$ be profinite.
\begin{enumerate}[label=(\alph*)]
\item If $k$ is discrete, a homomorphism $\rho\colon G\to\GL_n(k)$ is continuous iff $\ker\rho$ is open, iff $\rho$ factors through a finite quotient $G/U$, $U$ open
normal; then $\rho(G)$ is finite. For $G=G_F$, $F_\rho/F$ is then finite Galois and $\Gal(F_\rho/F)\simeq\rho(G_F)$. In particular a representation of a finite quotient
of $G$ is, by inflation, a representation of $G$.
\item A homomorphism $\rho\colon G\to\GL_n(\mathcal{O})$ is continuous iff every $\rho\bmod\varpi^m$ is. Then $\rho(G)$ is compact, hence closed, $\ker\rho$ is a closed
normal subgroup, and for $G=G_F$, $F_\rho/F$ is Galois (in general infinite) with $\Gal(F_\rho/F)\simeq\rho(G_F)$ as topological groups.
\item Every continuous $\rho\colon G\to\GL_n(E_\lambda)$ stabilises an $\mathcal{O}$-lattice $\Lambda\subset E_\lambda^n$ (a free $\mathcal{O}$-submodule of rank $n$ spanning $E_\lambda^n$);
in a basis of $\Lambda$, $\rho$ is a continuous $\rho\colon G\to\GL_n(\mathcal{O})$, and its reduction $\bar\rho_\Lambda$ on $\Lambda/\varpi\Lambda$ is continuous.
\item Every continuous $\rho\colon G\to\GL_n(\C)$ has finite image. In particular continuous characters $G_\Q\to\C^\times$ have finite order.
\item Subrepresentations, quotients, direct sums, tensor products, duals, $\Hom$ and $\ad$ of continuous representations, the semisimplification $\rho^{ss}$, and
twists $\rho\otimes\chi$ by continuous characters are continuous.
\end{enumerate}
\end{proposition}

\begin{proof}
(a) $\{1\}$ is open in $\GL_n(k)$, so continuity implies $\ker\rho$ open; conversely if $\ker\rho$ is open, the preimage of any set is a union of cosets of
$\ker\rho$, hence open. An open subgroup has finite index and contains an open normal subgroup (\Cref{ch:groups}), and $\ker\rho$ is itself normal. The statement
about $F_\rho$ is \Cref{thm:galois:krull}(d), (e) with the open normal subgroup $\ker\rho$. For inflation: $\rho'\circ\pi$ with $\pi\colon G\to G/U$ has open kernel.
(b) A homomorphism of topological groups is continuous iff it is continuous at $1$, and $\rho^{-1}(1+\varpi^m\mathrm{M}_n(\mathcal{O}))=\ker(\rho\bmod\varpi^m)$. The image
of a compact group is compact; the kernel is the intersection of the closed subgroups $\ker(\rho\bmod\varpi^m)$. Then \Cref{thm:galois:krull}(e) gives
$\Gal(F_\rho/F)\simeq G_F/\ker\rho$, and the continuous bijection $G_F/\ker\rho\to\rho(G_F)$ is a homeomorphism (compact to Hausdorff).
(c) $\GL_n(\mathcal{O})$ is open in $\GL_n(E_\lambda)$, so $U=\rho^{-1}(\GL_n(\mathcal{O}))$, the stabiliser of $\mathcal{O}^n$, is an open subgroup, of finite index. Put
$\Lambda=\sum_{gU\in G/U}\rho(g)\mathcal{O}^n$, a finite sum (well defined since $\rho(u)\mathcal{O}^n=\mathcal{O}^n$ for $u\in U$). It is a finitely generated
$\mathcal{O}$-submodule of $E_\lambda^n$ containing $\mathcal{O}^n$ and stable under $G$; finitely generated and torsion-free over the principal ideal domain $\mathcal{O}$, it is
free (\Cref{ch:modules}), of rank $n$ as it spans $E_\lambda^n$ and lies in $\varpi^{-m}\mathcal{O}^n$ for some $m$. In a basis of $\Lambda$ the matrices $\rho(g)$ lie in
$\GL_n(\mathcal{O})$, and continuity is unchanged by the change of basis. The reduction is continuous by (b) and (a).
(d) Let $\|\cdot\|$ be the operator norm and $W=\{A\in\GL_n(\C):\|A-1\|<1\}$. We claim that $W$ contains no subgroup other than $\{1\}$. Let $A$ generate a subgroup
inside $W$. For an eigenvalue $\nu$ of $A$ and every $k\in\Z$, $\nu^k-1$ is an eigenvalue of $A^k-1$, so $|\nu^k-1|\le\|A^k-1\|<1$. Then $|\nu|=1$
(otherwise $|\nu^k|$ is unbounded for $k\to\infty$ or $k\to-\infty$), and all powers of $\nu$ lie on the arc $|z-1|<1$ of the unit circle, that is in
$\{e^{i\theta}:|\theta|<\pi/3\}$; if $\nu=e^{i\theta}\ne1$ with $0<|\theta|<\pi/3$, some multiple $k\theta$ lies in $[\pi/3,\pi]$ modulo $2\pi$, a contradiction. So
all eigenvalues are $1$ and $A=1+N$ with $N$ nilpotent. If $N\ne0$, let $m\ge1$ be maximal with $N^m\ne0$; then $A^k-1=\sum_{j=1}^m\binom kjN^j$, and as $k\to\infty$ the
coefficient $\binom km$ dominates, so $\|A^k-1\|\to\infty$ (compare $\|\cdot\|$ with the norm of the coefficients in the basis $N,\dots,N^m$ of their span), a
contradiction. Hence $A=1$, proving the claim. Now $\rho^{-1}(W)$ is an open neighbourhood of $1$, so it contains an open normal subgroup $U$ (\Cref{ch:groups}); $\rho(U)$
is a subgroup inside $W$, so $\rho(U)=1$ and $\rho$ factors through the finite group $G/U$.
(e) Each construction is given in suitable bases by matrices whose entries are polynomials in the entries of $\rho(g)$, $\rho'(g)$ and $\det\rho(g)^{-1}$; these are
continuous. A subrepresentation (or quotient) is given, in a basis adapted to it, by a block of the matrix. The semisimplification is the direct sum of the
composition factors (Jordan-H\"older, \Cref{ch:modules} and \Cref{ch:reps}), which are subquotients.
\end{proof}

\begin{proposition}[Quadratic characters of $G_\Q$]\label{prop:galois:quadratic-characters}
\begin{enumerate}[label=(\alph*)]
\item $d\mapsto\chi_d$ is an isomorphism $\Q^\times/\Q^{\times2}\simeq\Hom_{\mathrm{cont}}(G_\Q,\{\pm1\})$, with $\ker\chi_d=G_{\Q(\sqrt d)}$; the nontrivial continuous quadratic
characters of $G_\Q$ correspond to the quadratic fields.
\item $\chi_{dd'}=\chi_d\chi_{d'}$. For example $\chi_{-15}=\chi_{-3}\chi_5$, and $\chi_{-4}\chi_{-3}=\chi_{12}$ is the character of $\Q(\sqrt3)$ (written by its discriminant $12$;
the symbol $\chi_3$ is never used, \Cref{conv:galois:characters}); and for a finite set $S$ of odd primes $\chi_S=\prod_{q\in S}\chi_{q^*}$ is the character of $\Q\bigl(\sqrt{\prod_{q\in S}q^*}\bigr)$.
\item For a prime $p\nmid D$, $D$ the discriminant of $\Q(\sqrt d)$: $\chi_d(\Frob_p)=\bigl(\frac Dp\bigr)$ (Kronecker symbol), which is $1$ iff $p$ splits in $\Q(\sqrt d)$. In
particular $\chi_{-3}(\Frob_p)=1$ iff $p\equiv1\pmod3$, $\chi_{-4}(\Frob_p)=1$ iff $p\equiv1\pmod4$, and for odd $\ell\ne p$, $\chi_{\ell^*}(\Frob_p)=\bigl(\frac p\ell\bigr)$.
\item As a Dirichlet character (\Cref{ch:rings}), $\chi_d$ has conductor $|D|$: $\chi_{q^*}$ has conductor $q$, $\chi_{-3}$ conductor $3$, $\chi_{-4}$ conductor $4$.
\end{enumerate}
\end{proposition}

\begin{proof}[Proof of (a), (b) and of the last formula of (c); (c), (d) otherwise stated]
(a) is \Cref{cor:galois:kummer-absolute} with $F=\Q$, $n=2$, and \Cref{prop:galois:quadratic}(c). (b) $\sqrt{dd'}=\pm\sqrt d\sqrt{d'}$, and $\chi_{dd'}$ does not depend on
the choice of square root. The last formula of (c) is a direct consequence of $\chi_{\ell^*}=\bigl(\frac{\bar\chi_\ell}\ell\bigr)$ (\Cref{prop:galois:sqrt-lstar}) and
$\bar\chi_\ell(\Frob_p)=p\bmod\ell$ (\Cref{thm:galois:frobenius}(c)). The general formula in (c) and (d) are proved in \Cref{ch:ant} and \Cref{ch:rings}; we do not prove
them here. Comparing $\chi_{\ell^*}(\Frob_p)=\bigl(\frac p\ell\bigr)$ with $\chi_{\ell^*}(\Frob_p)=\bigl(\frac{\ell^*}p\bigr)$ from (c) is quadratic reciprocity (\Cref{ch:rings}).
\end{proof}

The formulas of (c) were also checked numerically: for $5\le p<2000$, $x^2+3$ has a root modulo $p$ iff $p\equiv1\pmod 3$, $x^2+1$ iff $p\equiv1\pmod4$, and for
$\ell\in\{3,5,7,11,13\}$, $\ell^*$ is a square modulo $p$ iff $p$ is a square modulo $\ell$.

\begin{proposition}[Twists]\label{prop:galois:twist}
Let $\rho\colon G\to\GL_n(k)$ and $\chi\colon G\to k^\times$ be continuous ($k$ a field or $\mathcal{O}$ as above).
\begin{enumerate}[label=(\alph*)]
\item $\rho\otimes\chi\colon g\mapsto\chi(g)\rho(g)$ is continuous, $\tr(\rho\otimes\chi)(g)=\chi(g)\tr\rho(g)$ and $\det(\rho\otimes\chi)=\chi^n\det\rho$. For $G=G_F$ the field
cut out by $\rho\otimes\chi$ lies in $F_\rho F_\chi$.
\item If $n=2$ and $\rho\simeq\rho\otimes\chi$, then $\chi^2=1$. If moreover $k$ is a field, $\chi\ne1$ and $\rho$ is absolutely irreducible, every isomorphism
$\rho\to\rho\otimes\chi$ has trace $0$, and these isomorphisms form a line together with $0$.
\item Let $\ell$ be odd, $\rho\colon G_\Q\to\GL_2(\Zl)$ continuous and $k\ge1$. If $\rho\bmod\ell^k\simeq(\rho\otimes\chi_\ell)\bmod\ell^k$ (conjugate by an element of
$\GL_2(\Z/\ell^k)$), then $\chi_\ell(g)^2\equiv1\pmod{\ell^k}$ for all $g\in G_\Q$, and this forces $\ell=3$ and $k=1$. For $\ell=3$, $k=1$ the condition reads
$\bar\rho\simeq\bar\rho\otimes\bar\chi_3=\bar\rho\otimes\chi_{-3}$.
\end{enumerate}
\end{proposition}

\begin{proof}
(a) The formulas are immediate; continuity is \Cref{prop:galois:continuous-reps}(e). The kernel of $\rho\otimes\chi$ contains $\ker\rho\cap\ker\chi=G_{F_\rho F_\chi}$
(\Cref{thm:galois:krull}(c) and \Cref{thm:galois:fundamental}(e) in the limit).
(b) Taking determinants of $P\rho(g)P^{-1}=\chi(g)\rho(g)$ gives $\det\rho(g)=\chi(g)^2\det\rho(g)$, and $\det\rho(g)$ is a unit. The rest is the trace-zero lemma of
\Cref{ch:reps}: a $G$-map $\phi\colon\rho\to\rho\otimes\chi$ satisfies $(1-\chi(g))\tr\phi=0$, and $\Hom_G(\rho,\rho\otimes\chi)$ has dimension at most one with nonzero
elements isomorphisms when $\rho$ is absolutely irreducible.
(c) By (b) over the ring $\Z/\ell^k$ (the determinant argument works over any commutative ring), $\chi_\ell(g)^2\equiv1\pmod{\ell^k}$ for all $g$. By
\Cref{prop:galois:cyclotomic-character}(b) there is $g$ with $\chi_\ell(g)=2$ ($2\in\Zl^\times$ as $\ell$ is odd), so $\ell^k\mid4-1=3$: $\ell=3$ and $k=1$. Finally
$\bar\chi_3=\chi_{-3}$ by \Cref{prop:galois:sqrt-lstar}(c).
\end{proof}

\begin{paperlink}
Lemma lem:h0 identifies $h_\ell(E)$, the length of $H^0(\Q,\ad T(1)\otimes\Q_\ell/\Z_\ell)$, with the largest $k$ such that $\rho\bmod\ell^k\simeq(\rho\otimes\chi_\ell)
\bmod\ell^k$ (the module-theoretic part, Nakayama and cyclicity, is \Cref{ch:reps}). \Cref{prop:galois:twist}(c) is then the step ``Taking determinants,
$\chi_\ell(g)^2\equiv1\pmod{\ell^k}$ for all $g\in G_\Q$. Since $\chi_\ell(G_\Q)=\Z_\ell^\times$ contains $2$, \dots this forces $k=0$ for $\ell\ge5$ and $k\le1$ for
$\ell=3$'': $h_\ell(E)=0$ for $\ell\ge5$ and $h_3(E)\le1$, for every $E$ with $\bar\rho$ irreducible, with or without complex multiplication. It holds for every continuous
$\rho\colon G_\Q\to\GL_2(\Zl)$, irreducible or not. The original version of the paper argued instead with the quadratic character $\chi_{\ell^*}$, which differs from
$\bar\chi_\ell$ for $\ell\ge5$ (\Cref{prop:galois:sqrt-lstar}); the revised Lemma lem:h0 uses the argument above.
\end{paperlink}

\begin{proposition}[Restriction to open subgroups]\label{prop:galois:restriction}
Let $K/F$ be a finite separable extension inside $F^{\mathrm{sep}}$ and $\rho\colon G_F\to\GL_n(k)$ continuous.
\begin{enumerate}[label=(\alph*)]
\item $\Res\rho=\rho|_{G_K}$ is continuous, and $[\rho(G_F):\rho(G_K)]\le[K:F]$. If $K/F$ is Galois, $\rho(G_K)$ is normal in $\rho(G_F)$ and the index divides $[K:F]$.
\item (Image of a kernel.) If $\delta\colon\GL_n(k)\to A$ is a homomorphism to an abelian group and $G_K=\ker(\delta\circ\rho)$ (that is, $K$ is cut out by $\delta\circ\rho$),
then $\rho(G_K)=\rho(G_F)\cap\ker\delta$.
\item Let $E/\Q$ be an elliptic curve, $\ell$ odd and $\bar\rho=\bar\rho_{E,\ell}$, with $\det\bar\rho=\bar\chi_\ell$ (\Cref{ch:ellcurves}). Then
$\bar\rho(G_{\Q(\zeta_\ell)})=\bar\rho(G_\Q)\cap\SL_2(\F_\ell)$ and $\bar\rho(G_{\Q(\sqrt{\ell^*})})=\{g\in\bar\rho(G_\Q):\det g\in(\F_\ell^\times)^2\}$. For $\ell=3$ both are the
determinant-one subgroup of the image.
\item The field cut out by $\Res\rho$ is $KF_\rho$, and $\Gal(KF_\rho/K)\simeq\Gal(F_\rho/F_\rho\cap K)$.
\item $\rho$ is absolutely irreducible on $G_K$ iff the subgroup $\rho(G_K)$ fixes no proper nonzero subspace of $\bar k^n$; this depends only on $\rho(G_K)$.
\end{enumerate}
\end{proposition}

\begin{proof}
(a) Restriction of a continuous map. The cosets of $G_K$ (there are $[K:F]$ of them, \Cref{prop:galois:absolute}(a)) map onto the cosets of $\rho(G_K)$. If $K/F$ is Galois,
$G_K$ is normal, so $\rho(G_K)$ is normal in $\rho(G_F)$, and $\rho(G_F)/\rho(G_K)$ is a quotient of $G_F/G_K\simeq\Gal(K/F)$.
(b) If $g\in G_K$ then $\delta(\rho(g))=1$. Conversely if $\rho(g)\in\ker\delta$ then $g\in\ker(\delta\circ\rho)=G_K$; so every element of $\rho(G_F)\cap\ker\delta$ is in
$\rho(G_K)$.
(c) $\Q(\zeta_\ell)$ is cut out by $\bar\chi_\ell=\det\circ\bar\rho$ (\Cref{prop:galois:cyclotomic-character}(c)): take $\delta=\det$. $\Q(\sqrt{\ell^*})$ is cut out by
$\chi_{\ell^*}=\bigl(\frac{\bar\chi_\ell}{\ell}\bigr)=\bigl(\frac{\det\bar\rho}{\ell}\bigr)$ (\Cref{prop:galois:sqrt-lstar}): take $\delta=\bigl(\frac{\det}{\ell}\bigr)$. For
$\ell=3$, $(\F_3^\times)^2=\{1\}$.
(d) The kernel of $\Res\rho$ is $G_K\cap\ker\rho=G_K\cap G_{F_\rho}=G_{KF_\rho}$ (\Cref{thm:galois:krull}). The isomorphism is \Cref{cor:galois:translation-infinite}.
(e) Definition (\Cref{conv:galois:reps}); the stable subspaces of $\bar k^n$ under $G_K$ are those under $\rho(G_K)$.
\end{proof}

\begin{proposition}[Index two]\label{prop:galois:index-two}
Let $G$ be a profinite group, $H\subseteq G$ an open subgroup of index $2$, $\chi\colon G\to G/H\simeq\{\pm1\}$, $s\in G\setminus H$, $k$ a field with
$\operatorname{char}k\ne2$ (discrete, or a finite extension of $\Q_\ell$), and $\rho\colon G\to\GL_2(k)$ continuous and absolutely irreducible. The following are
equivalent:
\begin{enumerate}[label=(\roman*)]
\item $\rho|_H$ is not absolutely irreducible;
\item $\rho\simeq\rho\otimes\chi$ (over $k$; equivalently over $\bar k$);
\item $\rho\otimes\bar k\simeq\Ind_H^G\psi$ for a continuous character $\psi\colon H\to\bar k^\times$.
\end{enumerate}
When they hold, $\rho|_H\otimes\bar k\simeq\psi\oplus\psi^s$ with $\psi^s(h)=\psi(s^{-1}hs)$ and $\psi^s\ne\psi$, and every isomorphism $\rho\to\rho\otimes\chi$ has trace $0$.
For $G=G_\Q$ and $H=G_K$, $K$ quadratic, $\chi=\chi_d$ is the character of $K=\Q(\sqrt d)$.
\end{proposition}

\begin{proof}
This is the Clifford theory for a subgroup of index two of \Cref{ch:reps}, applied to the abstract group $G$; the equivalence (ii) over $k$ and over $\bar k$ holds
because spaces of $G$-maps commute with extension of scalars (\Cref{ch:reps}). Only continuity needs comment. In (iii), $\psi$ is the character by which $H$ acts on a
line $\bar kw\subset\bar k^2$ stable under $\rho(H)$. The coordinates of $w$ generate a finite extension $k'$ of $k$, so $\psi(h)=\rho(h)w/w\in k'$, and in a basis $w,w'$
of $k'^2$ the value $\psi(h)$ is the upper left entry of the matrix of $\rho(h)$, a continuous function of $h$ ($k'$ discrete if $k$ is, and with its $\ell$-adic
topology if $k$ is a finite extension of $\Q_\ell$). Conversely, $\Ind_H^G\psi$ for continuous $\psi$ is continuous: in the basis $\{v,sv\}$ its matrix entries are values of $\psi$ at
$h$, $s^{-1}hs$ or products with $\psi(s^2)$, all continuous functions of $h\in H$ and $g\in G$ ($H$ being open and closed).
\end{proof}

\begin{paperlink}
For $\bar\rho=\bar\rho_{E,\ell}$, irreducible, hence absolutely irreducible by \Cref{prop:galois:complex-conjugation}(d), and $K=\Q(\sqrt{\ell^*})$: condition (i) is the
failure of the hypothesis of Theorem thm:main (DFG04's $\ell\in S_f$). For $\ell=3$, $\chi_{-3}=\bar\chi_3$, and (ii) says that the space of $G_\Q$-maps
$\bar\rho\to\bar\rho\otimes\bar\chi_3$, which are automatically of trace zero, is nonzero: $H^0(\Q,\ad\bar\rho(1))\ne0$, the condition $h_3(E)=1$ of Lemma lem:h0; and by
(iii) it happens exactly when $\bar\rho\otimes\F_9$ is induced from $G_{\Q(\sqrt{-3})}$, the case $N_s(\F_3)$ of Proposition prop:ns (\Cref{prop:galois:quartic-on-lines}(f)
and \Cref{ch:reps}: the determinant-one subgroup $C_4$ of $N_s(\F_3)$ is abelian, so $\bar\rho|_{G_{\Q(\sqrt{-3})}}\otimes\F_9=\psi\oplus\psi^{-1}$ with $\psi$ of order $4$). For
$\ell\ge5$, (ii) for $\chi_{\ell^*}$ can hold (the dihedral case excluded by $S_f$; by Proposition prop:Sf only when $\ell\mid N$), but it is not the condition $\bar\rho\simeq\bar\rho\otimes\bar\chi_\ell$ of Lemma lem:h0,
which never holds for $\ell\ge5$ (\Cref{prop:galois:twist}(b): it would force $\bar\chi_\ell^2=1$).
\end{paperlink}

\begin{proposition}[Reduction and semisimplification]\label{prop:galois:reduction-ss}
Let $E_\lambda$ be a finite extension of $\Q_\ell$ with ring of integers $\mathcal{O}$, uniformiser $\varpi$ and residue field $k$, $G$ a profinite group, and
$\rho\colon G\to\GL_n(E_\lambda)$ continuous. For two $G$-stable lattices $\Lambda,\Lambda'$, the reductions $\bar\rho_\Lambda$, $\bar\rho_{\Lambda'}$ have the same characteristic
polynomials, namely the reductions of the characteristic polynomials of $\rho$, which lie in $\mathcal{O}[X]$; hence $\bar\rho_\Lambda^{ss}\simeq\bar\rho_{\Lambda'}^{ss}$, and
$\bar\rho^{ss}$ is well defined. If $\bar\rho_\Lambda$ is irreducible, then $\Lambda'=\varpi^m\Lambda$ for some $m\in\Z$.
\end{proposition}

\begin{proof}
In a basis of $\Lambda$ the matrix of $\rho(g)$ lies in $\GL_n(\mathcal{O})$, so $\det(X-\rho(g))\in\mathcal{O}[X]$, and it does not depend on the basis of $E_\lambda^n$; its
reduction is the characteristic polynomial of $\bar\rho_\Lambda(g)$, and likewise for $\Lambda'$. Brauer-Nesbitt (\Cref{ch:reps}, valid for any group and any field)
gives $\bar\rho_\Lambda^{ss}\simeq\bar\rho_{\Lambda'}^{ss}$. The last statement is the homothety lemma of \Cref{ch:reps}.
\end{proof}

\begin{proposition}[The $\ell$-adic Kummer cocycle]\label{prop:galois:l-adic-kummer}
Let $\ell\ne\operatorname{char}F$ be a prime and $a\in F^\times$. Choose a compatible system $(\alpha_k)_{k\ge0}$ in $F^{\mathrm{sep}}$ with $\alpha_0=a$ and
$\alpha_{k+1}^\ell=\alpha_k$, and define, in the additive notation of \Cref{conv:galois:characters}(c),
\[
\kappa_a\colon G_F\to\Zl(1),\qquad\kappa_a(\sigma)=\bigl(\sigma(\alpha_k)/\alpha_k\bigr)_{k\ge0}.
\]
\begin{enumerate}[label=(\alph*)]
\item $\kappa_a$ is continuous and $\kappa_a(\sigma\tau)=\kappa_a(\sigma)+\chi_\ell(\sigma)\kappa_a(\tau)$: a continuous $1$-cocycle (\Cref{ch:cohomology}). Equivalently, writing
$\kappa_a(\sigma)=b(\sigma)e$ for a generator $e$ of $\Zl(1)$,
\[
\sigma\longmapsto\begin{pmatrix}\chi_\ell(\sigma)&b(\sigma)\\0&1\end{pmatrix}
\]
is a continuous homomorphism $G_F\to\GL_2(\Zl)$: the action on $\Zl e_1\oplus\Zl e_2$ with $e_1=e$, $\sigma e_1=\chi_\ell(\sigma)e_1$, $\sigma e_2=e_2+b(\sigma)e_1$.
\item Another compatible system changes $\kappa_a$ by the coboundary $\sigma\mapsto(\chi_\ell(\sigma)-1)z$ for some $z\in\Zl(1)$. For the product systems,
$\kappa_{aa'}=\kappa_a+\kappa_{a'}$ for $a,a'\in F^\times$. The $k$-th component of $\kappa_a$, restricted to $G_{F(\mu_{\ell^k})}$, is the Kummer character $\chi_a$ with $n=\ell^k$.
\item On $\ker\chi_\ell=G_{F(\mu_{\ell^\infty})}$, $\kappa_a$ is a homomorphism, and for $\sigma\in G_F$, $\tau\in\ker\chi_\ell$:
$\kappa_a(\sigma\tau\sigma^{-1})=\chi_\ell(\sigma)\kappa_a(\tau)$.
\item If $K\supseteq F$ contains $\mu_{\ell^k}$, then $\kappa_a\bmod\ell^k$ vanishes on $G_K$ iff $a\in(K^\times)^{\ell^k}$.
\end{enumerate}
\end{proposition}

\begin{proof}
(a) $(\sigma(\alpha_k)/\alpha_k)^{\ell^k}=\sigma(a)/a=1$, and $(\sigma(\alpha_{k+1})/\alpha_{k+1})^\ell=\sigma(\alpha_k)/\alpha_k$, so $\kappa_a(\sigma)\in\Zl(1)$. Multiplicatively,
\[
\frac{\sigma\tau(\alpha_k)}{\alpha_k}=\sigma\Bigl(\frac{\tau(\alpha_k)}{\alpha_k}\Bigr)\cdot\frac{\sigma(\alpha_k)}{\alpha_k},
\]
and $\sigma$ acts on $\mu_{\ell^k}$ by raising to the power $\chi_\ell(\sigma)$; additively this is the cocycle relation. The $k$-th component is constant on the cosets
of the open subgroup $G_{F(\mu_{\ell^k},\alpha_k)}$ (for $\tau$ in it, $\kappa_a(\sigma\tau)_k=\kappa_a(\sigma)_k+\chi_\ell(\sigma)\kappa_a(\tau)_k=\kappa_a(\sigma)_k$), so $\kappa_a$ is
continuous. The matrix identity $\sm{\chi(\sigma)}{b(\sigma)}01\sm{\chi(\tau)}{b(\tau)}01=\sm{\chi(\sigma\tau)}{b(\sigma)+\chi(\sigma)b(\tau)}01$ is the cocycle relation.
(b) Another system is $(\epsilon_k\alpha_k)$ with $\epsilon=(\epsilon_k)\in\Zl(1)$, and $\sigma(\epsilon_k\alpha_k)/\epsilon_k\alpha_k=(\sigma(\epsilon_k)/\epsilon_k)\cdot\sigma(\alpha_k)/
\alpha_k$; additively the change is $\sigma(\epsilon)-\epsilon=(\chi_\ell(\sigma)-1)\epsilon$. Additivity in $a$ is clear. On $G_{F(\mu_{\ell^k})}$ the $k$-th component is
$\sigma(\alpha_k)/\alpha_k=\chi_a(\sigma)$ by definition (\Cref{prop:galois:kummer-character}, over $F(\mu_{\ell^k})$, with $n=\ell^k$, $\alpha=\alpha_k$).
(c) For $\tau,\tau'\in\ker\chi_\ell$ the cocycle relation reads $\kappa_a(\tau\tau')=\kappa_a(\tau)+\kappa_a(\tau')$. From $0=\kappa_a(1)=\kappa_a(\sigma\sigma^{-1})=\kappa_a(\sigma)+
\chi_\ell(\sigma)\kappa_a(\sigma^{-1})$ and $\chi_\ell(\tau)=1$:
\[
\kappa_a(\sigma\tau\sigma^{-1})=\kappa_a(\sigma)+\chi_\ell(\sigma)\bigl(\kappa_a(\tau)+\kappa_a(\sigma^{-1})\bigr)=\chi_\ell(\sigma)\kappa_a(\tau).
\]
(d) $G_K\subseteq G_{F(\mu_{\ell^k})}$, so by (b) $\kappa_a\bmod\ell^k$ on $G_K$ is the Kummer character of $a$ over $K$ with $n=\ell^k$; it is trivial iff $a\in(K^\times)^{\ell^k}$
by \Cref{prop:galois:kummer-character}(b).
\end{proof}

\begin{paperlink}
In the proof of Lemma lem:tamfactor, $F=\Q_q$ with $q\ne\ell$, $a=q_E$ is the Tate parameter, and over the unramified quadratic extension the action on $T_\ell E$ in the
basis $e_1=(\zeta_{\ell^k})_k$, $e_2=(q_E^{1/\ell^k})_k$ is the matrix of (a) (\Cref{ch:reduction}); for the nonsplit curve it is twisted by the unramified character
$\varepsilon$, giving $\sm{\chi_\ell\varepsilon}{\varepsilon\kappa}0\varepsilon$. Part (c) with $\sigma=\Frob$ and $\tau=t$ a generator of the pro-$\ell$ tame inertia quotient (on which
$\chi_\ell$ is trivial, \Cref{ch:local}) gives $\kappa(\Frob\,t\,\Frob^{-1})=q\kappa(t)$, the abstract form of the tame relation $\Frob\,t\,\Frob^{-1}=t^q$, equivalently
$\Frob^{-1}t\,\Frob=t^{q^{-1}}$, which the paper uses. By (b), $\kappa_{q_E}=n\kappa_q+\kappa_u$ for $q_E=q^nu$; the restriction of $\kappa_u$ to inertia vanishes (the unit $u$
has $\ell$-power roots in $\Q_q^{\mathrm{ur}}$, by Hensel, \Cref{ch:local}), and the restriction of $\kappa_q$ to inertia is the tame character, whence ``the restriction of $\kappa$
to $I$ is $n$ times the tame character''. In Lemma lem:tamagawa-exp, (d) with $K=\Q_q^{\mathrm{ur}}$ and $G_K=I_q$ says that inertia acts trivially on $E[\ell^t]$ iff
$q_E\in(\Q_q^{\mathrm{ur},\times})^{\ell^t}$; \Cref{ch:local} turns this into $\ell^t\mid n_q$. The value $c'=\varepsilon\kappa(\Frob)$ depends on the choices (by (b) and on
the lift of $\Frob$), and only its class modulo $(n,q-1)\Zl$ is meaningful, which does not affect the minimum in the lemma (\Cref{ch:galcoh}).
\end{paperlink}

\section{Frobenius elements and the Chebotarev density theorem}\label{sec:galois:chebotarev}

The step ``$\bar\rho_h\simeq\bar\rho$ by Chebotarev and Brauer-Nesbitt'' in Lemma lem:old compares two residual representations through the characteristic polynomials
of Frobenius elements, which are known only at the primes not dividing $\ell N$. This section states what Frobenius elements are (they are constructed in
\Cref{ch:ant}; no theorem of Sections 1 to 7 depends on them, and in Sections 6 and 8 they enter only through the stated \Cref{thm:galois:frobenius}, in
\Cref{ex:galois:ns-example} and \Cref{prop:galois:quadratic-characters}(c)), states the Chebotarev density theorem, and proves the consequences used in the paper: Frobenius elements
are dense, so a semisimple representation is determined by their characteristic polynomials, and congruent traces give isomorphic residual representations. A
worked example checks the factorisation patterns of $\psi_3$ against Frobenius, as the test of Section sec:data does.

\begin{theorem}[Frobenius elements]\label{thm:galois:frobenius}
Let $K/\Q$ be a finite Galois extension with ring of integers $\mathcal{O}_K$, and $p$ a prime unramified in $K$ (\Cref{ch:ant}).
\begin{enumerate}[label=(\alph*)]
\item For each prime $\mathfrak P$ of $\mathcal{O}_K$ above $p$ there is a unique $\Frob_{\mathfrak P}\in\Gal(K/\Q)$ with $\Frob_{\mathfrak P}(x)\equiv x^p\pmod{\mathfrak P}$
for all $x\in\mathcal{O}_K$. The $\Frob_{\mathfrak P}$, $\mathfrak P\mid p$, form one conjugacy class $\Frob_p$ ($\Frob_{\sigma\mathfrak P}=\sigma\Frob_{\mathfrak P}\sigma^{-1}$). The order
of $\Frob_{\mathfrak P}$ is the residue degree $[\mathcal{O}_K/\mathfrak P:\F_p]$, and $p$ splits completely in $K$ iff $\Frob_p=1$.
\item For $K'\subseteq K$ Galois over $\Q$, $\Frob_{\mathfrak P}$ restricts to $\Frob_{\mathfrak P\cap K'}$; so $\Frob_p(K/\Q)$ maps to $\Frob_p(K'/\Q)$.
\item In $\Q(\zeta_n)$ with $p\nmid n$, $\Frob_p=\sigma_p$; hence $\bar\chi_\ell(\Frob_p)=p\bmod\ell$ and $\chi_\ell(\Frob_p)=p$ for $p\ne\ell$. In $\Q(\sqrt d)$ with $p\nmid D$,
$\chi_d(\Frob_p)=\bigl(\frac Dp\bigr)$.
\item For $L/\Q$ Galois of any degree, unramified at $p$ (that is, every finite subextension is), the Frobenius elements at the primes of $L$ above $p$ (compatible
systems of primes of the finite subextensions) are the elements $\Frob_{\mathfrak P}\in\Gal(L/\Q)$ restricting to $\Frob_{\mathfrak P\cap K}$ for each finite Galois $K\subseteq L$;
they form one conjugacy class. In $G_\Q$ a Frobenius element at $p$ is well defined up to conjugation and up to the inertia group at a prime above $p$; so for a
continuous representation $\rho$ of $G_\Q$ unramified at $p$ (trivial on the inertia groups at $p$), the conjugacy class of $\rho(\Frob_p)$, and in particular its
characteristic polynomial, is well defined.
\item (Dedekind.) Let $f\in\Z[x]$ be separable with leading coefficient $a$, and $p\nmid a\disc(f)$. If $f\bmod p$ has irreducible factors of degrees $d_1,\dots,d_r$, then
$p$ is unramified in the splitting field of $f$ and $\Frob_p$ acts on the roots of $f$ with cycle type $(d_1,\dots,d_r)$.
\end{enumerate}
\end{theorem}

We do not prove this here; see \Cref{ch:ant}, and \cite[Chapter I, Section 9]{Neukirch} (decomposition groups and Frobenius) and \cite[Chapter I, Section 8]{Neukirch}
(the factorisation of $p$ through $f\bmod p$); the section numbers in this book are pointers we have not re-verified. The finite-field part is proved here: the
residue extension $(\mathcal{O}_K/\mathfrak P)/\F_p$ is cyclic generated by $\Frob_p$ (\Cref{thm:galois:finite-fields}), and $\Frob_{\mathfrak P}$ is the lift of that
generator, which exists and is unique when $p$ is unramified. Note the convention (\Cref{conv:galois:frobenius}): $\Frob_p$ is arithmetic.

\begin{example}[Frobenius on the four lines]\label{ex:galois:frobenius-on-lines}
Let $E/\Q$ have good reduction at a prime $p\ne3$ with $p\nmid\disc(\psi_3)$. Then $\bar\rho_{E,3}$ is unramified at $p$ and $\bar\rho_{E,3}(\Frob_p)$ has characteristic
polynomial $X^2-a_pX+p\bmod3$ (\Cref{ch:ellcurves}, \Cref{ch:galrep}). By \Cref{thm:galois:frobenius}(e) and \Cref{prop:galois:quartic-on-lines}, the degrees of the
irreducible factors of $\psi_3\bmod p$ are the orbit sizes of $g=\bar\rho_{E,3}(\Frob_p)$ on the four lines. These depend only on the characteristic polynomial, except
in the last case:
\begin{itemize}
\item $X^2\pm X-1$: the eigenvalues $\lambda,\lambda^3\in\F_9\setminus\F_3$ satisfy $\lambda^4=\det g=-1$, so $\lambda$ has order $8$ and $g$ has order $4$ in $\PGL_2(\F_3)$ with
no fixed line: a $4$-cycle, pattern $(4)$;
\item $X^2+1$: $g^2=-1$, so $g$ is an involution on the lines, without fixed line (no eigenvalue in $\F_3$): pattern $(2,2)$;
\item $X^2-1$: $g\sim\operatorname{diag}(1,-1)$ fixes its two eigenlines and swaps the other two: pattern $(1,1,2)$;
\item $(X\mp1)^2$: $g=\pm(1+N)$ with $N$ nilpotent; if $N\ne0$, $g$ fixes one line and has order $3$ in $\PGL_2(\F_3)$: pattern $(1,3)$; if $g=\pm1$: $(1,1,1,1)$.
\end{itemize}
For 11a1 ($a_2=-2$, $a_5=1$, $a_7=-2$, $a_{13}=4$, $a_{17}=-2$, $a_{19}=0$, $a_{23}=-1$, $a_{29}=0$, $a_{31}=7$, $a_{37}=3$; $\disc\psi_3=-3^3\cdot11^{10}$) we checked with PARI that
$\psi_3\bmod p$ has the pattern $(4)$ for $p=2,5,17,23$, $(1,3)$ for $p=7,13,31$, $(2,2)$ for $p=19,37$ and $(1,1,2)$ for $p=29$, as predicted. The pattern $(1,1,2)$ at
$p=29$ means that $\Frob_{29}$ fixes two lines, although the image is $\GL_2(\F_3)$: Frobenius elements see conjugacy classes, not the image.
\end{example}

\begin{definition}[Density]\label{def:galois:density}
Let $X$ be a set of primes. Its \emph{Dirichlet density} is
\[
\lim_{s\to1^+}\frac{\sum_{p\in X}p^{-s}}{\log\frac1{s-1}},
\]
and its \emph{natural density} is $\lim_{x\to\infty}\#\{p\le x:p\in X\}/\#\{p\le x\}$, when these limits exist. A set with a natural density has the same Dirichlet density
(we do not use the converse, which is false). Finite sets have density $0$, a subset of a set of density $0$ has density $0$ (the sums are nonnegative), and densities are
additive on disjoint unions of sets that have densities. $X$ has \emph{density one} if its complement has density $0$. ``\emph{Almost all} $p$'' means all but finitely
many. Unless stated otherwise, density means Dirichlet density.
\end{definition}

For the comparison of the two densities see \cite[Chapter VI, Section 4]{SerreCourse}; the normalisation uses $\sum_pp^{-s}\sim\log\frac1{s-1}$ as $s\to1^+$, which is
proved there.

\begin{theorem}[Chebotarev's density theorem]\label{thm:galois:chebotarev}
Let $K/\Q$ be a finite Galois extension with group $G$ and $C\subseteq G$ a conjugacy class. The set of primes $p$, unramified in $K$, with $\Frob_p=C$ has Dirichlet
density $\#C/\#G$; it also has natural density $\#C/\#G$. The same holds for a finite Galois extension of number fields, with the primes of the base field. In particular
every conjugacy class is $\Frob_p$ for infinitely many $p$.
\end{theorem}

We do not prove this; see \cite[Chapter VII, Section 13]{Neukirch} for the Dirichlet density and \cite[Chapter I, Section 2.2]{SerreAbelian} for the statement with
natural density; the book uses only the positivity of the Dirichlet density. Special cases: for $K=\Q(\zeta_n)$ (\Cref{thm:galois:frobenius}(c)) it is Dirichlet's
theorem, each class $a\in(\Z/n)^\times$ containing primes of density $1/\varphi(n)$; for $K=\Q(\sqrt d)$, half of the primes split. For the splitting field of $x^3-2$
(\Cref{ex:galois:x3minus2}, group $S_3$) and $p\nmid6$, \Cref{thm:galois:frobenius}(e) says that $x^3-2$ has exactly one root modulo $p$ iff $\Frob_p$ is a
transposition (density $1/2$), three roots iff $\Frob_p=1$ (density $1/6$), and none iff $\Frob_p$ is a $3$-cycle (density $1/3$). A transposition is odd, and the sign
character cuts out $\Q(\sqrt{-108})=\Q(\sqrt{-3})$, so ``exactly one root'' happens iff $\chi_{-3}(\Frob_p)=-1$, iff $p\equiv2\pmod3$. Among the $428$ primes $5\le p<3000$
the observed proportions are $0.516$, $0.159$ and $0.325$.

\begin{corollary}[Frobenius elements are dense]\label{cor:galois:chebotarev-dense}
Let $L/\Q$ be a Galois extension (possibly infinite) unramified outside a finite set $S$ of primes, and $X$ a set of primes of density one. Then the Frobenius elements
$\Frob_{\mathfrak P}\in\Gal(L/\Q)$ at the primes $\mathfrak P$ of $L$ above the $p\in X\setminus S$ form a dense subset of $\Gal(L/\Q)$.
\end{corollary}

\begin{proof}
A basic open set of $\Gal(L/\Q)$ is $\sigma\Gal(L/K)$ with $K\subseteq L$ finite Galois over $\Q$ (\Cref{prop:galois:infinite}(b)). Let $C$ be the conjugacy class of
$\sigma|_K$ in $\Gal(K/\Q)$. By \Cref{thm:galois:chebotarev} the set $Y$ of primes unramified in $K$ with $\Frob_p(K/\Q)=C$ has density $\#C/[K:\Q]>0$. The complement of
$X\setminus S$ has density $0$ (the complement of $X$ has density $0$ and $S$ is finite), so $Y\cap(X\setminus S)$ has the positive density of $Y$, and in particular contains
a prime $p$. As $\Frob_p(K/\Q)=C$ is the set of $\Frob_{\mathfrak p}$ for the primes $\mathfrak p$ of $K$ above $p$, there is $\mathfrak p$ with $\Frob_{\mathfrak p}=\sigma|_K$; choose a
prime $\mathfrak P$ of $L$ above $\mathfrak p$ (\Cref{ch:ant}). By \Cref{thm:galois:frobenius}(d), $\Frob_{\mathfrak P}$ restricts to $\Frob_{\mathfrak p}=\sigma|_K$ on $K$, that is
$\Frob_{\mathfrak P}\in\sigma\Gal(L/K)$.
\end{proof}

\begin{corollary}[Chebotarev and Brauer-Nesbitt]\label{cor:galois:chebotarev-bn}
Let $k$ be a finite field, $\Fbar_\ell$, or a finite extension of $\Q_\ell$, and $\rho_1,\rho_2\colon G_\Q\to\GL_n(k)$ continuous, semisimple, and unramified outside a finite
set $S$ of primes. If $\det(X-\rho_1(\Frob_p))=\det(X-\rho_2(\Frob_p))$ for all $p$ in a set $X$ of density one, $p\notin S$, then $\rho_1\simeq\rho_2$.
\end{corollary}

\begin{proof}
Let $L$ be the field cut out by $\rho_1\oplus\rho_2$, the fixed field of $\ker\rho_1\cap\ker\rho_2$ (\Cref{prop:galois:continuous-reps}); it is Galois over $\Q$ and unramified
outside $S$, since the inertia groups at $p\notin S$ act trivially through both $\rho_i$. Both $\rho_i$ factor through $\Gal(L/\Q)$. The map
\[
\Gal(L/\Q)\to k^n,\qquad g\mapsto\text{coefficients of }\det(X-\rho_1(g))-\det(X-\rho_2(g)),
\]
is continuous (polynomial in the matrix entries), so its zero set is closed; it contains the Frobenius elements at the primes above $X\setminus S$, which are dense by
\Cref{cor:galois:chebotarev-dense}. Hence $\rho_1(g)$ and $\rho_2(g)$ have the same characteristic polynomial for every $g\in G_\Q$. By Brauer-Nesbitt (\Cref{ch:reps}, for an
arbitrary group over an arbitrary field) $\rho_1^{ss}\simeq\rho_2^{ss}$, that is $\rho_1\simeq\rho_2$ as both are semisimple. (For finite $k$ the image is finite and
the finite-group form suffices.)
\end{proof}

\begin{corollary}[Congruent traces]\label{cor:galois:congruent-ap}
Let $\mathcal{O}$ be the ring of integers of a finite extension of $\Q_\ell$, $\lambda=\varpi\mathcal{O}$ its maximal ideal, $k=\mathcal{O}/\lambda$, and
$\rho_1,\rho_2\colon G_\Q\to\GL_2(\mathcal{O})$ continuous and unramified outside a finite set $S$. Suppose $\tr\rho_1(\Frob_p)\equiv\tr\rho_2(\Frob_p)$ and
$\det\rho_1(\Frob_p)\equiv\det\rho_2(\Frob_p)\pmod\lambda$ for all $p$ in a set of density one, $p\notin S$.
\begin{enumerate}[label=(\alph*)]
\item $\bar\rho_1^{ss}\simeq\bar\rho_2^{ss}$ over $k$.
\item If $\bar\rho_1$ is absolutely irreducible, then $\bar\rho_2\simeq\bar\rho_1$.
\item Conversely, if $\bar\rho_1^{ss}\simeq\bar\rho_2^{ss}$, the congruences hold for every $p\notin S$.
\item (Newforms; given the representations attached to newforms, \Cref{ch:modgalrep}.) Let $\ell$ be odd, $E/\Q$ an elliptic curve of conductor $N$ with
$\bar\rho_{E,\ell}$ irreducible, and $h$ a newform of weight $2$ on $\Gamma_0(M)$ (trivial character) whose coefficients lie in $\mathcal{O}$ (for a chosen embedding of its coefficient field).
If $a_p(h)\equiv a_p(E)\pmod\lambda$ for all $p\nmid\ell MN$ in a set of density one, then $\bar\rho_h\simeq\bar\rho_{E,\ell}\otimes_{\F_\ell}k$, for every choice of the stable
lattice defining $\bar\rho_h$.
\end{enumerate}
\end{corollary}

\begin{proof}
(a) The characteristic polynomial of $\bar\rho_i(\Frob_p)$ is $X^2-\tr X+\det$ reduced modulo $\lambda$, and it is also that of $\bar\rho_i^{ss}(\Frob_p)$ (characteristic
polynomials are multiplicative along a composition series). The $\bar\rho_i^{ss}$ are continuous, semisimple and unramified outside $S$ (inertia acts trivially on
$\bar\rho_i$, hence on its subquotients); apply \Cref{cor:galois:chebotarev-bn}.
(b) $\bar\rho_2^{ss}\simeq\bar\rho_1^{ss}=\bar\rho_1$ is irreducible of dimension $2$, so the two-dimensional $\bar\rho_2$ has an irreducible composition factor of
dimension $2$, hence is irreducible and equal to its semisimplification.
(c) Characteristic polynomials of $\bar\rho_i(\Frob_p)$ and of $\bar\rho_i^{ss}(\Frob_p)$ agree.
(d) The black box of \Cref{ch:modgalrep} gives a continuous $\rho_{h,\lambda}\colon G_\Q\to\GL_2(\mathcal{O})$ (after choosing a stable lattice,
\Cref{prop:galois:continuous-reps}(c)), unramified outside $M\ell$, with $\det\rho_{h,\lambda}=\chi_\ell$ and $\tr\rho_{h,\lambda}(\Frob_p)=a_p(h)$ for $p\nmid M\ell$. Take
$\rho_1=\rho_{E,\ell}\otimes_{\Zl}\mathcal{O}$, which is unramified outside $\ell N$ with $\det\rho_1=\chi_\ell$ and $\tr\rho_1(\Frob_p)=a_p(E)$ for $p\nmid\ell N$
(\Cref{ch:ellcurves}, \Cref{ch:galrep}), and $\rho_2=\rho_{h,\lambda}$, with $S$ the primes dividing $\ell MN$. The determinants of Frobenius agree and the traces are
congruent on a set of density one. Since $\bar\rho_{E,\ell}$ is irreducible and $\det\bar\rho_{E,\ell}=\bar\chi_\ell$ with $\ell$ odd, it is absolutely irreducible
(\Cref{prop:galois:complex-conjugation}(d)), and so is $\bar\rho_1=\bar\rho_{E,\ell}\otimes_{\F_\ell}k$. By (b), $\bar\rho_h=\bar\rho_2\simeq\bar\rho_1$; the argument
applies to every stable lattice (and \Cref{prop:galois:reduction-ss} gives the independence of $\bar\rho_h^{ss}$ directly).
\end{proof}

\begin{paperlink}
Lemma lem:old, case $p\nmid\ell N$: ``$\bar\rho_h\simeq\bar\rho$ over a common finite extension of $\F_\ell$, by Chebotarev and Brauer-Nesbitt, $\bar\rho$ being irreducible,
hence absolutely irreducible''. This is \Cref{cor:galois:congruent-ap}(d), after enlarging $\mathcal{O}$ so that it contains the coefficients of the newform $h$ of the
lemma; it rests on the black box of \Cref{ch:modgalrep} (the representation $\rho_{h,\lambda}$ with $\tr\rho_{h,\lambda}(\Frob_p)=a_p(h)$). The two residual representations
a priori live over different finite fields ($k$ and $\F_\ell$), so they are compared over $k$, as $\bar\rho_{E,\ell}\otimes_{\F_\ell}k$; and ``irreducible, hence absolutely
irreducible'' is \Cref{prop:galois:complex-conjugation}(d). The equivalence quoted in the proof of Theorem thm:qual
(Section sec:proof: ``$\bar\rho_g\simeq\bar\rho$, that is, $a_p(g)\equiv a_p(f)\pmod\lambda$ for all $p\nmid 3N$''), and in the Example of that section, is (c) in one
direction and (d) in the other.
\end{paperlink}

\section{Exercises}\label{sec:galois:exercises}

The exercises are graded from routine to harder. Exercises \ref{exer:galois:726} and \ref{exer:galois:121} check the paper's statements on the $3$-division
polynomial on curves not treated in the text; Exercise \ref{exer:galois:twist-power} generalises the determinant argument of Lemma lem:h0.

\begin{exercise}\label{exer:galois:biquadratic}
Show that $\Q(\sqrt2,\sqrt3)=\Q(\sqrt2+\sqrt3)$, that the minimal polynomial of $\sqrt2+\sqrt3$ is $x^4-10x^2+1$, and list the intermediate fields with the
corresponding subgroups of $\Gal(\Q(\sqrt2,\sqrt3)/\Q)\simeq V_4$. Deduce that $x^4-10x^2+1$ is irreducible over $\Q$ but reducible modulo every prime $p$. (Hint: modulo $p$
one of $2$, $3$, $6$ is a square, and then the polynomial factors over $\F_p$; compare with \Cref{prop:galois:quartic}(c), case $V_4$.)
\end{exercise}

\begin{exercise}\label{exer:galois:x4minus2}
Let $\alpha=2^{1/4}\in\R$. Show that the splitting field of $x^4-2$ over $\Q$ is $\Q(\alpha,i)$, of degree $8$, with Galois group $D_4$, and identify the three quadratic
subfields. Show that over $\Q(i)$ the polynomial $x^4-2$ stays irreducible with cyclic Galois group of order $4$ (use \Cref{prop:galois:kummer-character}(c) with $n=4$,
$\mu_4\subset\Q(i)$, after checking that $2$ is not a square in $\Q(i)$). Confirm the result with \Cref{prop:galois:quartic}(c): the resolvent cubic of $x^4-2$ is
$y^3+8y=y(y^2+8)$.
\end{exercise}

\begin{exercise}\label{exer:galois:cubic7}
Let $z=\zeta_7+\zeta_7^{-1}=2\cos(2\pi/7)$. Show that $z$ is a root of $x^3+x^2-2x-1$, that this cubic has discriminant $49$ and roots $2\cos(2\pi k/7)$, $k=1,2,3$, and that
its splitting field is the real cubic subfield of $\Q(\zeta_7)$. Which primes $p$ split completely in it? (Answer: $p\equiv\pm1\pmod7$.)
\end{exercise}

\begin{exercise}\label{exer:galois:disc-prime-power}
Let $p$ be an odd prime and $k\ge1$. Using $(x^{p^{k-1}}-1)\Phi_{p^k}(x)=x^{p^k}-1$ and the method of \Cref{prop:galois:discriminant}(e), show that
\[
\disc(\Phi_{p^k})=(-1)^{\varphi(p^k)/2}\,p^{\,p^{k-1}(pk-k-1)}.
\]
Check the cases $p^k=9$ ($-3^9$) and $25$ ($5^{35}$). Deduce that the unique quadratic subfield of $\Q(\zeta_{p^k})$ is $\Q(\sqrt{p^*})$.
\end{exercise}

\begin{exercise}\label{exer:galois:conductor}
Show that $\Q(\sqrt{-5})\subseteq\Q(\zeta_{20})$ and $\Q(\sqrt6)\subseteq\Q(\zeta_{24})$. Show that if $\Q(\sqrt{-5})\subseteq\Q(\zeta_m)$ then $20\mid m$. (Hint: by
\Cref{thm:galois:cyclotomic}(c), $\Q(\sqrt{-5})\subseteq\Q(\zeta_m)\cap\Q(\zeta_{20})=\Q(\zeta_{\gcd(m,20)})$; list the quadratic subfields of $\Q(\zeta_g)$ for the proper divisors
$g$ of $20$.) Prove the analogous statement for $\Q(\sqrt6)$ and $24$.
\end{exercise}

\begin{exercise}\label{exer:galois:nfactorial}
Prove that the degree of a splitting field of a polynomial of degree $n$ divides $n!$. (Hint: induction on $n$; if $f$ is reducible, $f=gh$, use that $\deg g!\,\deg h!$
divides $n!$; if $f$ is irreducible, adjoin one root.)
\end{exercise}

\begin{exercise}\label{exer:galois:artin-schreier-trace}
Let $q=p^r$ and $a\in\F_q$. Show that $x^p-x-a$ is irreducible over $\F_q$ iff $\Tr_{\F_q/\F_p}(a)\ne0$. (Hint: $\wp\colon\F_q\to\F_q$ is $\F_p$-linear with kernel $\F_p$, so its
image has index $p$; show that $\Tr_{\F_q/\F_p}(\wp(b))=0$, and compare sizes, using \Cref{prop:galois:norm-trace}(d).)
\end{exercise}

\begin{exercise}\label{exer:galois:f13}
Let $u\in\F_{13}^\times$ and $\ell=3$, so $v=v_3(12)=1$. Show that $u$ is a cube in $\F_{13}$ iff $u^4=1$, that the cubes are $\{1,5,8,12\}$, and that the class of $u=2$ in
$\F_{13}^\times/(\F_{13}^\times)^3\simeq\mu_3(\F_{13})$ is $2^4=3$. Compute $\Frob_{13}(\alpha)/\alpha$ for a cube root $\alpha$ of $2$ in $\F_{13^3}$, and explain why $x^3-2$ is
irreducible over $\F_{13}$.
\end{exercise}

\begin{exercise}\label{exer:galois:supernatural}
Let $\widehat\Z\simeq G_{\F_p}$. Show that the closed subgroups of $\widehat\Z=\prod_q\Z_q$ are the products $\prod_q q^{a_q}\Z_q$ with $a_q\in\{0,1,2,\dots,\infty\}$
($q^\infty\Z_q=0$), and that the closed subgroup topologically generated by $\Frob_p^n$ is $n\widehat\Z$. Describe the fixed fields: for instance, the union of the fields
$\F_{p^{\ell^k}}$, $k\ge1$, has Galois group $\Z_\ell$ over $\F_p$ and is the fixed field of $\prod_{q\ne\ell}\Z_q$.
\end{exercise}

\begin{exercise}\label{exer:galois:small-subgroups}
Let $\ell$ be odd. Show that $\chi_\ell(G_\Q)=\Zl^\times$ is infinite, and that every neighbourhood of $1$ in $\GL_1(\Q_\ell)=\Q_\ell^\times$ contains a nontrivial closed subgroup
(namely $1+\ell^k\Zl$ for $k$ large). Explain why the argument of \Cref{prop:galois:continuous-reps}(d) fails over $\Q_\ell$. Show also that a continuous homomorphism
$\widehat\Z\to\C^\times$ is determined by the image of $1$, which is a root of unity, and that every root of unity occurs.
\end{exercise}

\begin{exercise}\label{exer:galois:dihedral}
Let $\ell\ge5$ and $\bar\rho\colon G_\Q\to\GL_2(\F_\ell)$ continuous and absolutely irreducible, with image contained in the normaliser $N_s(\F_\ell)$ of the diagonal split
Cartan subgroup $C_s(\F_\ell)$ but not in $C_s(\F_\ell)$. Show: (a) the composite $G_\Q\to N_s(\F_\ell)/C_s(\F_\ell)\simeq\{\pm1\}$ is $\chi_d$ for a unique quadratic field $\Q(\sqrt d)$;
(b) $\bar\rho\simeq\bar\rho\otimes\chi_d$ and $\bar\rho\simeq\Ind_{G_{\Q(\sqrt d)}}^{G_\Q}\psi$ for a character $\psi$ of $G_{\Q(\sqrt d)}$ with values in $\F_\ell^\times$; (c) if $d=\ell^*$, the
hypothesis of Theorem thm:main fails for $\bar\rho$, although $\bar\rho\not\simeq\bar\rho\otimes\bar\chi_\ell$.
\end{exercise}

\begin{exercise}\label{exer:galois:726}
The paper's example 726a1 is $y^2+xy=x^3+x^2-35x-51$, with mod $3$ image 3Ns. Check with PARI (\texttt{elldivpol(E,3)}, \texttt{factor}) that
\[
\psi_3=3x^4+5x^3-210x^2-612x-1480=(x^2-x-74)(3x^2+8x+20),
\]
that the two factors have discriminants $297=3^3\cdot11$ and $-176=-2^4\cdot11$, so that $K_1=\Q(\sqrt{33})$, $K_2=\Q(\sqrt{-11})$ and the third quadratic subfield of $K_1K_2$ is
$\Q(\sqrt{-3})$, as \Cref{prop:galois:quartic-on-lines}(e) requires. Verify that $a_p\equiv0\pmod 3$ for the primes $p\equiv2\pmod3$, $p\nmid726$, $p<100$, and find the
primes $p<60$, $p\nmid726$, at which $\psi_3\bmod p$ splits into four distinct linear factors (answer: $31$, $37$). (At $p=2$ and $p=11$ it also splits into linear factors,
with repetitions; at $p=3$ its degree drops.)
\end{exercise}

\begin{exercise}\label{exer:galois:121}
For 121a1, $y^2+xy+y=x^3+x^2-30x-76$, check that $\psi_3=3x^4+5x^3-177x^2-909x-1249$ is irreducible with $\disc\psi_3=-3^3\cdot11^4$ and irreducible resolvent cubic, so
that its group is $S_4$ and the image of $\bar\rho_{E,3}$ is $\GL_2(\F_3)$. With $a_{19}=-6$ and $a_{29}=-9$, predict the factorisation patterns of $\psi_3$ modulo $19$ and $29$
from \Cref{ex:galois:frobenius-on-lines}, and verify them.
\end{exercise}

\begin{exercise}\label{exer:galois:twist-power}
Let $\ell$ be odd, $\rho\colon G_\Q\to\GL_2(\Zl)$ continuous and $m\ge1$. Show that if $\rho\bmod\ell^k\simeq(\rho\otimes\chi_\ell^m)\bmod\ell^k$ with $k\ge1$, then
$\ell^{k-1}(\ell-1)$ divides $2m$. (Hint: $\chi_\ell^{2m}\equiv1\pmod{\ell^k}$ on $G_\Q$, and $(\Z/\ell^k)^\times$ is cyclic of order $\ell^{k-1}(\ell-1)$, \Cref{ch:rings}.) Recover
\Cref{prop:galois:twist}(c) as the case $m=1$, and show that for $\ell=5$, $m=2$ the condition allows only $k=1$.
\end{exercise}

\begin{exercise}\label{exer:galois:x3-3x+1}
Using \Cref{thm:galois:frobenius} and \Cref{thm:galois:chebotarev}, show that $x^3-3x+1$ has either three roots or none modulo every prime $p\ne3$, that it has three roots
iff $p\equiv\pm1\pmod9$, and that the set of these primes has density $1/3$. Count the primes $5\le p<3000$ with three roots (there are $142$ of $428$).
\end{exercise}

\begin{exercise}\label{exer:galois:restriction-index}
Let $E/\Q$ be an elliptic curve with $\bar\rho_{E,5}$ surjective onto $\GL_2(\F_5)$. Compute the images of $G_{\Q(\sqrt5)}$ and of $G_{\Q(\zeta_5)}$ under $\bar\rho_{E,5}$ and their
orders ($240$ and $120$). Show that $\bar\rho_{E,5}$ is absolutely irreducible on $G_{\Q(\sqrt5)}$. (Use \Cref{prop:galois:restriction}(c) and the fact that $\SL_2(\F_5)$ acts
irreducibly on $\Fbar_5^2$, \Cref{ch:groups}.)
\end{exercise}

\section{Notes and references}\label{sec:galois:notes}

The finite theory (Sections 1 to 6) follows \cite[Chapters V and VI]{Lang}, \cite[Chapters 13 and 14]{DummitFoote} and \cite{MilneFields}; the proofs of the primitive
element theorem, of Artin's theorem and of Hilbert's Theorem 90 are the classical ones. The irreducibility proof of $\Phi_n$ is Dedekind's, as in \cite[Chapter VI,
Section 3]{Lang} and \cite[Chapter 2]{Washington}; Gauss sums and $g^2=\ell^*$ are \Cref{ch:rings} and \cite[Chapter 6]{IrelandRosen}. Kummer theory and Artin-Schreier theory
are \cite[Chapter VI, Sections 6 and 8]{Lang}; Capelli's criterion is \cite[Chapter VI, Section 9]{Lang}. The resolvent cubic and the Galois groups of quartics are
\cite[Section 14.6]{DummitFoote}; the $C_4$ versus $D_4$ test is due to Kappe and Warren \cite{KappeWarren89}. Infinite Galois theory and the Krull topology follow
\cite[Chapter IV, Section 1]{Neukirch}, \cite[Chapter I, Section 1]{NSW} and \cite[Chapter 7]{MilneFields}. Continuous $\ell$-adic representations, lattices, the
independence of $\bar\rho^{ss}$ from the lattice and the use of Chebotarev's theorem with Brauer-Nesbitt are in \cite[Chapter I]{SerreAbelian} and \cite[Section 1]{DDT}; the
Chebotarev density theorem itself is \cite[Chapter VII, Section 13]{Neukirch}, and the comparison of densities \cite[Chapter VI, Section 4]{SerreCourse}. The theorem of
Artin and Schreier is in \cite[Chapter VI, Section 9]{Lang}, and the real closed fields behind the statement on elements of finite order in \cite[Chapter XI]{Lang}. The Kummer cocycle of a Tate parameter and its matrix form are \cite[Chapter IV]{SerreAbelian}
and \Cref{ch:reduction}; its cohomological meaning is \cite[Chapter VI]{NSW} and \Cref{ch:galcoh}.

The chapter and section numbers of the books cited in this chapter (Lang, Dummit-Foote, Neukirch, Neukirch-Schmidt-Wingberg, Serre, Washington, Ireland-Rosen,
Milne) are pointers recalled from the editions listed in the bibliography and have not been checked against the printed books; the only theorem number given,
Lang's Theorem VI.9.1 for Capelli's criterion, is unverified in the same sense. The numerical statements (discriminants, factorisations of $\psi_3$, Frobenius patterns,
traces $a_p$, the subgroup table of $\GL_2(\F_3)$, the prime counts) were recomputed for this chapter with PARI/GP 2.17 through cypari2 and a short enumeration of
$\GL_2(\F_3)$; the curve data agree with Cremona's tables (6912d1 and 726a1 have mod $3$ image 3Ns there).

Departures from the paper: none of the statements of this chapter strengthens a statement of the paper, so no Commentary comparing a stronger statement with the paper's is needed. The
paper's $\chi_q$ is written $\chi_{q^*}$ throughout, and the paper's $\bar\chi_3^{\mathrm{cyc}}$ is $\bar\chi_3$. The revised Lemma lem:h0 and Section 1.3 of the paper use
$\bar\chi_\ell$ where the original version used $\chi_{\ell^*}$; \Cref{prop:galois:twist}(c) and \Cref{prop:galois:sqrt-lstar}(c) are the facts behind the correction.
